% Les fêtes traditionnelles — Allemand 1ère A — Cours signé OnBuch+
% Programme officiel MINESEC. Compiler avec : tectonic ALA21-les-fetes-traditionnelles.tex
\newcommand{\DOCMATIERE}{Allemand}
\newcommand{\DOCNIVEAU}{1ère A}
\newcommand{\DOCMODULE}{Chapitre 5 : Citoyenneté}
\newcommand{\DOCLECON}{ALA21}
\newcommand{\DOCTITRE}{Les fêtes traditionnelles}
% OnBuch+ — préambule « cours signés » ENRICHI (compilé avec tectonic / XeLaTeX)
% Système visuel « Pop » d'OnBuch+ : fond crème (jamais blanc), bordures encre
% épaisses, ombres « dures » décalées, titres Archivo Black en capitales.
% Version MATHÉMATIQUES : graphes (pgfplots/tikz), boîtes pédagogiques
% (définition, propriété, méthode, attention, le savais-tu, exercice résolu,
% exercice, TP, bilan), page de garde et signature OnBuch+ ancrée en bas.
%
% Macros attendues dans main.tex AVANT \input{preamble} :
%   \DOCMATIERE  \DOCNIVEAU  \DOCMODULE  \DOCLECON (numéro)  \DOCTITRE
\documentclass[11pt]{article}

\usepackage{fontspec}
\usepackage[ngerman,french]{babel} % français = langue principale ; l'allemand via \all{..} / textebox
\usepackage[a4paper,margin=2cm,top=3.4cm,bottom=2.8cm,headheight=1.4cm,footskip=1.3cm]{geometry}
\usepackage{amsmath,amssymb,amsfonts}
\usepackage{mathtools} % \xmapsto, \coloneqq... souvent utilisés par les modèles
\usepackage{cancel} % simplifications barrées (bilans de forces)
% \abs{z} (module, valeur absolue) et \norm{u} (norme), utilisés par les modèles
\providecommand{\abs}{}\let\abs\relax\DeclarePairedDelimiter{\abs}{\lvert}{\rvert}
\providecommand{\norm}{}\let\norm\relax\DeclarePairedDelimiter{\norm}{\lVert}{\rVert}
\usepackage[table]{xcolor} % [table] : fournit aussi \rowcolors (alternance de lignes)
\usepackage{pagecolor}
\usepackage{tikz}
\usetikzlibrary{arrows.meta,positioning,calc,shapes.geometric,shapes.misc,shapes.arrows,decorations.pathmorphing,decorations.pathreplacing,decorations.markings,patterns,shadows,mindmap,trees,fit,backgrounds,matrix,intersections,angles,quotes}
\usepackage{pgfplots}
\usepackage[version=4]{mhchem} % équations nucléaires \ce{^{235}_{92}U}
\usepackage[european,siunitx]{circuitikz} % schémas électriques
\pgfplotsset{compat=1.18}
\usepgfplotslibrary{fillbetween,groupplots} % aires entre courbes (calcul intégral)
\usepackage{enumitem}
\usepackage{fancyhdr}
% [most] charge toutes les bibliothèques tcolorbox (listings, minted, raster,
% documentation...) et gonfle inutilement la mémoire TeX sur un document long
% (~300 boîtes) : on ne charge que ce qui est réellement utilisé.
\usepackage[skins,breakable]{tcolorbox}
\usepackage{graphicx}
\usepackage{colortbl}
\usepackage{booktabs}
\usepackage{tabularx}
\usepackage{diagbox} % case d'en-tête diagonale des tableaux à double entrée
% Colonnes extensibles centrée / gauche / droite (C, L, R), souvent utilisées par les modèles
\newcolumntype{C}{>{\centering\arraybackslash}X}
\newcolumntype{L}{>{\raggedright\arraybackslash}X}
\newcolumntype{R}{>{\raggedleft\arraybackslash}X}
\usepackage{array}
\usepackage{multicol}
\usepackage{multirow}
\usepackage{siunitx}
% parse-numbers=false : accepte aussi \SI{2,5 \times 10^{-3}}{...}
\sisetup{locale=FR,output-decimal-marker={,},group-minimum-digits=4,parse-numbers=false}
\DeclareSIUnit\ohm{\ensuremath{\symup{\Omega}}} % U+2126 absent de la police
\DeclareSIUnit\division{div} % réglages d'oscilloscope (ms/div, V/div)
\DeclareSIUnit\tr{tr} % tours (vitesse de rotation, ex. \SI{1500}{\tr\per\minute})
\DeclareSIUnit\molar{M} % concentration molaire (ex. \SI{3}{\molar}, \SI{1}{\micro\molar})
\DeclareSIUnit\an{an} % année (le modèle confond parfois avec \year, un compteur TeX) — ex. \SI{5}{\kilo\meter\per\an}
\DeclareSIUnit\FCFA{FCFA} % franc CFA (ex. \SI{500}{\FCFA\per\kilo\gram})
\DeclareSIUnit\degreeBrix{\textdegree Brix} % teneur en sucre d'un sirop/jus (réfractométrie)
\DeclareSIUnit\month{mois} % le modèle utilise parfois \month, un compteur TeX, comme unité de durée
\usepackage{qrcode}
% \degree : commande fréquemment utilisée par les modèles pour un angle ou une
% température en dehors de \SI{}{} (non fournie par les paquets chargés).
\providecommand{\degree}{^{\circ}}
% \qed (fin de démonstration), utilisé par les modèles sans amsthm
\providecommand{\qed}{\hfill$\square$}
% \barre : double barre des valeurs interdites dans un tableau de variations (tkz-tab)
\providecommand{\barre}{\ensuremath{\|}}
% environnement proof (amsthm non chargé) utilisé par les modèles
\ifdefined\proof\else\newenvironment{proof}[1][Démonstration]{\par\noindent\textit{#1.}\ }{\qed\par}\fi
\usepackage{lastpage}
\usepackage{hyperref}
\hypersetup{hidelinks}

% ============================================================
% Palette « Pop » OnBuch+ — mêmes hex que la classe P (plus_ui.dart)
% ============================================================
\definecolor{poporange}{HTML}{F59321}
\definecolor{popdark}{HTML}{E07A0C}
\definecolor{popcream}{HTML}{FFF6E8}
\definecolor{popink}{HTML}{1C1714}
\definecolor{poppurple}{HTML}{7A5AE0}
\definecolor{popgreen}{HTML}{2E9E6A}
\definecolor{poppink}{HTML}{E0457B}
\definecolor{popblue}{HTML}{2F7DE1}
\definecolor{popgold}{HTML}{C98A12}
\definecolor{popmuted}{HTML}{8A7F76}
\definecolor{popline}{HTML}{EADFD2}
\definecolor{popgray}{HTML}{8A7F76}
\colorlet{popgrey}{popgray}
% Alias tolérants (le modèle nomme parfois les couleurs différemment)
\colorlet{popwhite}{white}
\colorlet{popblack}{popink}
\colorlet{popred}{poppink}
\colorlet{popyellow}{popgold}
% Teintes claires (fonds de tableaux/encadrés) — le modèle nomme parfois
% les couleurs avec un suffixe L (ex. popblueL) sans les avoir définies.
\colorlet{poporangeL}{poporange!20}
\colorlet{popdarkL}{popdark!20}
\colorlet{poppurpleL}{poppurple!20}
\colorlet{popgreenL}{popgreen!20}
\colorlet{poppinkL}{poppink!20}
\colorlet{popblueL}{popblue!20}
\colorlet{popgoldL}{popgold!20}
\colorlet{popredL}{popred!20}
\colorlet{popgrayL}{popgray!20}
% Teintes claires pour fonds de tableaux / zones
\colorlet{popblueL}{popblue!10!white}
\colorlet{poporangeL}{poporange!14!white}
\colorlet{popgreenL}{popgreen!12!white}
\colorlet{poppurpleL}{poppurple!12!white}
\colorlet{poppinkL}{poppink!12!white}
\colorlet{popgoldL}{popgold!14!white}

\pagecolor{popcream}

% ============================================================
% Typographie
% ============================================================
\defaultfontfeatures{Ligatures=TeX}
\setmainfont{PlusJakartaSans-Medium.ttf}[Path=../../../fonts/,BoldFont=PlusJakartaSans-Bold.ttf]
% Symboles, API et emojis absents de la police principale : repli sur DejaVu Sans (caractères actifs)
\newfontface\symfont{DejaVuSans.ttf}[Path=../../../fonts/]
\makeatletter
\newcommand\symactive[1]{\catcode#1=13 \begingroup\lccode`\~=#1 \lowercase{\endgroup\def~}{{\symfont\char#1}}}
\newcommand\symblank[1]{\catcode#1=13 \begingroup\lccode`\~=#1 \lowercase{\endgroup\def~}{}}
\newcommand\symhyph[1]{\catcode#1=13 \begingroup\lccode`\~=#1 \lowercase{\endgroup\def~}{\mbox{-}}}
\newcount\symc
\newcommand\symrange[2]{\symc=#1 \@whilenum\symc<#2 \do{\expandafter\symactive\expandafter{\the\symc}\advance\symc by 1 }}
\newcommand\symblankrange[2]{\symc=#1 \@whilenum\symc<#2 \do{\expandafter\symblank\expandafter{\the\symc}\advance\symc by 1 }}
\makeatother
\symrange{"0250}{"02FF}\symactive{"2713}\symactive{"2714}\symactive{"2717}\symactive{"2718}\symactive{"2610}\symactive{"2611}\symactive{"2612}
\symhyph{"2011}\symhyph{"2E1A}\symblankrange{"1F300}{"1FAFF}\symblank{"FE0F}
% Maths en sans-serif (Fira Math) avec chiffres et lettres droites de Plus
% Jakarta, pour que les formules se fondent dans le texte.
\usepackage[math-style=ISO,bold-style=ISO]{unicode-math}
\setmathfont{FiraMath-Regular.otf}[Path=../../../fonts/]
\setmathfont{PlusJakartaSans-Medium.ttf}[Path=../../../fonts/,range={up/num,up/{Latin,latin},\mathup}]
\usepackage{icomma} % virgule décimale sans espace en mode math
% Caractères mathématiques Unicode absents des polices de texte (Plus Jakarta,
% Archivo Black) : rendus par leur équivalent mathématique, en texte comme en
% formule — sinon ils s'affichent en carré vide (ex. « Arithmétique dans ℤ »).
\usepackage{newunicodechar}
\newunicodechar{ℝ}{\ensuremath{\mathbb{R}}}
\newunicodechar{ℤ}{\ensuremath{\mathbb{Z}}}
\newunicodechar{ℂ}{\ensuremath{\mathbb{C}}}
\newunicodechar{ℕ}{\ensuremath{\mathbb{N}}}
\newunicodechar{ℚ}{\ensuremath{\mathbb{Q}}}
\newunicodechar{𝔻}{\ensuremath{\mathbb{D}}}
\newunicodechar{α}{\ensuremath{\alpha}}
\newunicodechar{β}{\ensuremath{\beta}}
\newunicodechar{γ}{\ensuremath{\gamma}}
\newunicodechar{δ}{\ensuremath{\delta}}
\newunicodechar{ε}{\ensuremath{\varepsilon}}
\newunicodechar{θ}{\ensuremath{\theta}}
\newunicodechar{λ}{\ensuremath{\lambda}}
\newunicodechar{μ}{\ensuremath{\mu}}
\newunicodechar{σ}{\ensuremath{\sigma}}
\newunicodechar{τ}{\ensuremath{\tau}}
\newunicodechar{φ}{\ensuremath{\varphi}}
\newunicodechar{ψ}{\ensuremath{\psi}}
\newunicodechar{ω}{\ensuremath{\omega}}
\newunicodechar{Σ}{\ensuremath{\Sigma}}
% majuscules produites par \MakeUppercase dans les titres (α-aminés -> Α)
\newunicodechar{Α}{\ensuremath{\alpha}}
\newunicodechar{Β}{\ensuremath{\beta}}
\newunicodechar{Γ}{\ensuremath{\Gamma}}
\newunicodechar{Δ}{\ensuremath{\Delta}}
\newunicodechar{Ε}{\ensuremath{\varepsilon}}
\newunicodechar{Θ}{\ensuremath{\Theta}}
\newunicodechar{Λ}{\ensuremath{\Lambda}}
\newunicodechar{Μ}{\ensuremath{\mu}}
\newunicodechar{Τ}{\ensuremath{\tau}}
\newunicodechar{Φ}{\ensuremath{\Phi}}
\newunicodechar{Ψ}{\ensuremath{\Psi}}
\newunicodechar{Ω}{\ensuremath{\Omega}}
\newunicodechar{↦}{\ensuremath{\mapsto}}
\newunicodechar{∈}{\ensuremath{\in}}
\newunicodechar{∉}{\ensuremath{\notin}}
\newunicodechar{∧}{\ensuremath{\wedge}}
\newunicodechar{⇔}{\ensuremath{\Leftrightarrow}}
\newunicodechar{⟺}{\ensuremath{\Longleftrightarrow}}
\newunicodechar{⇒}{\ensuremath{\Rightarrow}}
\newunicodechar{⟹}{\ensuremath{\Longrightarrow}}
\newunicodechar{∘}{\ensuremath{\circ}}
\newunicodechar{∪}{\ensuremath{\cup}}
\newunicodechar{∩}{\ensuremath{\cap}}
\newunicodechar{≡}{\ensuremath{\equiv}}
\newunicodechar{⊂}{\ensuremath{\subset}}
\newunicodechar{⊥}{\ensuremath{\perp}}
\newunicodechar{∃}{\ensuremath{\exists}}
\newunicodechar{∀}{\ensuremath{\forall}}
\newunicodechar{∛}{\ensuremath{\sqrt[3]{\,}}}
\newunicodechar{∜}{\ensuremath{\sqrt[4]{\,}}}
\newunicodechar{⃗}{\ensuremath{{}^{\to}}}
\newunicodechar{ⁿ}{\ifmmode^{n}\else\textsuperscript{\ensuremath{n}}\fi}
\newunicodechar{ˣ}{\ifmmode^{x}\else\textsuperscript{\ensuremath{x}}\fi}
\newunicodechar{ᵇ}{\ifmmode^{b}\else\textsuperscript{\ensuremath{b}}\fi}
\newunicodechar{ʳ}{\ifmmode^{r}\else\textsuperscript{\ensuremath{r}}\fi}
\newunicodechar{ᵘ}{\ifmmode^{u}\else\textsuperscript{\ensuremath{u}}\fi}
\newunicodechar{ʸ}{\ifmmode^{y}\else\textsuperscript{\ensuremath{y}}\fi}
\newunicodechar{ᵃ}{\ifmmode^{a}\else\textsuperscript{\ensuremath{a}}\fi}
\newunicodechar{ᵉ}{\ifmmode^{e}\else\textsuperscript{\ensuremath{e}}\fi}
\newunicodechar{ᵏ}{\ifmmode^{k}\else\textsuperscript{\ensuremath{k}}\fi}
\newunicodechar{ᵖ}{\ifmmode^{p}\else\textsuperscript{\ensuremath{p}}\fi}
\newunicodechar{ᴺ}{\ifmmode^{N}\else\textsuperscript{\ensuremath{N}}\fi}
\newunicodechar{ᶜ}{\ifmmode^{c}\else\textsuperscript{\ensuremath{c}}\fi}
\newunicodechar{ʰ}{\ifmmode^{h}\else\textsuperscript{\ensuremath{h}}\fi}
\newunicodechar{ᵠ}{\ifmmode^{q}\else\textsuperscript{\ensuremath{q}}\fi}
\newunicodechar{ₙ}{\ifmmode_{n}\else\textsubscript{\ensuremath{n}}\fi}
\newunicodechar{ₐ}{\ifmmode_{a}\else\textsubscript{\ensuremath{a}}\fi}
\newunicodechar{ᵢ}{\ifmmode_{i}\else\textsubscript{\ensuremath{i}}\fi}
\newunicodechar{ᵣ}{\ifmmode_{r}\else\textsubscript{\ensuremath{r}}\fi}
\newunicodechar{ₖ}{\ifmmode_{k}\else\textsubscript{\ensuremath{k}}\fi}
\newunicodechar{ᵦ}{\ifmmode_{\beta}\else\textsubscript{\ensuremath{\beta}}\fi}
\newunicodechar{ₓ}{\ifmmode_{x}\else\textsubscript{\ensuremath{x}}\fi}
\newfontfamily\popdisplay{ArchivoBlack-Regular.ttf}[Path=../../../fonts/]
\newfontfamily\poplogo{SpaceGrotesk-Bold.ttf}[Path=../../../fonts/]
\newfontfamily\popsemi{PlusJakartaSans-SemiBold.ttf}[Path=../../../fonts/]
\newfontfamily\popxbold{PlusJakartaSans-ExtraBold.ttf}[Path=../../../fonts/]
\newfontfamily\popxboldls{PlusJakartaSans-ExtraBold.ttf}[Path=../../../fonts/,LetterSpace=3.0]

\setlist[enumerate]{leftmargin=1.7em,itemsep=5pt,topsep=4pt}
\setlist[itemize]{leftmargin=1.7em,itemsep=4pt,topsep=4pt,label={\color{poporange}\textbullet}}
\setlength{\parindent}{0pt}
\setlength{\parskip}{5pt}
\renewcommand{\arraystretch}{1.25}

% pgfplots : style Pop par défaut
\pgfplotsset{
  every axis/.append style={axis line style={popink,line width=1pt},tick style={popink},
    grid style={popline,line width=0.5pt},label style={font=\small},tick label style={font=\footnotesize}},
  popcurve/.style={poporange,line width=1.8pt,no markers,smooth},
}
% Même style disponible dans un \draw TikZ simple (hors pgfplots)
\tikzset{popcurve/.style={poporange,line width=1.8pt}}
\tikzset{popgrid/.style={popline,very thin}}
\colorlet{popcurve}{poporange} % le nom du style est parfois utilisé comme couleur

% ============================================================
% Pill / squiggle
% ============================================================
\newcommand{\poppill}[3]{%
  \tikz[baseline=-0.55ex]\node[draw=popink,line width=1.2pt,rounded corners=2.6pt,%
    fill=#2,inner xsep=6.5pt,inner ysep=3pt]{%
    \popxboldls\fontsize{7}{7}\selectfont\color{#3}\MakeUppercase{#1}};%
}
\newcommand{\popsquiggle}[1]{%
  \tikz[baseline=-0.4ex]{
    \draw[#1,line width=2.6pt,line cap=round]
      plot [smooth] coordinates {(0,0.09) (0.34,0.01) (0.68,0.21) (1.02,0.01) (1.36,0.10)};
  }%
}
% Mot-clé surligné (marqueur orange sous le texte)
\newcommand{\cle}[1]{{\popxbold\color{popdark}#1}}

% ============================================================
% En-tête / pied
% ============================================================
\pagestyle{fancy}
\fancyhf{}
\renewcommand{\headrulewidth}{0pt}
\renewcommand{\footrulewidth}{0pt}
\lhead{\raisebox{-0.15ex}{\poplogo\fontsize{13}{13}\selectfont\color{popink}OnBuch\color{poporange}+}}
\rhead{\poppill{\DOCMATIERE\ \textbullet\ \DOCNIVEAU\ \textbullet\ Leçon \DOCLECON}{white}{popink}}
\lfoot{\popxbold\fontsize{7.6}{7.6}\selectfont\color{popmuted}\MakeUppercase{Cours signé OnBuch+}\ \textbullet\ {\popsemi\DOCTITRE}}
\rfoot{\tikz[baseline=-0.6ex]\node[rounded corners=8.5pt,fill=popink,minimum height=17pt,minimum width=17pt,inner xsep=5pt,inner sep=0]{\popxbold\fontsize{8}{8}\selectfont\color{popcream}\thepage\,/\,\pageref*{LastPage}};}

% ============================================================
% Titres
% ============================================================
\newcommand{\chaptitre}[2]{%
  \begin{tcolorbox}[enhanced,colback=poporange,colframe=popink,boxrule=2.6pt,arc=11pt,
      drop shadow=popink,left=15pt,right=15pt,top=12pt,bottom=11pt,before skip=3pt,after skip=18pt]
    \poppill{#2}{popink}{popcream}\par\vspace{9pt}
    {\raggedright\hyphenpenalty=10000\exhyphenpenalty=10000\popdisplay\fontsize{22}{24}\selectfont\color{popcream}\MakeUppercase{#1}\par}
  \end{tcolorbox}
}
% Section numérotée : pastille orange avec le numéro + titre Archivo
\newcounter{popsec}
\newcommand{\coursec}[1]{%
  \stepcounter{popsec}\setcounter{popsub}{0}%
  \par\vspace{16pt}\noindent
  \begin{minipage}{\linewidth}
  \tikz[baseline=-0.7ex]\node[draw=popink,line width=1.6pt,fill=poporange,rounded corners=4pt,minimum size=22pt,inner sep=0]{\popdisplay\fontsize{12}{12}\selectfont\color{popcream}\thepopsec};%
  \hspace{8pt}{\popdisplay\fontsize{15}{17}\selectfont\color{popink}\MakeUppercase{#1}}\par
  \vspace{4pt}\noindent{\color{poporange}\rule{36pt}{3.6pt}}
  \end{minipage}\par\nopagebreak\vspace{8pt}
}
\newcounter{popsub}
\newcommand{\courssub}[1]{%
  \stepcounter{popsub}%
  \par\vspace{9pt}\noindent{\popxbold\fontsize{11.5}{13}\selectfont\color{popink}{\color{poporange}\thepopsec.\thepopsub}\hspace{6pt}#1}\par\nopagebreak\vspace{4pt}
}

% ============================================================
% Boîtes pédagogiques — même recette : fond blanc, bordure épaisse colorée,
% ombre dure encre, badge plein. Titre optionnel [#1].
% ============================================================
\tcbset{popbase/.style={breakable,enhanced,colback=white,boxrule=2.3pt,arc=9pt,
  shadow={3pt}{-3pt}{0pt}{popink},left=14pt,right=14pt,top=11pt,bottom=11pt,before skip=11pt,after skip=15pt,
  fontupper=\popsemi\fontsize{10.6}{14.5}\selectfont\color{popink},
  fontlower=\popsemi\fontsize{10.6}{14.5}\selectfont\color{popink}}}
% Coche dessinée (le glyphe ✓ n'existe pas dans les polices du document)
\newcommand{\popcheck}[1]{\tikz[baseline=-0.5ex]\draw[#1,line width=1.6pt,line cap=round,line join=round](-0.13,0.0)--(-0.04,-0.09)--(0.13,0.1);}
\providecommand{\coche}{\popcheck{popgreen}}
\providecommand{\resultat}[1]{\par\smallskip\noindent\fcolorbox{popgreen}{popgreenL}{\parbox{\dimexpr\linewidth-2\fboxsep-2\fboxrule\relax}{\textbf{Résultat.} #1}}\par\smallskip}
\newcommand{\popboxhead}[3]{\poppill{#1}{#2}{white}\ifx\relax#3\relax\else\hspace{6pt}{\popxbold\fontsize{10.6}{12}\selectfont\color{popink}#3}\fi\par\vspace{7pt}}

\newtcolorbox{aretenir}[1][]{popbase,colframe=popgreen,before upper={\popboxhead{À retenir}{popgreen}{#1}}}
% Texte en allemand (document de lecture, dialogue, extrait) : cadre
% d'étude. Un titre optionnel (source, auteur) s'affiche dans l'en-tête.
\newtcolorbox{textebox}[1][]{popbase,colframe=popblue,before upper={\popboxhead{Text}{popblue}{#1}\begin{otherlanguage}{ngerman}},after upper={\end{otherlanguage}}}
% Marques ✓ / ✗ (forme correcte / fautive) souvent utilisées par les modèles
\providecommand{\cross}{\textcolor{poppink}{\ensuremath{\times}}}
\providecommand{\xmark}{\cross}\providecommand{\crossmark}{\cross}\providecommand{\cmark}{\ensuremath{\checkmark}}
% Ligne de réponse / justification à compléter (inventée par les modèles)
\providecommand{\justification}[1]{\par\noindent\textit{Justification~:} \dotfill\par}
% \textipa (package tipa non chargé) : la police ne couvre pas l'API -> simple passage
\providecommand{\textipa}[1]{#1}
% Soulignement (ulem non chargé) : \uline, \ul
\providecommand{\uline}[1]{\underline{#1}}\providecommand{\ul}[1]{\underline{#1}}
% \glossaire{\item ...} inventée par les modèles : liste de glossaire
\providecommand{\glossaire}[1]{\begin{itemize}#1\end{itemize}}
% \alert (beamer) inventée par les modèles : texte mis en évidence
\providecommand{\alert}[1]{\textbf{\textcolor{poppink}{#1}}}
% variantes ASCII d'ordinaux écrites en macro
\providecommand{\iere}{\textsuperscript{re}}\providecommand{\ieme}{\textsuperscript{e}}\providecommand{\ere}{\textsuperscript{re}}\providecommand{\eme}{\textsuperscript{e}}\providecommand{\er}{\textsuperscript{er}}\providecommand{\ie}{\textsuperscript{e}}
% Ordinaux français écrits en macro par les modèles : 1\ière, 3\ième
\newcommand{\ière}{\textsuperscript{re}}\newcommand{\ième}{\textsuperscript{e}}
% \exemple (inventée par les modèles) : étiquette « Exemple : »
\providecommand{\exemple}{\textbf{Exemple~:}\ }
% \square utilisable aussi hors mode math (cases à cocher)
\AtBeginDocument{\let\squareMath\square\renewcommand{\square}{\ensuremath{\squareMath}}}
% Mot ou phrase en allemand à l'intérieur d'un paragraphe français
\newcommand{\all}[1]{\ifmmode\text{\textit{#1}}\else\begingroup\selectlanguage{ngerman}\itshape #1\endgroup\fi}
\let\esp\all % alias toléré
\newtcolorbox{exemplebox}[1][]{popbase,colframe=poporange,before upper={\popboxhead{Exemple}{poporange}{#1}}}
\newtcolorbox{definition}[1][]{popbase,colframe=popblue,before upper={\popboxhead{Définition}{popblue}{#1}}}
\newtcolorbox{propriete}[1][]{popbase,colframe=poppurple,before upper={\popboxhead{Propriété}{poppurple}{#1}}}
\newtcolorbox{methode}[1][]{popbase,colframe=popgold,before upper={\popboxhead{Méthode}{popgold}{#1}}}
\newtcolorbox{attention}[1][]{popbase,colframe=poppink,before upper={\popboxhead{Attention}{poppink}{#1}}}
\newtcolorbox{experience}[1][]{popbase,colframe=popblue,colback=popblueL,before upper={\popboxhead{Expérience}{popblue}{#1}}}
% « Le savais-tu ? » : carte violette pleine (contexte, culture, Cameroun)
\newtcolorbox{savaistu}[1][]{popbase,colback=poppurple,colframe=popink,
  fontupper=\popsemi\fontsize{10.4}{14.2}\selectfont\color{white},
  before upper={\poppill{Le savais-tu\,?}{white}{poppurple}\ifx\relax#1\relax\else\hspace{6pt}{\popxbold\color{white}#1}\fi\par\vspace{7pt}}}
% Exercice résolu : énoncé en haut, \tcblower puis la solution sur fond crème
\newcounter{popexo}\newcounter{popexr}
\newtcolorbox{exoresolu}[1][]{popbase,colframe=poporange,colbacklower=popcream,
  segmentation style={popink,line width=1.2pt,dashed},
  before upper={\stepcounter{popexr}\popboxhead{Exercice résolu \thepopexr}{poporange}{#1}},
  before lower={\poppill{Solution}{popink}{popcream}\par\vspace{6pt}}}
% Exercice (non résolu en place) — la correction est dans la partie Corrigés
\newtcolorbox{exercice}[1][]{popbase,colframe=popink,
  before upper={\stepcounter{popexo}\popboxhead{Exercice \thepopexo}{popink}{#1}}}
% Corrigé
\newtcolorbox{corrige}[1][]{popbase,colframe=popgreen,colback=popgreenL,
  before upper={\popboxhead{Corrigé}{popgreen}{#1}}}
% Objectifs / prérequis / bilan
\newtcolorbox{objectifs}{popbase,colframe=popink,colback=poporangeL,
  before upper={\popboxhead{Objectifs de la leçon}{popink}{}}}
\newtcolorbox{prerequis}{popbase,colframe=popink,colback=popblueL,
  before upper={\popboxhead{Prérequis}{popblue}{}}}
\newtcolorbox{bilan}{popbase,colframe=popink,colback=white,boxrule=2.8pt,
  before upper={\popboxhead{Fiche bilan}{poporange}{L'essentiel à connaître pour l'examen}}}
% Figure encadrée avec légende
\newtcolorbox{popfigure}[1][]{enhanced,breakable=false,colback=white,colframe=popline,boxrule=1.4pt,arc=8pt,
  left=10pt,right=10pt,top=10pt,bottom=8pt,before skip=10pt,after skip=12pt,halign=center,
  lower separated=false,halign lower=center,
  fontlower=\popsemi\fontsize{9}{11.5}\selectfont\color{popmuted},
  #1}
\newcommand{\legende}[1]{\par\vspace{6pt}{\popsemi\fontsize{9}{11.5}\selectfont\color{popmuted}#1}}

% ============================================================
% Signature OnBuch+
% ============================================================
\newcommand{\poplogoinline}[1]{{\poplogo\fontsize{#1}{#1}\selectfont\color{popink}OnBuch\color{poporange}+}}
\newcommand{\signatureonbuch}{%
  \begin{tcolorbox}[enhanced,colback=popink,colframe=popink,boxrule=2.2pt,arc=10pt,
      drop shadow=poporange,left=14pt,right=14pt,top=10pt,bottom=10pt,before skip=0pt,after skip=0pt]
    \tikz[baseline=-0.7ex]\node[circle,fill=popgreen,draw=popcream,line width=1.3pt,minimum size=22pt,inner sep=0]{\popcheck{white}};%
    \hspace{9pt}\begin{minipage}[c]{0.62\linewidth}
      {\popxbold\fontsize{12}{13}\selectfont\color{popcream}Signé\ \textbullet\ {\poplogo OnBuch\color{poporange}+}}\par\vspace{2pt}
      {\raggedright\popsemi\fontsize{8}{10.5}\selectfont\color{popline}\MakeUppercase{Équipe pédagogique OnBuch+}\\\MakeUppercase{Conforme au programme officiel MINESEC}\par}
    \end{minipage}\hfill
    {\poplogo\fontsize{22}{22}\selectfont\color{popcream}OnBuch\color{poporange}+}
  \end{tcolorbox}%
}

% ============================================================
% Page de garde
% #1 titre · #2 matière · #3 niveau · #4 module · #5 n° leçon · #6 durée ·
% #7 description / objectifs (texte ou itemize)
% ============================================================
\newcommand{\pagegarde}[7]{%
  \thispagestyle{empty}
  \newgeometry{a4paper,margin=1.7cm,top=1.6cm,bottom=1.4cm}%
  \begin{tikzpicture}[remember picture,overlay]
    \fill[poporange!18] ([xshift=-3.2cm,yshift=-3.6cm]current page.north east) circle (4.6cm);
    \fill[poppurple!14] ([xshift=2.4cm,yshift=7.6cm]current page.south west) circle (3.8cm);
  \end{tikzpicture}%
  {\poplogo\fontsize{30}{30}\selectfont\color{popink}OnBuch\color{poporange}+}\hfill
  \raisebox{0.6ex}{\poppill{Cours signé \textbullet\ Premium}{popink}{popcream}}\par
  \vspace{1pt}\popsquiggle{poppurple}\par
  \vspace{8pt}
  \begin{tcolorbox}[enhanced,colback=poporange,colframe=popink,boxrule=2.8pt,arc=13pt,
      drop shadow=popink,left=18pt,right=18pt,top=12pt,bottom=13pt,after skip=10pt]
    \poppill{#2 \textbullet\ #3}{popink}{popcream}\hspace{5pt}\poppill{#4}{white}{popink}\par\vspace{8pt}
    {\popdisplay\fontsize{14}{14}\selectfont\color{popink}LEÇON #5}\par\vspace{4pt}
    {\raggedright\hyphenpenalty=10000\exhyphenpenalty=10000\popdisplay\fontsize{25}{27}\selectfont\color{popcream}\MakeUppercase{#1}\par}
    \vspace{8pt}
    {\popxbold\fontsize{9.5}{9.5}\selectfont\color{popink}\MakeUppercase{Durée conseillée : #6}}
  \end{tcolorbox}
  \begin{tcolorbox}[enhanced,colback=white,colframe=popink,boxrule=2.2pt,arc=10pt,
      drop shadow=popink,left=16pt,right=16pt,top=10pt,bottom=10pt,after skip=8pt]
    \poppill{Dans cette leçon}{poppurple}{white}\par\vspace{5pt}
    \popsemi\fontsize{9.3}{12.6}\selectfont\color{popink}#7
  \end{tcolorbox}
  \noindent\begin{minipage}[t]{0.487\linewidth}
    \begin{tcolorbox}[enhanced,colback=popgreen,colframe=popink,boxrule=2.2pt,arc=10pt,
        drop shadow=popink,left=11pt,right=11pt,top=10pt,bottom=10pt,height=3.1cm,valign=center]
      \tcbox[colback=white,colframe=popink,boxrule=1.3pt,arc=5pt,boxsep=0pt,nobeforeafter,
        left=3pt,right=3pt,top=3pt,bottom=3pt,valign=center]{\qrcode[height=1.9cm]{https://play.google.com/store/apps/details?id=com.luvvix.onbuch&hl=fr}}%
      \hspace{8pt}\begin{minipage}[c]{0.5\linewidth}
        {\popdisplay\fontsize{11}{12.5}\selectfont\color{white}\MakeUppercase{Télécharge OnBuch}}\par\vspace{4pt}
        {\raggedright\popsemi\fontsize{8.4}{11}\selectfont\color{white}Gratuit \textbullet\ Google Play\\Tuteur IA, annales, concours, résultats\par}
      \end{minipage}
    \end{tcolorbox}
  \end{minipage}\hfill
  \begin{minipage}[t]{0.487\linewidth}
    \begin{tcolorbox}[enhanced,colback=poppurple,colframe=popink,boxrule=2.2pt,arc=10pt,
        drop shadow=popink,left=11pt,right=11pt,top=10pt,bottom=10pt,height=3.1cm,valign=center]
      \tikz[baseline=-0.7ex]\node[circle,fill=white,draw=popink,line width=1.3pt,minimum size=1.6cm,inner sep=0]{\popdisplay\fontsize{24}{24}\selectfont\color{poppurple}+};%
      \hspace{8pt}\begin{minipage}[c]{0.6\linewidth}
        {\popdisplay\fontsize{11}{12.5}\selectfont\color{white}\MakeUppercase{Passe à OnBuch+}}\par\vspace{4pt}
        {\raggedright\popsemi\fontsize{8.4}{11}\selectfont\color{white}Tous les cours signés, Tuteur IA illimité, fiches PDF — onglet OnBuch+ de l'app\par}
      \end{minipage}
    \end{tcolorbox}
  \end{minipage}
  \vspace{8pt}
  \popcontenu
  \vfill
  \signatureonbuch
  \restoregeometry
  \newpage
}

% Rangée « Ce cours contient » (page de garde)
\newcommand{\poptile}[3]{% #1 couleur #2 gros texte #3 légende
  \begin{tcolorbox}[enhanced,colback=#1,colframe=popink,boxrule=2pt,arc=9pt,shadow={2.5pt}{-2.5pt}{0pt}{popink},
      left=6pt,right=6pt,top=9pt,bottom=9pt,halign=center,valign=center,height=1.75cm,width=\linewidth,nobeforeafter]
    {\popdisplay\fontsize{12.5}{13}\selectfont\color{white}\MakeUppercase{#2}}\par\vspace{3pt}
    {\centering\popsemi\fontsize{7.8}{9.5}\selectfont\color{white}#3\par}
  \end{tcolorbox}}
\newcommand{\popcontenu}{%
  \par\noindent{\popxboldls\fontsize{7.5}{7.5}\selectfont\color{popmuted}CE COURS SIGNÉ CONTIENT}\par\vspace{7pt}
  \noindent\begin{minipage}[t]{0.235\linewidth}\poptile{popblue}{Cours}{complet, illustré, pas à pas}\end{minipage}\hfill
  \begin{minipage}[t]{0.235\linewidth}\poptile{poporange}{Méthodes}{et exercices résolus}\end{minipage}\hfill
  \begin{minipage}[t]{0.235\linewidth}\poptile{poppink}{Exercices}{type examen + corrigés}\end{minipage}\hfill
  \begin{minipage}[t]{0.235\linewidth}\poptile{popgreen}{Fiche bilan}{l'essentiel à retenir}\end{minipage}\par}

% Sommaire
\newtcolorbox{sommaire}{enhanced,colback=white,colframe=popink,boxrule=2.2pt,arc=10pt,
  drop shadow=popink,left=18pt,right=18pt,top=13pt,bottom=13pt,before skip=6pt,after skip=20pt,
  before upper={\poppill{Sommaire}{poppurple}{white}\par\vspace{9pt}\popsemi\fontsize{10.6}{14.5}\selectfont\color{popink}}}

% Signature de fin de cours — poussée tout en bas de la dernière page
\newcommand{\findecours}{%
  \par\vfill\nopagebreak
  \begin{center}\popsquiggle{poporange}\end{center}\vspace{4pt}
  \signatureonbuch
}
\AtBeginDocument{\def\llbracket{\mathopen{[\mkern-3.5mu[}}\def\rrbracket{\mathclose{]\mkern-3.5mu]}}} % intervalles d'entiers (glyphe ⟦ absent de la police)
% Glyphes absents de Fira Math : redéfinis à partir de glyphes disponibles
\AtBeginDocument{%
  \def\setminus{\mathbin{\backslash}}\let\smallsetminus\setminus
  \def\vdots{\mathord{\vbox{\baselineskip3.2pt\lineskiplimit0pt\kern4pt\hbox{.}\hbox{.}\hbox{.}}}}%
  \def\ddots{\mathinner{\mkern1mu\raise6.5pt\hbox{.}\mkern2mu\raise3.6pt\hbox{.}\mkern2mu\raise0.7pt\hbox{.}\mkern1mu}}%
  \def\triangle{\mathord{\scalebox{0.9}{$\symup{\Delta}$}}}%
  \def\nabla{\mathord{\raisebox{\depth}{\scalebox{1}[-1]{$\symup{\Delta}$}}}}%
  \def\checkmark{\popcheck{popdark}}%
  \def\star{\mathbin{\ast}}\def\bigstar{\mathord{\ast}}%
  \def\diamond{\mathbin{\scalebox{0.6}{$\Diamond$}}}%
  \def\bigcup{\mathop{\mathchoice{\raisebox{-0.3em}{\scalebox{1.7}{$\cup$}}}{\scalebox{1.25}{$\cup$}}{\cup}{\cup}}\displaylimits}%
  \def\bigcap{\mathop{\mathchoice{\raisebox{-0.3em}{\scalebox{1.7}{$\cap$}}}{\scalebox{1.25}{$\cap$}}{\cap}{\cap}}\displaylimits}%
  \def\lesseqqgtr{\mathrel{\lessgtr}}\def\bigoplus{\mathop{\oplus}\displaylimits}%
}
\providecommand{\ssi}{si et seulement si} % abréviation fréquente
\AtBeginDocument{\providecommand{\wideparen}{\overparen}} % arc de cercle

%% FIGFIX-BEGIN v3 — correctifs globaux des figures (docs/onbuch-generation/tools/figfix)
\usepackage{adjustbox}\usepackage{etoolbox}
% toute figure TikZ/pgfplots plus large que la ligne est réduite à la largeur disponible
\BeforeBeginEnvironment{tikzpicture}{\begin{adjustbox}{max width=\linewidth}}
\AfterEndEnvironment{tikzpicture}{\end{adjustbox}}
% légendes pgfplots lisibles par-dessus les courbes
\pgfplotsset{every axis/.append style={legend style={fill=white,fill opacity=0.92,text opacity=1,draw=black!15,inner sep=3pt}}}
% étiquettes d'arêtes (syntaxe edge["..."]) avec fond blanc
\tikzset{every edge quotes/.append style={fill=white,inner sep=1.2pt}}
% bulles de cartes mentales : texte à la ligne dans la bulle, bulle un peu plus grande
\tikzset{every concept/.append style={text width=2.0cm,align=center,minimum size=2.3cm,inner sep=2pt}}
%% FIGFIX-END
\begin{document}

\pagegarde{Les fêtes traditionnelles}{Allemand}{1ère A}{Chapitre 5 : Citoyenneté}{ALA21}{2 h}{Cette leçon de type texte propose la lecture et le commentaire d'un texte allemand sur les fêtes traditionnelles (Oktoberfest, carnaval, fêtes villageoises), enrichi d'un lexique thématique complet. Elle consolide deux points de grammaire essentiels : le datif avec les prépositions de lieu et le passé composé, pour aboutir à une production guidée sur les fêtes traditionnelles camerounaises.
\vspace{4pt}
\begin{multicols}{2}\small
\begin{itemize}
\item À la fin de cette leçon, je sais comprendre globalement puis en détail un texte allemand sur les fêtes traditionnelles
\item À la fin de cette leçon, je sais employer le vocabulaire du thème (das Fest, die Tradition, der Tanz, das Essen, die Kleidung, feiern) avec articles et pluriels
\item À la fin de cette leçon, je sais reconnaître et utiliser le datif après les prépositions de lieu (in, an, auf, zu, bei, mit, von, nach)
\item À la fin de cette leçon, je sais conjuguer le passé composé (Perfekt) avec haben ou sein et former les participes passés
\item À la fin de cette leçon, je sais répondre par des phrases complètes à des questions de compréhension
\item À la fin de cette leçon, je sais résumer un texte en quelques phrases
\end{itemize}\end{multicols}}

\begin{sommaire}
\begin{multicols}{2}
\begin{enumerate}
\item Texte support et première lecture
\item Compréhension du texte
\item Lexique thématique : les fêtes traditionnelles
\item Grammaire 1 : le datif et les prépositions de lieu
\item Grammaire 2 : le passé composé (das Perfekt)
\item Méthode du commentaire et commentaire modèle
\item Production guidée : ma fête traditionnelle
\item Activité d'intégration
\item Méthodes et exercices résolus
\item Exercices
\item Corrigés des exercices
\item Fiche bilan
\end{enumerate}
\end{multicols}
\end{sommaire}

\begin{objectifs}
\begin{itemize}
\item À la fin de cette leçon, je sais comprendre globalement puis en détail un texte allemand sur les fêtes traditionnelles
\item À la fin de cette leçon, je sais employer le vocabulaire du thème (das Fest, die Tradition, der Tanz, das Essen, die Kleidung, feiern) avec articles et pluriels
\item À la fin de cette leçon, je sais reconnaître et utiliser le datif après les prépositions de lieu (in, an, auf, zu, bei, mit, von, nach)
\item À la fin de cette leçon, je sais conjuguer le passé composé (Perfekt) avec haben ou sein et former les participes passés
\item À la fin de cette leçon, je sais répondre par des phrases complètes à des questions de compréhension
\item À la fin de cette leçon, je sais résumer un texte en quelques phrases
\item À la fin de cette leçon, je sais rédiger un court paragraphe sur une fête traditionnelle camerounaise au passé composé
\item À la fin de cette leçon, je sais présenter oralement une fête traditionnelle en 5 à 8 phrases
\end{itemize}
\end{objectifs}

\begin{prerequis}
\begin{itemize}
\item Les articles définis et indéfinis au nominatif et à l'accusatif (Seconde)
\item Le présent des verbes réguliers et irréguliers courants (Seconde)
\item Le vocabulaire de base de la famille, de la nourriture et des loisirs (Seconde)
\item La structure de la phrase allemande : verbe conjugué en 2e position (Seconde)
\item Les auxiliaires haben et sein au présent (début d'année)
\end{itemize}
\end{prerequis}

\newpage

\coursec{Texte support et première lecture}

\courssub{Situation d'accroche et problématique}

En décembre, à Yaoundé et à Douala, les élèves préparent avec fierté les grandes fêtes traditionnelles du pays : le \cle{Ngondo} chez les Sawa, le festival \cle{Nguon} chez les Bamoun. Partout : des danses, des plats traditionnels, des habits colorés, des tambours, des familles réunies. Mais savez-vous que les pays germanophones ont, eux aussi, leurs fêtes traditionnelles ? En Allemagne, on fête l'\all{Oktoberfest} à Munich ; en Autriche, on danse dans les bals de Vienne ; en Suisse, on célèbre la Fête des Vignerons.

\begin{center}
\begin{tabularx}{\linewidth}{|X|X|X|}\hline
\rowcolor{popblueL} Cameroun & Allemagne / Autriche / Suisse & Point commun \\ \hline
le Ngondo (Sawa) & das Oktoberfest (München) & danses, musique \\ \hline
le Nguon (Bamoun) & der Kölner Karneval (Köln) & déguisements, défilés \\ \hline
les danses traditionnelles & die Bälle in Wien (Autriche) & habits traditionnels \\ \hline
les plats de fête & die Fête des Vignerons (Suisse) & repas en famille \\ \hline
\end{tabularx}
\end{center}

\textbf{Problématique :} comment parler de ces fêtes en allemand ? Comment comprendre un texte qui les décrit ? Et comment raconter, au passé, une fête à laquelle on a participé ? Cette leçon nous donne le vocabulaire du thème \all{das Fest} et deux outils grammaticaux essentiels : le \cle{datif} avec les prépositions de lieu et le \cle{passé composé} (\all{das Perfekt}).

\courssub{Présentation du texte : « Traditionelle Feste in Deutschland »}

Le texte que nous allons lire est un court article de magazine jeunesse. Il présente trois types de fêtes traditionnelles allemandes : l'\all{Oktoberfest} à Munich (Bavière), le carnaval de Cologne (\all{der Kölner Karneval}) et les fêtes villageoises comme le \all{Schützenfest} (fête des tireurs), le \all{Weinfest} (fête du vin) et l'\all{Erntedankfest} (fête de la moisson). Avant de lire, observons le titre : \all{traditionelle Feste}. Le mot \all{Fest} est un mot transparent : il ressemble au français « fête ». \all{Traditionell} ressemble à « traditionnel ». Nous pouvons donc déjà deviner le sujet du texte : \emph{les fêtes traditionnelles}.

\begin{textebox}[Traditionelle Feste in Deutschland]
In Deutschland gibt es viele traditionelle Feste. Das bekannteste ist das Oktoberfest in München. Es findet jedes Jahr im September auf der Theresienwiese statt und dauert zwei Wochen. Millionen Besucher aus der ganzen Welt kommen nach München. Sie sitzen in großen Bierzelten, essen Brezeln und Hähnchen und hören traditionelle Musik. Viele tragen Trachten : die Männer Lederhosen, die Frauen Dirndl.

Im Februar feiert man in Köln Karneval. Die Menschen verkleiden sich und tanzen in den Straßen. Beim großen Umzug werfen die Leute Süßigkeiten und Blumen. Auch in den Dörfern gibt es Feste : das Schützenfest, das Weinfest oder das Erntedankfest. Dort trifft sich die ganze Familie, man tanzt, singt und isst zusammen. Diese Feste sind eine wichtige Tradition : sie bringen die Menschen zusammen.
\end{textebox}

\textbf{Glossaire (mots nouveaux) :}

\begin{center}
\begin{tabularx}{\linewidth}{|X|X|}\hline
\rowcolor{popblueL} Deutsch & Français \\ \hline
die Wiese, -n & la prairie \\ \hline
das Bierzelt, -e & la tente à bière \\ \hline
der Umzug, „-e & le défilé, le cortège \\ \hline
die Tracht, -en & le costume traditionnel \\ \hline
das Feuerwerk, -e & le feu d'artifice \\ \hline
stattfinden (findet statt, fand statt, hat stattgefunden) & avoir lieu \\ \hline
dauern & durer \\ \hline
sich verkleiden & se déguiser \\ \hline
der Brauch, „-e & la coutume \\ \hline
die Musikkapelle, -n & la fanfare, l'orchestre \\ \hline
die Süßigkeiten (plur.) & les bonbons, les friandises \\ \hline
zusammenbringen (bringt zusammen) & rassembler, réunir \\ \hline
\end{tabularx}
\end{center}

\begin{attention}[Prononciation des mots nouveaux]
Quelques repères en lettres françaises : \all{die Wiese} « vi-zé », \all{das Bierzelt} « bir-tseult », \all{der Umzug} « oum-tsouk », \all{die Tracht} « trakht », \all{stattfinden} « chtat-fin-den », \all{sich verkleiden} « zikh fèr-klaï-den », \all{die Süßigkeiten} « zu-sikh-kaï-ten ». Rappel : \all{w} se prononce « v », \all{v} se prononce « f », \all{z} se prononce « ts », \all{s} devant voyelle se prononce « z ».
\end{attention}

\courssub{Première lecture : globale, sans dictionnaire}

\begin{methode}[Comment lire un texte allemand en quatre étapes]
\begin{enumerate}
\item \textbf{Lire le titre} et repérer les mots transparents (mots qui ressemblent au français) : \all{Fest}, \all{Tradition}, \all{Musik}, \all{Tanz}, \all{Familie}. Formuler une première hypothèse : « de quoi va parler le texte ? »
\item \textbf{Première lecture globale, sans dictionnaire.} On lit tout le texte d'un trait pour saisir l'idée générale. On ne s'arrête pas sur les mots inconnus.
\item \textbf{Deuxième lecture avec le glossaire.} On relie chaque mot nouveau à sa traduction et on vérifie ses hypothèses.
\item \textbf{Souligner les mots-clés} qui répondent aux questions : \all{wer ?} (qui ?), \all{wo ?} (où ?), \all{wann ?} (quand ?), \all{was macht man ?} (que fait-on ?).
\end{enumerate}
\end{methode}

\textbf{Hypothèses de sens à partir des mots transparents.} Avant même la lecture, les mots \all{Fest}, \all{traditionelle}, \all{Musik}, \all{tanzen}, \all{Familie}, \all{Millionen Besucher} nous permettent de deviner : le texte parle de fêtes, de traditions, de musique et de réunions de famille, avec beaucoup de visiteurs. C'est exactement le contenu du texte : les mots transparents sont nos meilleurs amis !

\begin{attention}[Ne traduis pas mot à mot !]
L'erreur classique du francophone est de traduire chaque mot dans l'ordre. Or l'allemand place souvent le verbe ailleurs : \all{Es findet jedes Jahr im September auf der Theresienwiese statt.} ne se traduit pas « il trouve chaque année en septembre sur la Theresienwiese place », mais « il a lieu chaque année en septembre sur la Theresienwiese ». \textbf{Stratégie :} lire la phrase entière, repérer le verbe (et sa particule séparable à la fin !), puis chercher le sens global.
\end{attention}

\courssub{Deuxième lecture : lecture à voix haute et prononciation}

L'enseignant lit maintenant le texte à voix haute ; les élèves écoutent, puis répètent phrase par phrase. Trois points de prononciation et d'intonation à travailler :

\begin{itemize}
\item \textbf{Les voyelles longues et brèves :} \all{Wiese} (i long, « vi-zé ») contre \all{sitzen} (i bref). La double consonne après une voyelle indique une voyelle brève.
\item \textbf{Le « ch » :} après a, o, u, il se prononce « rkh » (dans la gorge) : \all{Tracht}, \all{auch}, \all{Hähnchen} non — attention, après ä, i, e, il se prononce « hy » doux : \all{Hähnchen} « hèn-hyèn ».
\item \textbf{L'intonation :} dans une phrase déclarative allemande, la voix descend à la fin. Dans les énumérations (\all{man tanzt, singt und isst}), la voix monte sur chaque élément et descend sur le dernier.
\end{itemize}

\courssub{Repérage des informations globales}

Répondons maintenant aux quatre questions fondamentales :

\begin{center}
\begin{tabularx}{\linewidth}{|X|X|}\hline
\rowcolor{popblueL} Frage (question) & Antwort (réponse trouvée dans le texte) \\ \hline
\all{Wovon handelt der Text ?} (de quoi parle le texte ?) & \all{Von traditionellen Festen in Deutschland.} \\ \hline
\all{Welche Feste ?} (quelles fêtes ?) & \all{das Oktoberfest, der Karneval in Köln, das Schützenfest, das Weinfest, das Erntedankfest} \\ \hline
\all{Wo ?} (où ?) & \all{in München, in Köln, in den Dörfern} \\ \hline
\all{Wann ?} (quand ?) & \all{im September (Oktoberfest), im Februar (Karneval)} \\ \hline
\all{Was macht man ?} (que fait-on ?) & \all{man isst, trinkt, tanzt, singt, verkleidet sich} \\ \hline
\end{tabularx}
\end{center}

\begin{popfigure}
\begin{tikzpicture}[mindmap, grow cyclic, every node/.style={concept, align=center, font=\small}, root concept/.append style={concept color=popblue, font=\bfseries}, level 1/.append style={level distance=3.2cm, sibling angle=90}]
\node[root concept]{das Fest}
  child[concept color=poporange]{ node[align=center]{wo ?\\ München, Köln,\\ die Dörfer} }
  child[concept color=popgreen]{ node[align=center]{wann ?\\ im September,\\ im Februar} }
  child[concept color=poppurple]{ node[align=center]{wer ?\\ die Besucher,\\ die Familie,\\ die Leute} }
  child[concept color=poppink]{ node[align=center]{was macht\\ man ?\\ essen, trinken,\\ tanzen, singen,\\ sich verkleiden} };
\end{tikzpicture}
\legende{Carte mentale : les quatre questions pour comprendre un texte sur une fête.}
\end{popfigure}

\begin{popfigure}
\begin{tikzpicture}[scale=1.0]
\draw[fill=popcream, draw=popline, rounded corners=6pt] (0,0) .. controls (0.4,1.2) and (-0.2,2.4) .. (0.6,3.4) .. controls (1.2,4.2) and (1.0,5.0) .. (1.8,5.6) .. controls (2.8,6.2) and (4.0,6.0) .. (4.6,5.2) .. controls (5.2,4.4) and (4.8,3.4) .. (5.2,2.4) .. controls (5.6,1.4) and (4.8,0.4) .. (3.8,0.1) .. controls (2.6,-0.2) and (1.0,-0.3) .. (0,0) -- cycle;
\node[font=\bfseries\small, popdark] at (2.6,6.4) {Deutschland};
\fill[poporange] (1.5,4.6) circle (0.09);
\node[right, font=\small, align=left] at (1.7,4.6) {\textbf{Köln}\\ Karneval (Februar)};
\fill[poppurple] (3.9,0.9) circle (0.09);
\node[right, font=\small, align=left] at (4.1,0.9) {\textbf{München}\\ Oktoberfest (September)};
\draw[->, thick, popmuted] (1.5,4.4) -- (3.7,1.1);
\node[font=\scriptsize, popmuted, align=center] at (2.2,2.6) {ca. 570 km};
\end{tikzpicture}
\legende{Mini-carte de l'Allemagne : Köln (carnaval) au nord-ouest, München (Oktoberfest) au sud-est.}
\end{popfigure}

\begin{savaistu}[L'Oktoberfest et la « cinquième saison »]
L'Oktoberfest commence en réalité en \textbf{septembre} et non en octobre : on l'a avancé pour profiter d'un temps plus doux. Il attire environ 6 millions de visiteurs chaque année sur la \all{Theresienwiese} (« la prairie de Thérèse »), du nom de la princesse Thérèse de Saxe-Hildburghausen, dont le mariage en 1810 est à l'origine de la fête. Le carnaval de Cologne, lui, est surnommé \all{die fünfte Jahreszeit} — « la cinquième saison » — tant il occupe la vie de la ville ! Et au Cameroun ? Le Ngondo rassemble chaque année les Sawa sur les bords du Wouri : comme en Allemagne, la fête traditionnelle est un moment où la communauté se retrouve.
\end{savaistu}

\courssub{Vérification de la compréhension globale}

\begin{exoresolu}[Richtig oder falsch ?]
Énoncé : lis les phrases suivantes et dis si elles sont vraies (\all{richtig}) ou fausses (\all{falsch}) d'après le texte. Corrige les phrases fausses.
\begin{enumerate}
\item \all{Das Oktoberfest findet im Oktober statt.}
\item \all{Die Besucher essen Brezeln und hören Musik.}
\item \all{In Köln verkleiden sich die Menschen.}
\item \all{In den Dörfern gibt es keine Feste.}
\end{enumerate}
\tcblower
Solution détaillée :
\begin{enumerate}
\item \textbf{Falsch.} Le texte dit : \all{Es findet jedes Jahr im September ... statt.} La fête a lieu en \emph{septembre}, pas en octobre. Piège classique : le nom \all{Oktoberfest} trompe les francophones !
\item \textbf{Richtig.} Le texte dit : \all{Sie ... essen Brezeln und Hähnchen und hören traditionelle Musik.}
\item \textbf{Richtig.} Le texte dit : \all{Die Menschen verkleiden sich und tanzen in den Straßen.}
\item \textbf{Falsch.} Le texte dit : \all{Auch in den Dörfern gibt es Feste : das Schützenfest, das Weinfest oder das Erntedankfest.} Il y a \emph{aussi} des fêtes dans les villages.
\end{enumerate}
\end{exoresolu}

\begin{exercice}[À toi de jouer : premières questions]
Réponds en allemand, par des phrases complètes et simples :
\begin{enumerate}
\item \all{Wo findet das Oktoberfest statt ?}
\item \all{Wie lange dauert das Oktoberfest ?}
\item \all{Was tragen die Männer und die Frauen ?}
\item \all{Wann feiert man in Köln Karneval ?}
\item \all{Was werfen die Leute beim Umzug ?}
\end{enumerate}
\end{exercice}

\begin{corrige}[Exercice « À toi de jouer »]
\begin{enumerate}
\item \all{Das Oktoberfest findet in München statt, auf der Theresienwiese.} (Il a lieu à Munich, sur la Theresienwiese.)
\item \all{Es dauert zwei Wochen.} (Il dure deux semaines.)
\item \all{Die Männer tragen Lederhosen, die Frauen tragen Dirndl.} (Les hommes portent des culottes de cuir, les femmes des dirndl.)
\item \all{Man feiert im Februar Karneval.} (On fête le carnaval en février.)
\item \all{Die Leute werfen Süßigkeiten und Blumen.} (Les gens jettent des bonbons et des fleurs.)
\end{enumerate}
Remarque de méthode : dans chaque réponse, le \textbf{verbe conjugué est en deuxième position} — c'est la règle d'or de la phrase allemande.
\end{corrige}

\begin{aretenir}[Ce qu'il faut retenir de cette première étape]
\begin{itemize}
\item Un texte se lit en \textbf{quatre étapes} : titre et mots transparents, lecture globale, lecture avec glossaire, repérage des mots-clés (\all{wer, wo, wann, was}).
\item Les \textbf{mots transparents} (\all{Fest, Tradition, Musik, Familie}) permettent de deviner le sens avant la lecture.
\item On ne traduit \textbf{jamais mot à mot} : on repère le verbe et le sens global de la phrase.
\item Le texte présente trois types de fêtes : l'\all{Oktoberfest} (München, septembre), le \all{Karneval} (Köln, février) et les fêtes villageoises (\all{Schützenfest, Weinfest, Erntedankfest}).
\end{itemize}
\end{aretenir}

Maintenant que le sens global du texte est clair, nous passerons à la compréhension détaillée, puis au lexique thématique et aux deux outils grammaticaux de la leçon : le datif avec les prépositions de lieu et le passé composé, qui nous permettront de raconter notre propre fête traditionnelle.

\coursec{Compréhension du texte}

Dans la section précédente, nous avons découvert le texte support et effectué une première lecture globale : nous savons déjà qu'il parle des \cle{fêtes traditionnelles} en Allemagne. Il est temps maintenant de passer à la \cle{compréhension détaillée} : lire ligne par ligne, repérer les informations précises et apprendre à répondre à des questions par des phrases complètes et correctes. C'est une compétence essentielle pour l'examen : elle sera évaluée à l'écrit comme à l'oral.

\courssub{Rappel du texte support}

Avant de répondre aux questions, relisons attentivement le texte. Pour faciliter ton travail, les lignes sont repérées par paragraphes.

\begin{textebox}[Traditionelle Feste in Deutschland]
Deutschland ist ein Land mit vielen traditionellen Festen. Jedes Fest hat seine Geschichte, seine Musik, sein Essen und seine Kleidung.

Das berühmteste Fest ist das Oktoberfest in München. Es findet jedes Jahr im September statt und dauert etwa zwei Wochen. Die Besucher sitzen in großen Zelten an langen Tischen. Sie essen Brezeln und Hähnchen und trinken Bier. Die Männer tragen Lederhosen, die Frauen tragen das Dirndl, ein traditionelles Kleid. Viele Touristen aus aller Welt kommen nach München, um dieses Fest zu erleben.

Im Februar feiert man in Köln Karneval. Die Menschen verkleiden sich und tanzen in den Straßen. Beim großen Umzug am Rosenmontag werfen die Leute Bonbons und Blumen von den Wagen. Kinder und Erwachsene lachen und singen zusammen.

In vielen Dörfern gibt es auch kleinere Feste, zum Beispiel das Schützenfest im Sommer. Die Schützen schießen auf eine Figur aus Holz, und der beste Schütze wird König des Festes. Am Abend essen die Familien zusammen und tanzen bis in die Nacht.

Diese Feste sind sehr wichtig für die Menschen : sie bringen die Familien und die Nachbarn zusammen, und sie geben die Traditionen von den Großeltern an die Kinder weiter. So bleibt die Kultur lebendig.
\end{textebox}

\textbf{Glossaire (Wortschatz)}

\begin{center}
\begin{tabularx}{\linewidth}{|X|X|X|}
\hline
\rowcolor{popblueL} Allemand & Forme / Remarque & Français \\
\hline
das Fest & die Feste & la fête \\
\hline
die Tradition & die Traditionen & la tradition \\
\hline
berühmt & adjectif & célèbre \\
\hline
stattfinden (findet statt) & verbe séparable & avoir lieu \\
\hline
dauern & verbe régulier & durer \\
\hline
das Zelt & die Zelte & la tente, le chapiteau \\
\hline
die Brezel & die Brezeln & le bretzel \\
\hline
das Hähnchen & die Hähnchen & le poulet (rôti) \\
\hline
die Lederhose & die Lederhosen & la culotte de cuir (bavaroise) \\
\hline
das Dirndl & die Dirndl & la robe traditionnelle bavaroise \\
\hline
sich verkleiden & verbe pronominal & se déguiser \\
\hline
der Umzug & die Umzüge & le défilé, la procession \\
\hline
der Wagen & die Wagen & le char (du défilé) \\
\hline
das Schützenfest & die Schützenfeste & la fête des archers / des tireurs \\
\hline
schießen auf & verbe fort (schießt, schoss, hat geschossen) & tirer sur \\
\hline
die Figur aus Holz & die Figuren & la figurine en bois \\
\hline
der Nachbar & die Nachbarn & le voisin \\
\hline
weitergeben (gibt weiter) & verbe séparable & transmettre \\
\hline
lebendig & adjectif & vivant \\
\hline
\end{tabularx}
\end{center}

\courssub{Compréhension globale du texte}

Avant les questions de détail, l'examinateur vérifie toujours que tu as compris \emph{l'essentiel}. Trois tâches classiques :

\begin{enumerate}
\item \textbf{Choisir un titre.} Le titre doit résumer le sujet général, pas un détail. Ici : \all{Traditionelle Feste in Deutschland} est le meilleur titre. Un titre comme \all{Das Oktoberfest} serait trop étroit : le texte parle aussi du Carnaval et du Schützenfest.
\item \textbf{Résumer en une phrase.} Une seule phrase, avec un sujet, un verbe en deuxième position et les informations essentielles. Modèle : \all{Der Text beschreibt drei traditionelle Feste in Deutschland : das Oktoberfest, den Karneval und das Schützenfest.} « Le texte décrit trois fêtes traditionnelles en Allemagne : l'Oktoberfest, le Carnaval et le Schützenfest. »
\item \textbf{Identifier le thème.} Le thème est : \all{die traditionellen Feste und ihre Bedeutung für die Menschen} « les fêtes traditionnelles et leur importance pour les gens ».
\end{enumerate}

\begin{methode}[Répondre à une question de compréhension en trois étapes]
\begin{enumerate}
\item \textbf{Repérer} la ligne ou le paragraphe du texte qui contient la réponse (souligne-la au brouillon).
\item \textbf{Reformuler} avec tes propres mots : ne recopie jamais tout le paragraphe, prends seulement l'information demandée.
\item \textbf{Rédiger} une phrase complète en allemand : sujet + verbe conjugué en \cle{deuxième position} + compléments. Reprends les mots de la question dans ta réponse.
\end{enumerate}
\end{methode}

\begin{propriete}[La technique de la réponse complète]
En allemand comme en français, une bonne réponse \cle{reprend les mots de la question}. La question donne déjà la structure de la phrase :

\all{Wann findet das Oktoberfest statt ?} → \all{Das Oktoberfest findet im September statt.}

On remplace le mot interrogatif (\all{wann}, \all{wo}, \all{was}, \all{wer}) par l'information du texte, et le verbe revient en deuxième position. C'est la transformation question → réponse.
\end{propriete}

\begin{popfigure}
\begin{tikzpicture}[node distance=0.6cm and 1.2cm]
\node[draw=popblue, fill=popblueL, rounded corners, align=center, text width=5.6cm] (q) {\all{Wann findet das Oktoberfest statt ?}\\ \footnotesize question : verbe en 2\up{e} position};
\node[draw=popgreen, fill=popgreenL, rounded corners, align=center, text width=6.4cm, right=of q] (r) {\all{Das Oktoberfest findet im September statt.}\\ \footnotesize réponse : sujet d'abord, verbe en 2\up{e} position};
\draw[-{Stealth[length=3mm]}, very thick, poporange] (q) -- (r) node[midway, above, align=center, font=\footnotesize, text=popdark]{on remplace\\ \all{wann} par la date};
\end{tikzpicture}
\legende{De la question à la réponse complète : le verbe reste en deuxième position.}
\end{popfigure}

\courssub{Questions de compréhension détaillée}

Voici les questions types de l'examen, avec leurs réponses modèles. Observe comment chaque réponse reprend les mots de la question.

\begin{enumerate}
\item \all{Wann findet das Oktoberfest statt ?}\\
→ \all{Das Oktoberfest findet jedes Jahr im September statt.} « L'Oktoberfest a lieu chaque année en septembre. »
\item \all{Wo sitzen die Besucher ?}\\
→ \all{Die Besucher sitzen in großen Zelten an langen Tischen.} « Les visiteurs sont assis dans de grandes tentes, à de longues tables. »
\item \all{Was essen und trinken sie ?}\\
→ \all{Sie essen Brezeln und Hähnchen und trinken Bier.} « Ils mangent des bretzels et du poulet et boivent de la bière. »
\item \all{Was tragen die Männer und die Frauen ?}\\
→ \all{Die Männer tragen Lederhosen und die Frauen tragen das Dirndl.} « Les hommes portent la culotte de cuir et les femmes portent le dirndl. »
\item \all{Wann feiert man in Köln Karneval ?}\\
→ \all{In Köln feiert man im Februar Karneval.} « À Cologne, on fête le Carnaval en février. »
\item \all{Was machen die Leute beim Umzug ?}\\
→ \all{Beim Umzug werfen die Leute Bonbons und Blumen von den Wagen.} « Pendant le défilé, les gens jettent des bonbons et des fleurs depuis les chars. »
\item \all{Welche Feste gibt es in den Dörfern ?}\\
→ \all{In den Dörfern gibt es kleinere Feste, zum Beispiel das Schützenfest.} « Dans les villages, il y a des fêtes plus petites, par exemple le Schützenfest. »
\item \all{Warum sind diese Feste wichtig ?}\\
→ \all{Diese Feste sind wichtig, weil sie die Familien und die Nachbarn zusammenbringen und die Traditionen weitergeben.} « Ces fêtes sont importantes parce qu'elles rassemblent les familles et les voisins et transmettent les traditions. »
\end{enumerate}

\begin{exemplebox}[Observer la transformation question → réponse]
\all{Wann findet das Oktoberfest statt ?} « Quand a lieu l'Oktoberfest ? »\\
→ \all{Es findet jedes Jahr im September statt.} « Il a lieu chaque année en septembre. »

\all{Was machen die Menschen beim Karneval ?} « Que font les gens au Carnaval ? »\\
→ \all{Die Menschen verkleiden sich und tanzen in den Straßen.} « Les gens se déguisent et dansent dans les rues. » (verbe pronominal \all{sich verkleiden} = se déguiser ; le pronom réfléchi \all{sich} se place juste après le verbe conjugué.)
\end{exemplebox}

\begin{attention}[Le verbe séparable \all{stattfinden}]
\all{stattfinden} « avoir lieu » est un verbe à \cle{particule séparable}. Au présent, la particule \all{statt} va \emph{à la fin de la phrase} :

\all{Das Fest findet im September statt.} « La fête a lieu en septembre. »

Erreur typique des francophones : écrire \emph{\all{Das Fest stattfindet im September}} — c'est FAUX. La particule se détache toujours du verbe conjugué. Autre exemple : \all{Der Umzug findet am Montag statt.} « Le défilé a lieu le lundi. »
\end{attention}

\courssub{Vrai ou faux justifié (Richtig / Falsch)}

À l'examen, chaque réponse \all{richtig} ou \all{falsch} doit être \cle{justifiée par une citation du texte} entre guillemets allemands. Sans citation, la réponse ne vaut rien.

\begin{exoresolu}[Richtig oder falsch ?]
Énoncé. Entoure \all{richtig} ou \all{falsch} et justifie avec une phrase du texte :
\begin{enumerate}
\item \all{Das Oktoberfest dauert einen Monat.}
\item \all{Die Frauen tragen beim Oktoberfest das Dirndl.}
\item \all{In Köln feiert man Karneval im Sommer.}
\item \all{Der beste Schütze wird König des Festes.}
\end{enumerate}
\tcblower
Solution détaillée :
\begin{enumerate}
\item \all{Falsch.} Citation : „Es findet jedes Jahr im September statt und dauert etwa zwei Wochen.“ Le texte dit \emph{deux semaines}, pas un mois.
\item \all{Richtig.} Citation : „die Frauen tragen das Dirndl, ein traditionelles Kleid.“
\item \all{Falsch.} Citation : „Im Februar feiert man in Köln Karneval.“ Le Carnaval de Cologne a lieu en \emph{février}, c'est-à-dire en hiver, pas en été.
\item \all{Richtig.} Citation : „der beste Schütze wird König des Festes.“
\end{enumerate}
Méthode : pour chaque phrase, je repère d'abord le mot-clé (\all{dauern}, \all{Dirndl}, \all{Köln}, \all{Schütze}), je retrouve la ligne du texte, je compare, puis je cite la phrase exacte entre guillemets „…“.
\end{exoresolu}

\courssub{Exercice de correspondance : fête, mois, ville, activité}

Un exercice très fréquent consiste à relier les informations entre elles. Voici le tableau récapitulatif du texte — mémorise-le, il te servira aussi pour ta production écrite.

\begin{center}
\begin{tabularx}{\linewidth}{|X|X|X|X|}
\hline
\rowcolor{popblueL} das Fest & der Ort & die Zeit & die Aktivitäten \\
\hline
Oktoberfest & München & September & essen, trinken, in Zelten sitzen, Trachten tragen \\
\hline
Karneval & Köln & Februar & sich verkleiden, tanzen, Bonbons werfen \\
\hline
Schützenfest & die Dörfer & Sommer & schießen, essen, bis in die Nacht tanzen \\
\hline
\end{tabularx}
\end{center}

\begin{exercice}[À toi de jouer : complète le tableau]
Complète le tableau ci-dessous avec les informations du texte (certaines cases sont déjà remplies pour t'aider).
\end{exercice}

\begin{popfigure}
\begin{tikzpicture}
% Dimensions : largeur 12cm, hauteur 3.2cm (4 lignes x 0.8cm)
% Colonnes : 0, 3, 6, 9, 12
% Lignes : 0, 0.8, 1.6, 2.4, 3.2
\draw[thick, popdark] (0,0) rectangle (12,3.2);
% Lignes horizontales
\foreach \y in {0.8, 1.6, 2.4} {
    \draw[thick, popdark] (0,\y) -- (12,\y);
}
% Lignes verticales
\foreach \x in {3, 6, 9} {
    \draw[thick, popdark] (\x,0) -- (\x,3.2);
}
% En-tête
\fill[popblueL] (0,2.4) rectangle (12,3.2);
\node[font=\bfseries, align=center] at (1.5,2.8) {das Fest};
\node[font=\bfseries, align=center] at (4.5,2.8) {der Ort};
\node[font=\bfseries, align=center] at (7.5,2.8) {die Zeit};
\node[font=\bfseries, align=center] at (10.5,2.8) {die Aktivitäten};
% Ligne 1 : Oktoberfest
\node[align=center] at (1.5,2.0) {Oktoberfest};
\node[align=center] at (4.5,2.0) {München};
\node[align=center, text=poporange] at (7.5,2.0) {?};
\node[align=center, text=poporange] at (10.5,2.0) {?};
% Ligne 2 : Karneval
\node[align=center] at (1.5,1.2) {Karneval};
\node[align=center, text=poporange] at (4.5,1.2) {?};
\node[align=center] at (7.5,1.2) {Februar};
\node[align=center, text=poporange] at (10.5,1.2) {?};
% Ligne 3 : Schützenfest
\node[align=center] at (1.5,0.4) {Schützenfest};
\node[align=center, text=poporange] at (4.5,0.4) {?};
\node[align=center, text=poporange] at (7.5,0.4) {?};
\node[align=center] at (10.5,0.4) {schießen, tanzen};
\end{tikzpicture}
\legende{Tableau à compléter : relie chaque fête à son lieu, sa date et ses activités.}
\end{popfigure}

\begin{corrige}[Exercice « À toi de jouer »]
\begin{itemize}
\item Oktoberfest : \textbf{Zeit} = \all{September} ; \textbf{Aktivitäten} = \all{Brezeln und Hähnchen essen, Bier trinken, Trachten tragen}.
\item Karneval : \textbf{Ort} = \all{Köln} ; \textbf{Aktivitäten} = \all{sich verkleiden, in den Straßen tanzen, Bonbons und Blumen werfen}.
\item Schützenfest : \textbf{Ort} = \all{die Dörfer} ; \textbf{Zeit} = \all{Sommer}.
\end{itemize}
Vérification : chaque réponse se trouve dans le texte — relis les paragraphes 2, 3 et 4 pour contrôler.
\end{corrige}

\begin{exercice}[Questions personnelles (préparation orale)]
Réponds par des phrases complètes, en reprenant les mots de la question :
\begin{enumerate}
\item \all{Welches Fest findest du interessant ? Warum ?}
\item \all{Gibt es in Kamerun ähnliche Feste ? Welche ?}
\item \all{Was isst man bei Festen in deiner Familie ?}
\end{enumerate}
\end{exercice}

\begin{attention}[Pièges fréquents dans les réponses de compréhension]
\begin{itemize}
\item \textbf{Recopier tout le paragraphe} au lieu de sélectionner l'information demandée : l'examinateur sanctionne le hors-sujet.
\item \textbf{Oublier la deuxième position du verbe} : après un complément en tête de phrase, le verbe vient tout de suite : \all{In Köln feiert man im Februar Karneval.} (et non \emph{\all{In Köln man feiert...}}).
\item \textbf{Traduire mot à mot depuis le français} : « avoir lieu » ne se traduit jamais par \emph{\all{haben}} ; il faut le verbe séparable \all{stattfinden}.
\item \textbf{Oublier la majuscule} aux noms : \all{das Oktoberfest}, \all{der Karneval}, \all{die Brezeln}.
\end{itemize}
\end{attention}

\begin{aretenir}[L'essentiel de la section]
\begin{itemize}
\item Compréhension globale : titre général + résumé en une phrase + thème.
\item Compréhension détaillée : repérer la ligne, reformuler, rédiger une phrase complète avec le verbe en deuxième position.
\item La réponse reprend les mots de la question ; le mot interrogatif est remplacé par l'information du texte.
\item \all{Richtig / Falsch} : toujours justifier par une citation exacte entre guillemets „…“.
\item \all{stattfinden} est séparable : \all{Das Fest findet im September statt.}
\end{itemize}
\end{aretenir}

Maintenant que tu comprends le texte en détail, passons à son lexique : dans la section suivante, nous étudierons systématiquement le vocabulaire des fêtes traditionnelles, avec les articles, les pluriels et les familles de mots.

\coursec{Lexique thématique : les fêtes traditionnelles}

Dans la section précédente, nous avons lu et compris le texte support sur les fêtes traditionnelles. Pour bien parler de ce thème, à l'écrit comme à l'oral, il faut maintenant en maîtriser le vocabulaire de manière précise : chaque mot avec son \cle{article} (der, die, das), son \cle{pluriel} et sa prononciation. C'est l'objectif de cette section : construire, champ par champ, un lexique solide et utilisable immédiatement dans vos réponses aux questions, vos résumés et vos productions écrites.

\begin{propriete}[La règle d'or du vocabulaire allemand]
En allemand, on n'apprend \textbf{jamais} un nom seul. On l'apprend toujours avec trois informations :
\begin{enumerate}
\item son \cle{article} : \all{der} (masculin), \all{die} (féminin), \all{das} (neutre) ;
\item son \cle{pluriel} : \all{das Fest, die Feste} ; \all{der Tanz, die Tänze} ; \all{die Tradition, die Traditionen} ;
\item un \cle{exemple} en contexte : \all{Wir feiern das Fest mit der Familie.}
\end{enumerate}
Le genre d'un nom allemand est imprévisible pour un francophone : \all{das Fest} est neutre alors que « la fête » est féminin en français. Seule la mémorisation systématique de l'article évite les erreurs de déclinaison (au datif : \all{auf dem Fest}, et non \all{auf der Fest}).
\end{propriete}

\courssub{Champ 1 : la fête et la tradition}

Observons d'abord ces phrases tirées de situations authentiques :

\all{Das Oktoberfest findet jedes Jahr in München statt.} « L'Oktoberfest a lieu chaque année à Munich. »

\all{Die Feier dauert drei Tage.} « La fête dure trois jours. »

\all{Das ist eine alte Tradition in unserem Dorf.} « C'est une vieille tradition dans notre village. »

\begin{center}
\begin{tabularx}{\linewidth}{|l|l|X|}\hline
\rowcolor{popblueL} \textbf{Allemand} & \textbf{Français} & \textbf{Exemple} \\ \hline
das Fest, die Feste & la fête & \all{Das Fest beginnt um 18 Uhr.} \\ \hline
die Feier, die Feiern & la fête, la cérémonie & \all{Die Feier ist sehr schön.} \\ \hline
die Tradition, die Traditionen & la tradition & \all{Jede Region hat ihre Traditionen.} \\ \hline
der Brauch, die Bräuche & la coutume & \all{Ein alter Brauch in Bayern.} \\ \hline
der Feiertag, die Feiertage & le jour férié & \all{Der 3. Oktober ist ein Feiertag.} \\ \hline
feiern (feierte, hat gefeiert) & fêter, célébrer & \all{Wir feiern zusammen.} \\ \hline
stattfinden (fand statt, hat stattgefunden) & avoir lieu & \all{Das Fest findet im September statt.} \\ \hline
dauern (dauerte, hat gedauert) & durer & \all{Die Feier dauert bis Mitternacht.} \\ \hline
\end{tabularx}
\end{center}

\begin{attention}[das Fest ou fest ?]
Ne confondez pas \all{das Fest} (nom neutre : la fête) et \all{fest} (adjectif : solide, ferme). Comparez : \all{Wir besuchen das Fest.} « Nous visitons la fête. » et \all{Der Tisch ist fest.} « La table est solide. » La majuscule fait toute la différence : en allemand, \textbf{tous les noms prennent une majuscule}, les adjectifs non. Autre piège : \all{stattfinden} est un verbe à \cle{particule séparable} : \all{Das Fest findet morgen statt.} (la particule \all{statt} va en fin de phrase).
\end{attention}

\begin{exemplebox}[feiern et ses amis]
\all{Wir feiern das Fest mit der ganzen Familie.} « Nous fêtons la fête avec toute la famille. »

\all{Der Feiertag ist am Montag.} « Le jour férié est lundi. »

\all{In Douala findet jedes Jahr ein großes Fest statt.} « À Douala, une grande fête a lieu chaque année. »
\end{exemplebox}

\courssub{Champ 2 : la danse et la musique}

\begin{center}
\begin{tabularx}{\linewidth}{|l|l|X|}\hline
\rowcolor{popgreenL} \textbf{Allemand} & \textbf{Français} & \textbf{Exemple} \\ \hline
der Tanz, die Tänze & la danse & \all{Der Tanz ist sehr alt.} \\ \hline
die Musik (meist Sg.) & la musique & \all{Die Musik ist laut und fröhlich.} \\ \hline
das Lied, die Lieder & la chanson & \all{Alle singen ein altes Lied.} \\ \hline
die Musikkapelle, die Musikkapellen & la fanfare & \all{Die Musikkapelle spielt im Zelt.} \\ \hline
der Tänzer, die Tänzer / die Tänzerin, -nen & le danseur / la danseuse & \all{Die Tänzer tragen Masken.} \\ \hline
tanzen & danser & \all{Die Jugendlichen tanzen die ganze Nacht.} \\ \hline
singen (sang, hat gesungen) & chanter & \all{Wir singen traditionelle Lieder.} \\ \hline
spielen & jouer (d'un instrument) & \all{Sie spielt Gitarre.} \\ \hline
\end{tabularx}
\end{center}

\begin{exemplebox}[En contexte]
\all{Die Musikkapelle spielt, und alle Gäste tanzen.} « La fanfare joue, et tous les invités dansent. »

\all{Die Kinder singen ein traditionelles Lied.} « Les enfants chantent une chanson traditionnelle. »
\end{exemplebox}

\courssub{Champ 3 : la nourriture et les boissons}

\begin{center}
\begin{tabularx}{\linewidth}{|l|l|X|}\hline
\rowcolor{poporangeL} \textbf{Allemand} & \textbf{Français} & \textbf{Exemple} \\ \hline
das Essen (meist Sg.) & le repas, la nourriture & \all{Das Essen schmeckt sehr gut.} \\ \hline
das Getränk, die Getränke & la boisson & \all{Die Getränke sind kalt.} \\ \hline
das Bier, die Biere & la bière & \all{Die Männer trinken Bier.} \\ \hline
die Brezel, die Brezeln & le bretzel & \all{Ich kaufe eine Brezel.} \\ \hline
das Hähnchen, die Hähnchen & le poulet (rôti) & \all{Ein Hähnchen mit Kartoffeln.} \\ \hline
der Wein, die Weine & le vin & \all{Der Wein kommt aus Österreich.} \\ \hline
die Süßigkeiten (Pl.) & les bonbons, les sucreries & \all{Die Kinder lieben Süßigkeiten.} \\ \hline
essen (aß, hat gegessen) & manger & \all{Wir essen zusammen.} \\ \hline
trinken (trank, hat getrunken) & boire & \all{Sie trinkt Wasser.} \\ \hline
\end{tabularx}
\end{center}

\begin{attention}[Verbes irréguliers à mémoriser]
\all{essen} et \all{trinken} sont des verbes forts : au présent, \all{essen} change sa voyelle (\all{ich esse, du isst, er isst}) ; au passé composé : \all{ich habe gegessen}, \all{ich habe getrunken}. Ne dites jamais \all{ich habe geesst} ! Notez aussi le participe passé avec \all{ge-} et la voyelle changée : \all{essen → gegessen}, \all{trinken → getrunken}.
\end{attention}

\courssub{Champ 4 : les vêtements}

\begin{definition}[die Tracht, die Trachten]
La \all{Tracht} est le vêtement traditionnel porté lors des fêtes régionales dans les pays germanophones. En Bavière et en Autriche, les femmes portent le \all{Dirndl} (robe traditionnelle avec tablier) et les hommes les \all{Lederhosen} (culotte de cuir). Chaque région a ses propres couleurs et motifs.
\end{definition}

\begin{center}
\begin{tabularx}{\linewidth}{|l|l|X|}\hline
\rowcolor{poppurpleL} \textbf{Allemand} & \textbf{Français} & \textbf{Exemple} \\ \hline
die Kleidung (meist Sg.) & les vêtements & \all{Die Kleidung ist bunt.} \\ \hline
die Tracht, die Trachten & le costume traditionnel & \all{Sie trägt eine Tracht.} \\ \hline
das Dirndl, die Dirndl & le dirndl & \all{Das Dirndl ist rot und weiß.} \\ \hline
die Lederhosen (Pl.) & la culotte de cuir & \all{Er trägt Lederhosen.} \\ \hline
das Kostüm, die Kostüme & le costume, le déguisement & \all{Die Kinder tragen Kostüme.} \\ \hline
tragen (trug, hat getragen) & porter (vêtement) & \all{Sie trägt ein Dirndl.} \\ \hline
sich verkleiden & se déguiser & \all{Wir verkleiden uns.} \\ \hline
(sich) anziehen & (s')habiller & \all{Er zieht die Tracht an.} \\ \hline
\end{tabularx}
\end{center}

\begin{attention}[die Kleidung : singulier en allemand, pluriel en français]
En français on dit « les vêtements » (pluriel), mais en allemand \all{die Kleidung} est un nom \textbf{singulier} : \all{Die Kleidung ist schön.} « Les vêtements sont beaux. » Le verbe reste au singulier. Autres pièges : \all{tragen} (porter un vêtement) est irrégulier : \all{ich trage, du trägst, er trägt} ; \all{anziehen} est séparable : \all{Ich ziehe die Tracht an.} ; \all{sich verkleiden} est réfléchi : \all{Ich verkleide mich.}
\end{attention}

\begin{exemplebox}[En contexte]
\all{Auf dem Fest tragen viele Menschen Trachten.} « À la fête, beaucoup de gens portent des costumes traditionnels. »

\all{Die Frauen ziehen ihre Dirndl an, die Männer ihre Lederhosen.} « Les femmes enfilent leur dirndl, les hommes leur culotte de cuir. »
\end{exemplebox}

\courssub{Champ 5 : les activités et les personnes}

\begin{center}
\begin{tabularx}{\linewidth}{|l|l|X|}\hline
\rowcolor{popgoldL} \textbf{Allemand} & \textbf{Français} & \textbf{Exemple} \\ \hline
der Umzug, die Umzüge & le défilé, la procession & \all{Der Umzug beginnt um 10 Uhr.} \\ \hline
das Feuerwerk, die Feuerwerke & le feu d'artifice & \all{Das Feuerwerk ist spektakulär.} \\ \hline
der Besucher, die Besucher & le visiteur & \all{Viele Besucher kommen aus Europa.} \\ \hline
der Gast, die Gäste & l'invité & \all{Die Gäste essen und trinken.} \\ \hline
die Familie, die Familien & la famille & \all{Die ganze Familie ist da.} \\ \hline
sich treffen (traf, hat getroffen) & se retrouver & \all{Wir treffen uns am Marktplatz.} \\ \hline
werfen (warf, hat geworfen) & lancer, jeter & \all{Die Kinder werfen Süßigkeiten.} \\ \hline
zusammenbringen (sép.) & réunir & \all{Das Fest bringt alle zusammen.} \\ \hline
\end{tabularx}
\end{center}

\begin{exemplebox}[En contexte]
\all{Beim Umzug werfen die Kinder Süßigkeiten.} « Pendant le défilé, les enfants lancent des bonbons. »

\all{Das Fest bringt die ganze Familie zusammen.} « La fête réunit toute la famille. »

\all{Am Abend treffen sich die Besucher vor dem Feuerwerk.} « Le soir, les visiteurs se retrouvent devant le feu d'artifice. »
\end{exemplebox}

\courssub{Carte mentale du lexique}

\begin{popfigure}
\begin{center}
\begin{tikzpicture}[mindmap, grow cyclic,
  every node/.style={concept},
  root concept/.append style={concept color=popblue, font=\bfseries},
  level 1/.append style={level distance=3.6cm, sibling angle=72},
  level 2/.append style={level distance=2.2cm, sibling angle=45, font=\small}]
\node[root concept]{Traditionelle Feste}
  child[concept color=poporange]{ node {Fête / Tradition}
    child { node {das Fest} }
    child { node {die Feier} }
    child { node {der Brauch} } }
  child[concept color=popgreen]{ node {Danse / Musique}
    child { node {der Tanz} }
    child { node {das Lied} }
    child { node {tanzen} } }
  child[concept color=poppurple]{ node {Nourriture}
    child { node {das Essen} }
    child { node {die Brezel} }
    child { node {das Getränk} } }
  child[concept color=poppink]{ node {Vêtements}
    child { node {die Tracht} }
    child { node {das Dirndl} }
    child { node {tragen} } }
  child[concept color=popgold]{ node {Activités}
    child { node {der Umzug} }
    child { node {das Feuerwerk} }
    child { node {sich treffen} } };
\end{tikzpicture}
\end{center}
\legende{Carte mentale du lexique des fêtes traditionnelles : cinq champs à mémoriser.}
\end{popfigure}

\courssub{Familles de mots : un verbe, plusieurs noms}

L'allemand construit facilement des noms à partir de verbes. Connaître ces familles multiplie votre vocabulaire sans effort :

\begin{center}
\begin{tabularx}{\linewidth}{|X|X|X|}\hline
\rowcolor{popblueL} \textbf{Verbe} & \textbf{Noms dérivés} & \textbf{Français} \\ \hline
feiern & die Feier, das Fest, der Feiertag & fêter → la fête, le jour férié \\ \hline
tanzen & der Tanz, der Tänzer, die Tänzerin & danser → la danse, le danseur, la danseuse \\ \hline
trinken & das Getränk & boire → la boisson \\ \hline
singen & das Lied, der Sänger, die Sängerin & chanter → la chanson, le chanteur, la chanteuse \\ \hline
besuchen & der Besucher, der Besuch & visiter → le visiteur, la visite \\ \hline
\end{tabularx}
\end{center}

\begin{exemplebox}[Expressions utiles à réutiliser]
\all{ein Fest feiern} « fêter une fête » — \all{Wir feiern ein Fest im Dorf.}

\all{an einem Fest teilnehmen} « participer à une fête » — \all{Viele Touristen nehmen am Fest teil.} (attention : \all{teilnehmen} est séparable et se construit avec le datif : \all{an einem Fest})

\all{sich mit der Familie treffen} « se retrouver avec la famille » — \all{Ich treffe mich mit meiner Familie.} (\all{mit} + datif)

\all{traditionelle Kleidung tragen} « porter des vêtements traditionnels » — \all{Die Gäste tragen traditionelle Kleidung.}
\end{exemplebox}

\courssub{Texte d'application : un festival au Cameroun}

\begin{textebox}[Das Ngondo-Fest in Douala]
Jedes Jahr im Dezember findet in Douala das Ngondo-Fest statt. Es ist eine alte Tradition des Sawa-Volkes und dauert mehrere Tage. Tausende Besucher kommen ans Ufer des Wouri-Flusses, um an diesem Fest teilzunehmen.

Am Morgen beginnt der Umzug: Die Musikkapellen spielen, die Tänzer tanzen in ihren bunten Trachten, und die Kinder singen traditionelle Lieder. Viele Menschen tragen traditionelle Kleidung; einige verkleiden sich sogar als alte Könige. Die Atmosphäre ist fröhlich und laut.

Am Mittag treffen sich die Familien am Fluss. Es gibt viel zu essen: gegrillten Fisch, Hähnchen, Reis und süße Süßigkeiten für die Kinder. Die Erwachsenen trinken Palmwein oder ein kaltes Getränk. Das Essen schmeckt allen sehr gut.

Der Höhepunkt des Festes ist die Zeremonie auf dem Wasser: Ein Taucher springt in den Fluss und bringt eine Botschaft der Ahnen zurück. Alle Besucher schauen gespannt zu. Am Abend gibt es manchmal ein kleines Feuerwerk, und die Jugendlichen tanzen bis spät in die Nacht.

Das Ngondo-Fest bringt die ganze Gemeinschaft zusammen: Jung und Alt, Arm und Reich. Es zeigt, dass Traditionen lebendig sind — in Kamerun genauso wie in Deutschland, in Österreich oder in der Schweiz.
\end{textebox}

\textbf{Glossaire :} das Ufer (-) : la rive ; der Fluss, die Flüsse : le fleuve ; der Tänzer (-) : le danseur ; bunt : coloré ; sich verkleiden als : se déguiser en ; gegrillt : grillé ; der Palmwein : le vin de palme ; der Höhepunkt, die Höhepunkte : le point culminant ; der Taucher (-) : le plongeur ; die Botschaft, die Botschaften : le message ; der Ahne, die Ahnen : l'ancêtre ; gespannt : avec impatience ; die Gemeinschaft, die Gemeinschaften : la communauté ; lebendig : vivant.

\textbf{Questions de vocabulaire :}
\begin{enumerate}
\item Trouvez dans le texte trois noms de vêtements ou d'accessoires de fête.
\item Quels verbes du texte appartiennent au champ de la musique ?
\item Relevez deux expressions avec un verbe séparable et recopiez la phrase complète.
\end{enumerate}

\begin{exoresolu}[Compléter avec le bon mot]
Complétez avec : \all{Tracht, Umzug, Feuerwerk, statt, dauert, Getränk}.

1. Das Fest \dots\ drei Tage. \quad 2. Der \dots\ beginnt am Morgen im Dorf. \quad 3. Die Frauen tragen eine bunte \dots. \quad 4. Am Abend gibt es ein großes \dots. \quad 5. Das Oktoberfest findet in München \dots. \quad 6. Ein Bier ist ein kaltes \dots.
\tcblower
\textbf{Solution :} 1. \all{dauert} (le sujet \all{das Fest} est à la 3\textsuperscript{e} personne du singulier : durer). 2. \all{Umzug} (masculin, \all{der} : le défilé du matin). 3. \all{Tracht} (féminin, \all{eine} : le costume traditionnel porté). 4. \all{Feuerwerk} (neutre, \all{ein} : le feu d'artifice du soir). 5. \all{statt} (verbe séparable \all{stattfinden} : la particule va en fin de phrase). 6. \all{Getränk} (famille de \all{trinken} : la boisson).
\end{exoresolu}

\begin{exercice}[À vous de jouer]
\begin{enumerate}
\item Donnez l'article et le pluriel de : \all{Fest, Tanz, Lied, Brauch, Feiertag, Kostüm, Besucher, Umzug}.
\item Traduisez en français : \all{Die Musikkapelle spielt, und die Besucher tanzen in ihren Trachten.}
\item Traduisez en allemand : « La fête a lieu en décembre et dure trois jours. »
\item Complétez avec la bonne forme du verbe : \all{Die Kinder (werfen) Süßigkeiten. — Wir (sich treffen) am Fluss. — Er (tragen) ein Dirndl.}
\end{enumerate}
\end{exercice}

\begin{corrige}[Exercice : À vous de jouer]
\begin{enumerate}
\item \all{das Fest, die Feste} ; \all{der Tanz, die Tänze} ; \all{das Lied, die Lieder} ; \all{der Brauch, die Bräuche} ; \all{der Feiertag, die Feiertage} ; \all{das Kostüm, die Kostüme} ; \all{der Besucher, die Besucher} ; \all{der Umzug, die Umzüge}.
\item « La fanfare joue, et les visiteurs dansent dans leurs costumes traditionnels. »
\item \all{Das Fest findet im Dezember statt und dauert drei Tage.} (attention : particule \all{statt} en fin de première proposition ; \all{im Dezember} = in + dem).
\item \all{Die Kinder werfen Süßigkeiten. — Wir treffen uns am Fluss. — Er trägt ein Dirndl.} (\all{tragen} : voyelle \all{a → ä} aux 2\textsuperscript{e} et 3\textsuperscript{e} personnes du singulier ; \all{sich treffen} : pronom réfléchi \all{uns} après le verbe conjugué.)
\end{enumerate}
\end{corrige}

\begin{savaistu}[L'Oktoberfest, de Munich à Yaoundé]
Le mot \all{Oktoberfest} est connu dans le monde entier : c'est la plus grande fête populaire d'Allemagne, célébrée à Munich depuis 1810. Le saviez-vous ? Il existe même des « Oktoberfests » au Cameroun, organisés par la communauté germanophone de Yaoundé et de Douala, avec de la musique allemande, des bretzels et parfois des Trachten. Les fêtes traditionnelles voyagent, elles aussi !
\end{savaistu}

\begin{aretenir}[L'essentiel du lexique]
\begin{itemize}
\item Cinq champs : la fête (\all{das Fest, die Feier, die Tradition}), la musique (\all{der Tanz, das Lied, tanzen}), la nourriture (\all{das Essen, die Brezel, trinken}), les vêtements (\all{die Tracht, das Dirndl, tragen}), les activités (\all{der Umzug, das Feuerwerk, sich treffen}).
\item Toujours mémoriser article + pluriel : \all{das Fest, die Feste}.
\item Familles de mots : \all{feiern → die Feier} ; \all{tanzen → der Tanz} ; \all{trinken → das Getränk}.
\item Pièges : \all{das Fest} $\neq$ \all{fest} ; \all{die Kleidung} est singulier ; \all{teilnehmen}, \all{stattfinden} et \all{zusammenbringen} sont séparables ; \all{tragen} et \all{essen} changent de voyelle au présent (\all{er trägt}, \all{er isst}).
\end{itemize}
\end{aretenir}

\coursec{Grammaire 1 : le datif et les prépositions de lieu}

Dans la section précédente, nous avons enrichi notre vocabulaire des fêtes traditionnelles : \all{das Fest}, \all{die Tradition}, \all{der Umzug}, \all{das Bierzelt}, \all{feiern}. Mais pour employer ces mots dans des phrases correctes, il faut savoir les relier entre eux. En allemand, ce sont les \cle{cas} qui jouent ce rôle. Après l'accusatif (le complément d'objet direct), voici le deuxième grand cas du programme de Première : \cle{le datif}, le cas du complément indirect, et son compagnon inséparable : les \cle{prépositions de lieu}.

\courssub{Observation : des exemples tirés du texte}

Relisons quelques phrases du texte support sur la fête de la bière et soulignons les groupes prépositionnels :

\begin{exemplebox}[Phrases du texte à observer]
\all{Das Oktoberfest findet in München statt.} — La fête de la bière a lieu à Munich.\par
\all{Die Wiesn liegt auf der Theresienwiese.} — La Wiesn se trouve sur la prairie de Thérèse.\par
\all{Die Besucher sitzen in großen Bierzelten.} — Les visiteurs sont assis dans de grandes tentes à bière.\par
\all{Die Menschen tanzen in den Straßen.} — Les gens dansent dans les rues.\par
\all{Beim Umzug werfen die Leute Süßigkeiten.} — Lors du défilé, les gens jettent des bonbons.\par
\all{Viele Familien kommen mit der Familie und mit Freunden.} — Beaucoup viennent avec la famille et avec des amis.\par
\all{Die Touristen fahren nach München.} — Les touristes vont à Munich.
\end{exemplebox}

Observons les articles : \all{der Theresienwiese}, \all{den Straßen}, \all{beim} (= \all{bei dem}), \all{mit der Familie}. Aucun de ces groupes n'est un complément d'objet direct : ce sont des compléments de \emph{lieu}, de \emph{temps} ou de \emph{moyen}, introduits par une préposition. En allemand, la préposition \emph{commande} le cas de l'article qui la suit. Ici, c'est le \cle{datif}.

\courssub{Les formes du datif}

\begin{definition}[Le datif]
Le \cle{datif} (\all{der Dativ}) est le cas du complément d'objet indirect (COI) et des compléments introduits par certaines prépositions. Il répond aux questions \all{wem ?} (à qui ?), \all{wo ?} (où ?), \all{wann ?} (quand ?).
\end{definition}

\begin{center}
\begin{tabular}{|l|l|l|l|}
\hline
\rowcolor{popblueL} & Masculin & Féminin & Neutre \\
\hline
Nominatif & der Vater & die Mutter & das Kind \\
\hline
Accusatif & den Vater & die Mutter & das Kind \\
\hline
\rowcolor{popgoldL} \textbf{Datif} & \textbf{dem} Vater & \textbf{der} Mutter & \textbf{dem} Kind \\
\hline
\end{tabular}
\end{center}

\begin{propriete}[Les articles au datif]
\begin{itemize}
\item Article défini : masculin \all{dem}, féminin \all{der}, neutre \all{dem}, pluriel \all{den}.
\item Article indéfini : masculin \all{einem}, féminin \all{einer}, neutre \all{einem}, pluriel : pas d'article.
\item \textbf{Règle d'or du pluriel :} au datif pluriel, le nom prend un \textbf{-n} final (s'il n'en a pas déjà un) : \all{die Straße} → \all{in den Straß\emph{en}} ; \all{das Fest} → \all{auf den Fest\emph{en}} ; \all{das Kind} → \all{mit den Kind\emph{ern}} (déjà en -n).
\item L'adjectif épithète prend \all{-en} au datif (toutes déclinaisons confondues) : \all{in groß\emph{en} Bierzelten}, \all{mit mein\emph{em} best\emph{en} Freund}, \all{mit gut\emph{en} Freunden}.
\end{itemize}
\end{propriete}

\begin{popfigure}
\begin{center}
\begin{tikzpicture}
\node[draw=popblue, fill=popblueL, rounded corners, font=\bfseries, minimum width=3.2cm, minimum height=0.9cm] (titre) at (0,0) {DER DATIV};
\node[draw=popgreen, fill=popgreenL, rounded corners, align=center, minimum width=2.6cm] (m) at (-4.5,-1.6) {Masculin\\ dem / einem};
\node[draw=poppink, fill=poppinkL, rounded corners, align=center, minimum width=2.6cm] (f) at (-1.5,-1.6) {Féminin\\ der / einer};
\node[draw=poporange, fill=poporangeL, rounded corners, align=center, minimum width=2.6cm] (n) at (1.5,-1.6) {Neutre\\ dem / einem};
\node[draw=poppurple, fill=poppurpleL, rounded corners, align=center, minimum width=2.6cm] (p) at (4.5,-1.6) {Pluriel\\ den + NOM-n};
\draw[-{Stealth}, thick, popblue] (titre) -- (m);
\draw[-{Stealth}, thick, popblue] (titre) -- (f);
\draw[-{Stealth}, thick, popblue] (titre) -- (n);
\draw[-{Stealth}, thick, popblue] (titre) -- (p);
\node[align=center, font=\itshape] at (0,-2.8) {mit dem Bruder — mit der Schwester — mit dem Kind — mit den Kindern};
\end{tikzpicture}
\end{center}
\legende{Tableau mental des articles au datif : dem / der / dem / den + n}
\end{popfigure}

\courssub{Les prépositions toujours suivies du datif}

\begin{propriete}[aus, bei, mit, nach, seit, von, zu, gegenüber + datif]
Huit prépositions très fréquentes sont \textbf{toujours} suivies du datif, sans exception :\par
\begin{center}
\begin{tabular}{|l|l|l|}
\hline
\rowcolor{popblueL} Préposition & Sens & Exemple traduit \\
\hline
\all{aus} & de (origine, matière) & \all{Sie kommt aus Kamerun.} — Elle vient du Cameroun. \\
\hline
\all{bei} & chez, lors de, près de & \all{Beim Umzug singen alle.} — Lors du défilé, tous chantent. \\
\hline
\all{mit} & avec (compagnie, moyen) & \all{Er feiert mit der Familie.} — Il fête avec la famille. \\
\hline
\all{nach} & à, vers (ville, pays) ; après (temps) & \all{Wir fahren nach München.} — Nous allons à Munich. \\
\hline
\all{seit} & depuis (temps) & \all{Seit einem Jahr lernt sie Deutsch.} — Elle apprend l'allemand depuis un an. \\
\hline
\all{von} & de (provenance, auteur, possession) & \all{Das Geschenk ist von meinem Freund.} — Le cadeau est de mon ami. \\
\hline
\all{zu} & à, chez (personne, but) & \all{Ich gehe zu meiner Tante.} — Je vais chez ma tante. \\
\hline
\all{gegenüber} & en face de (lieu) & \all{Das Zelt ist gegenüber dem Markt.} — La tente est en face du marché. \\
\hline
\end{tabular}
\end{center}
\emph{Astuce de mémorisation :} la phrase mnémotechnique « \textbf{a}u \textbf{b}al, \textbf{m}arie \textbf{n}'a \textbf{s}u \textbf{v}aler \textbf{z}éro \textbf{g}avotte » (aus, bei, mit, nach, seit, von, zu, gegenüber).
\end{propriete}

\begin{attention}[nach et zu : deux façons de dire « à »]
Pour le mouvement, l'allemand distingue : \all{nach} + nom de ville ou de pays \textbf{sans article} (\all{nach München}, \all{nach Kamerun}, \all{nach Hause}) ; \all{zu} + personne, lieu précis ou activité \textbf{avec article} (\all{zu meiner Tante}, \all{zum Markt}, \all{zum Essen}). On dit donc \all{Ich fahre nach Douala} mais \all{Ich gehe zum Bäcker} (je vais chez le boulanger).
\end{attention}

\courssub{Les prépositions de lieu à deux cas (Wechselpräpositionen)}

Neuf prépositions de lieu peuvent être suivies du datif \emph{ou} de l'accusatif : \all{in, an, auf, über, unter, vor, hinter, neben, zwischen}. Comment choisir ? Tout dépend de la question que l'on pose au verbe.

\begin{propriete}[La règle d'or : wo ? ou wohin ?]
\begin{itemize}
\item \textbf{WO ?} (où ? = position, état, pas de déplacement) → \textbf{DATIF}\par
\all{Das Fest findet auf der Wiese statt.} — La fête a lieu sur la prairie. (où ? → \all{der Wiese}, datif)
\item \textbf{WOHIN ?} (vers où ? = mouvement, direction, déplacement) → \textbf{ACCUSATIF}\par
\all{Die Besucher kommen auf die Wiese.} — Les visiteurs viennent sur la prairie. (vers où ? → \all{die Wiese}, accusatif)
\end{itemize}
\end{propriete}

\begin{popfigure}
\begin{center}
\begin{tikzpicture}
\node[draw=popdark, fill=popcream, rounded corners, font=\bfseries, align=center, minimum width=4.5cm, minimum height=1cm] (q) at (0,0) {Préposition de lieu\\ in, an, auf, vor, hinter...};
\node[draw=popgreen, fill=popgreenL, rounded corners, align=center, minimum width=4.2cm, minimum height=1.4cm] (wo) at (-4,-2.4) {\textbf{WO ?} (position)\\ $\rightarrow$ DATIF\\ in der Stadt};
\node[draw=poporange, fill=poporangeL, rounded corners, align=center, minimum width=4.2cm, minimum height=1.4cm] (wohin) at (4,-2.4) {\textbf{WOHIN ?} (mouvement)\\ $\rightarrow$ AKKUSATIV\\ in die Stadt};
\draw[-{Stealth}, very thick, popgreen] (q.south west) ++(0.3,0) -- (wo.north) node[midway, left, font=\itshape] {où ?};
\draw[-{Stealth}, very thick, poporange] (q.south east) ++(-0.3,0) -- (wohin.north) node[midway, right, font=\itshape] {vers où ?};
\node[align=center, font=\itshape] at (0,-4) {Ich bin in der Stadt. (je suis en ville) — Ich gehe in die Stadt. (je vais en ville)};
\end{tikzpicture}
\end{center}
\legende{La règle d'or : wo ? → datif (position) ; wohin ? → accusatif (mouvement)}
\end{popfigure}

\begin{exemplebox}[Exemples traduits : le même lieu, deux cas]
\all{Sie sitzen in großen Bierzelten.} — Ils sont assis dans de grandes tentes à bière. (wo ? datif)\par
\all{Sie gehen in große Bierzelte.} — Ils entrent dans de grandes tentes à bière. (wohin ? accusatif)\par
\all{Die Menschen tanzen in den Straßen.} — Les gens dansent dans les rues. (wo ? datif pluriel : \all{den Straßen})\par
\all{Die Musiker ziehen durch die Straßen.} — Les musiciens défilent à travers les rues. (mouvement, mais \all{durch} exige toujours l'accusatif)\par
\all{Das Bier steht auf dem Tisch.} — La bière est sur la table. (wo ? datif)\par
\all{Der Kellner stellt das Bier auf den Tisch.} — Le serveur pose la bière sur la table. (wohin ? accusatif)
\end{exemplebox}

\courssub{Les contractions fréquentes}

En allemand parlé comme écrit, certaines prépositions fusionnent avec l'article défini. Il faut les reconnaître et les employer systématiquement :

\begin{center}
\begin{tabularx}{\linewidth}{|X|X|X|}
\hline
\rowcolor{popblueL} Contraction & Forme complète & Exemple traduit \\
\hline
\all{beim} & bei dem & \all{Beim Umzug tanzen alle.} — Lors du défilé, tous dansent. \\
\hline
\all{vom} & von dem & \all{Er kommt vom Fest.} — Il revient de la fête. \\
\hline
\all{zum} & zu dem & \all{Wir gehen zum Markt.} — Nous allons au marché. \\
\hline
\all{zur} & zu der & \all{Sie geht zur Schule.} — Elle va à l'école. \\
\hline
\all{im} & in dem & \all{Im Zelt ist es warm.} — Dans la tente, il fait chaud. (position) \\
\hline
\all{am} & an dem & \all{Am Abend gibt es Feuerwerk.} — Le soir, il y a un feu d'artifice. \\
\hline
\all{ins} & in das & \all{Wir gehen ins Zelt.} — Nous entrons dans la tente. (mouvement !) \\
\hline
\end{tabularx}
\end{center}

\begin{attention}[Les pièges du francophone]
\begin{itemize}
\item \textbf{Le -n du datif pluriel :} on écrit \all{in den Straß\emph{en}}, \all{mit den Kind\emph{ern}}, jamais « in den Straße ». C'est l'erreur la plus fréquente en Première !
\item \textbf{zu Hause / nach Hause :} \all{Ich bin zu Hause.} (je suis à la maison, position, datif figé) mais \all{Ich gehe nach Hause.} (je rentre à la maison, mouvement, accusatif figé). Deux expressions figées à apprendre par cœur.
\item \textbf{Pas de cas en français :} le français dit « dans les rues » sans rien changer ; l'allemand doit marquer la fonction par l'article (\all{in den Straßen}). Ne traduis jamais mot à mot : demande-toi toujours « quelle préposition ? quel cas ? ».
\item \textbf{Faux ami :} \all{bei} ne signifie pas « par » mais « chez / lors de / près de » : \all{bei meinen Eltern} = chez mes parents ; \all{beim Essen} = pendant le repas.
\item \textbf{Prépositions \all{in der Nähe von} / \all{weit von} :} elles régissent le datif (via \all{von}) : \all{in der Nähe vom Markt} (= \all{von dem Markt}).
\end{itemize}
\end{attention}

\begin{center}
\begin{tabularx}{\linewidth}{|X|X|X|}
\hline
\rowcolor{popblueL} Français & Deutsch & Remarque \\
\hline
à Munich (ville) & nach München & \all{nach} + ville/pays sans article \\
\hline
chez mon ami & bei meinem Freund & \all{bei} + datif \\
\hline
avec la famille & mit der Familie & \all{mit} + datif \\
\hline
dans les rues (position) & in den Straßen & wo ? → datif + -n pluriel \\
\hline
dans la tente (mouvement) & ins Zelt / in das Zelt & wohin ? → accusatif \\
\hline
à la maison (position / mouvement) & zu Hause / nach Hause & expressions figées \\
\hline
en face du marché & gegenüber dem Markt & \all{gegenüber} + datif \\
\hline
\end{tabularx}
\end{center}

\courssub{Texte d'application : un dimanche de fête à Yaoundé}

\begin{textebox}[Ein Festtag in Jaunde]
\all{Am Sonntag feiert die Familie Mbarga ein traditionelles Fest. Schon am Morgen fahren viele Gäste mit dem Auto oder mit dem Taxi zum Haus in der Nähe vom Markt. Die Frauen kochen in der Küche, die Männer sitzen unter einem großen Baum und sprechen über Fußball. Die Kinder spielen auf der Straße vor dem Haus. Um zwölf Uhr steht das Essen auf dem Tisch: Ndolé, Poulet DG und Reis. Alle sitzen an einem langen Tisch im Garten. Nach dem Essen gibt es einen Umzug durch die Straßen des Viertels. Die Musiker spielen auf ihren Trommeln, und die Leute tanzen in den Straßen. Beim Umzug tragen die Frauen ihre schönsten Kleider aus buntem Stoff. Am Abend sitzt die ganze Familie noch zusammen und erzählt von alten Traditionen. Die Großmutter erzählt von den Festen in ihrem Dorf. „Früher haben wir drei Tage gefeiert“, sagt sie. Die Kinder hören zu und lernen viel von ihrer Großmutter. Gegen Mitternacht gehen die Gäste nach Hause. Alle sind müde, aber glücklich.}\par
\vspace{0.3em}
\emph{Glossaire :} \all{der Gast, die Gäste} : l'invité ; \all{die Küche, -n} : la cuisine ; \all{der Baum, die Bäume} : l'arbre ; \all{der Umzug, -e} : le défilé ; \all{die Trommel, -n} : le tambour ; \all{das Kleid, die Kleider} : la robe, le vêtement ; \all{bunt} : coloré ; \all{der Stoff, -e} : le tissu ; \all{das Viertel, -} : le quartier ; \all{erzählen} : raconter ; \all{zuhören} (trennbar) : écouter (attentivement).
\end{textebox}

Questions de compréhension : 1) \all{Wo kochen die Frauen?} 2) \all{Wo spielen die Kinder?} 3) \all{Wohin gehen die Gäste gegen Mitternacht?} 4) \all{Wovon erzählt die Großmutter?}

\courssub{Méthode et entraînement}

\begin{methode}[Comment choisir le bon cas après une préposition de lieu]
\begin{enumerate}
\item \textbf{Repère la préposition} : est-ce une préposition \emph{toujours} datif (\all{aus, bei, mit, nach, seit, von, zu, gegenüber}) ? → datif automatique.
\item Si c'est une \emph{Wechselpräposition} (\all{in, an, auf, vor, hinter, über, unter, neben, zwischen}), \textbf{pose la question au verbe} : « wo ? » (où est-ce que ça se passe ?) → datif ; « wohin ? » (vers où est-ce qu'on va ?) → accusatif.
\item \textbf{Choisis l'article} selon le genre du nom : datif = \all{dem / der / dem / den} ; accusatif = \all{den / die / das / die}.
\item \textbf{Vérifie le pluriel} : au datif, ajoute \textbf{-n} au nom (\all{in den Straßen}) s'il n'en a pas déjà un.
\item \textbf{Pense aux contractions} : \all{bei dem → beim}, \all{in dem → im}, \all{zu dem → zum}, \all{von dem → vom}, \all{in das → ins}.
\end{enumerate}
\end{methode}

\begin{exoresolu}[Complète avec l'article au bon cas]
Énoncé. Complète : 1) Die Familie feiert mit \_\_\_ (die Großeltern). 2) Das Fest findet in \_\_\_ (das Dorf) statt. 3) Die Gäste gehen in \_\_\_ (das Haus). 4) Die Kinder tanzen auf \_\_\_ (die Straße). 5) Nach \_\_\_ (das Essen) tanzen alle.
\tcblower
Solution détaillée.\par
1) \all{mit den Großeltern} — \all{mit} est toujours suivi du datif ; pluriel : \all{den} (le nom \all{Großeltern} a déjà son -n). « La famille fête avec les grands-parents. »\par
2) \all{in dem Dorf} (ou \all{im Dorf}) — \all{findet statt} = « a lieu » : position, question « wo ? » → datif ; neutre : \all{dem}. « La fête a lieu dans le village. »\par
3) \all{in das Haus} (ou \all{ins Haus}) — \all{gehen} exprime un mouvement : « wohin ? » → accusatif ; neutre : \all{das}. « Les invités entrent dans la maison. »\par
4) \all{auf der Straße} — les enfants dansent, ils sont déjà sur la rue : position « wo ? » → datif ; féminin : \all{der}. « Les enfants dansent dans la rue. »\par
5) \all{Nach dem Essen} — \all{nach} (temporel) est toujours suivi du datif ; neutre : \all{dem}. « Après le repas, tout le monde danse. »
\end{exoresolu}

\begin{exercice}[À ton tour]
1) Mets à la bonne forme (article + nom, ou contraction) : a) Wir fahren mit \_\_\_ (der Bus) zum Fest. b) Die Musiker spielen in \_\_\_ (die Straßen, pluriel). c) Sie kommt von \_\_\_ (ihre Tante). d) Das Zelt steht neben \_\_\_ (der Markt). e) Die Touristen fahren nach \_\_\_ (Deutschland).\par
2) Traduis en allemand : a) Les gens dansent dans les rues. b) Nous allons chez nos amis. c) Lors du défilé, les enfants reçoivent des bonbons. d) Après la fête, nous rentrons à la maison.
\end{exercice}

\begin{corrige}[Exercice « À ton tour »]
1) a) \all{mit dem Bus} (\all{mit} + datif masculin). b) \all{in den Straßen} (position, datif pluriel + -n). c) \all{von ihrer Tante} (\all{von} + datif féminin : possessif \all{ihrer}). d) \all{neben dem Markt} (position \all{wo?}, datif masculin). e) \all{nach Deutschland} (pays sans article).\par
2) a) \all{Die Leute tanzen in den Straßen.} b) \all{Wir gehen zu unseren Freunden.} (\all{zu} + datif pluriel \all{unseren}). c) \all{Beim Umzug bekommen die Kinder Süßigkeiten.} (\all{beim} = \all{bei dem}). d) \all{Nach dem Fest gehen wir nach Hause.}
\end{corrige}

\begin{aretenir}[L'essentiel du datif et des prépositions de lieu]
\begin{itemize}
\item Datif (articles définis) : \all{dem / der / dem / den} + \textbf{-n} au pluriel du nom ; indéfini : \all{einem / einer / einem}.
\item Toujours datif après : \all{aus, bei, mit, nach, seit, von, zu, gegenüber}.
\item Wechselpräpositionen (\all{in, an, auf, vor, hinter, über, unter, neben, zwischen}) : \textbf{wo ? → datif ; wohin ? → accusatif}.
\item Contractions obligatoires : \all{beim, vom, zum, zur, im, am, ins}.
\item Expressions figées : \all{zu Hause} (position) / \all{nach Hause} (mouvement).
\end{itemize}
\end{aretenir}

\coursec{Grammaire 2 : le passé composé (das Perfekt)}

Dans la section précédente, nous avons appris à dire \emph{où} les choses se passent grâce au datif et aux prépositions de lieu : \all{auf dem Festplatz}, \all{in der Stadt}, \all{bei der Familie}. Mais notre texte sur les fêtes traditionnelles racontait surtout des événements \emph{passés} : la fête a eu lieu, les gens ont dansé, les visiteurs sont venus. Pour raconter tout cela, l'allemand utilise le temps que nous allons étudier maintenant : le \cle{Perfekt}, le passé composé allemand. C'est le temps roi de la conversation et du récit personnel — exactement celui qu'il vous faut pour raconter votre dernière fête de village ou vos vacances.

\courssub{Observation : découvrir le Perfekt dans un texte}

\begin{textebox}[Mein Wochenende auf dem Fest]
Letztes Wochenende war ein großes Fest in unserem Dorf. Ich habe das Fest mit meiner Familie gefeiert. Am Morgen haben wir unsere Trachten getragen und wir sind zum Umzug gegangen. Viele Besucher sind aus der ganzen Region gekommen. Einige sind mit dem Bus gefahren, andere sind zu Fuß gekommen. Auf dem Festplatz haben wir gesungen und getanzt. Meine Mutter hat ein traditionelles Essen gekocht, und wir haben zusammen gegessen und getrunken. Das Fest hat zwei Tage gedauert. Am Abend sind wir müde nach Hause gegangen, aber wir waren sehr glücklich. Nächstes Jahr findet das Fest wieder statt — ich freue mich schon!
\end{textebox}

\textbf{Glossaire :} \all{letztes Wochenende} le week-end dernier ; \all{die Tracht, -en} le costume traditionnel ; \all{der Umzug, „-e} le défilé ; \all{der Besucher, -} le visiteur ; \all{zu Fuß} à pied ; \all{der Festplatz, „-e} la place de fête ; \all{müde} fatigué ; \all{stattfinden} (hat stattgefunden) avoir lieu ; \all{sich freuen} se réjouir.

\textbf{Questions de compréhension :}
\begin{enumerate}
\item \all{Wann war das Fest?} (Quand la fête a-t-elle eu lieu ?)
\item \all{Was hat die Mutter gekocht?} (Qu'est-ce que la mère a cuisiné ?)
\item \all{Wie lange hat das Fest gedauert?} (Combien de temps la fête a-t-elle duré ?)
\end{enumerate}

Observons maintenant les verbes de ce texte. Tous les événements passés sont exprimés avec \textbf{deux mots} :

\begin{center}
\begin{tabularx}{\linewidth}{|X|X|X|}
\hline
\rowcolor{popblueL} Phrase allemande & Auxiliaire & Participe passé \\
\hline
\all{Ich habe das Fest gefeiert.} & habe & gefeiert \\
\hline
\all{Wir sind zum Umzug gegangen.} & sind & gegangen \\
\hline
\all{Die Besucher sind gekommen.} & sind & gekommen \\
\hline
\all{Meine Mutter hat gekocht.} & hat & gekocht \\
\hline
\all{Das Fest hat zwei Tage gedauert.} & hat & gedauert \\
\hline
\end{tabularx}
\end{center}

Deux constats essentiels : d'une part, l'auxiliaire est tantôt \all{haben} (avoir), tantôt \all{sein} (être) ; d'autre part — et c'est la grande différence avec le français — le participe passé ne se place pas juste après l'auxiliaire, mais \textbf{tout à la fin de la phrase}.

\courssub{Formation du Perfekt : auxiliaire + participe passé}

\begin{definition}[Das Perfekt]
Le \cle{Perfekt} est un temps composé formé de deux éléments : l'auxiliaire \all{haben} ou \all{sein}, conjugué au présent et placé en \textbf{2\textsuperscript{e} position}, et le \cle{participe passé} (Partizip II) du verbe principal, placé \textbf{à la toute dernière position} de la phrase.
\end{definition}

\begin{propriete}[La place du participe passé : toujours à la fin]
Le participe passé se place \textbf{TOUJOURS} à la fin de la phrase, quelle que soit sa longueur. Tout le reste (compléments de temps, de lieu, d'objet) se loge entre l'auxiliaire et le participe : c'est la fameuse « construction en cadre » (\all{Satzklammer}).

\all{Wir haben das Fest mit der Familie gefeiert.} « Nous avons fêté la fête avec la famille. »

\all{Ich habe am Wochenende mit meinen Freunden auf dem Festplatz getanzt.} « J'ai dansé le week-end avec mes amis sur la place de fête. »
\end{propriete}

\begin{popfigure}
\begin{tikzpicture}[
box/.style={draw=popblue, thick, rounded corners=4pt, fill=popblueL, align=center, inner sep=5pt, font=\small},
pbox/.style={draw=poporange, thick, rounded corners=4pt, fill=poporangeL, align=center, inner sep=5pt, font=\small},
fleche/.style={-{Stealth[length=2.5mm]}, thick, popdark}]
\node[box, align=center] (s) at (0,0){Sujet\\ \all{Wir}};
\node[box, align=center] (aux) at (3.2,0){Auxiliaire conjugué\\ (2\textsuperscript{e} position)\\ \all{haben / sind}};
\node[box, fill=popgreenL, draw=popgreen, align=center] (mil) at (7.2,0){Compléments\\ (temps, lieu, objet)\\ \all{das Fest mit der Familie}};
\node[pbox, align=center] (part) at (11.6,0){Participe passé\\ (dernière position)\\ \all{gefeiert}};
\draw[fleche] (s) -- (aux);
\draw[fleche] (aux) -- (mil);
\draw[fleche] (mil) -- (part);
\draw[fleche, poporange, dashed, bend right=35] (aux.south) to node[below, font=\footnotesize, text=poporange]{le cadre se referme à la fin} (part.south);
\end{tikzpicture}
\legende{La construction en cadre : l'auxiliaire ouvre la phrase, le participe passé la ferme.}
\end{popfigure}

\courssub{Le participe passé des verbes faibles : ge- + radical + -t}

Les \cle{verbes faibles} (verbes réguliers) forment leur participe passé de manière mécanique : \textbf{ge- + radical + -t}. Le radical s'obtient en retirant la terminaison \all{-en} de l'infinitif.

\begin{center}
\begin{tabularx}{\linewidth}{|X|X|X|}
\hline
\rowcolor{popblueL} Infinitif & Participe passé & Exemple traduit \\
\hline
\all{feiern} (fêter) & \all{gefeiert} & \all{Ich habe das Fest gefeiert.} J'ai fêté la fête. \\
\hline
\all{tanzen} (danser) & \all{getanzt} & \all{Wir haben getanzt.} Nous avons dansé. \\
\hline
\all{machen} (faire) & \all{gemacht} & \all{Was hast du gemacht?} Qu'as-tu fait ? \\
\hline
\all{kochen} (cuisiner) & \all{gekocht} & \all{Mama hat gekocht.} Maman a cuisiné. \\
\hline
\all{fragen} (demander) & \all{gefragt} & \all{Er hat mich gefragt.} Il m'a demandé. \\
\hline
\all{arbeiten} (travailler) & \all{gearbeitet} & \all{Sie hat viel gearbeitet.} Elle a beaucoup travaillé. \\
\hline
\end{tabularx}
\end{center}

Remarque : les verbes dont le radical se termine en \all{-t} ou \all{-d} (comme \all{arbeiten}, \all{warten}) prennent \textbf{-et} au lieu de \textbf{-t} pour faciliter la prononciation : \all{gearbeitet}, \all{gewartet}.

\courssub{Le participe passé des verbes forts : ge- + (radical modifié) + -en}

Les \cle{verbes forts} (verbes irréguliers) forment leur participe avec \textbf{ge- + radical souvent modifié + -en}. Il faut les apprendre par cœur, comme les verbes irréguliers anglais. Voici les plus utiles pour raconter une fête :

\begin{center}
\begin{tabularx}{\linewidth}{|X|X|X|}
\hline
\rowcolor{popblueL} Infinitif & Participe passé & Exemple traduit \\
\hline
\all{essen} (manger) & \all{gegessen} & \all{Wir haben zusammen gegessen.} Nous avons mangé ensemble. \\
\hline
\all{trinken} (boire) & \all{getrunken} & \all{Sie haben Tee getrunken.} Ils ont bu du thé. \\
\hline
\all{singen} (chanter) & \all{gesungen} & \all{Wir haben gesungen.} Nous avons chanté. \\
\hline
\all{tragen} (porter) & \all{getragen} & \all{Ich habe eine Tracht getragen.} J'ai porté un costume traditionnel. \\
\hline
\all{sehen} (voir) & \all{gesehen} & \all{Hast du den Umzug gesehen?} As-tu vu le défilé ? \\
\hline
\all{gehen} (aller) & \all{gegangen} & \all{Sie sind zum Umzug gegangen.} Ils sont allés au défilé. \\
\hline
\all{kommen} (venir) & \all{gekommen} & \all{Die Gäste sind gekommen.} Les invités sont venus. \\
\hline
\all{fahren} (aller en véhicule) & \all{gefahren} & \all{Wir sind nach München gefahren.} Nous sommes allés à Munich. \\
\hline
\all{bleiben} (rester) & \all{geblieben} & \all{Er ist zu Hause geblieben.} Il est resté à la maison. \\
\hline
\end{tabularx}
\end{center}

\courssub{Les exceptions : les participes sans ge-}

\begin{propriete}[Pas de ge- pour les verbes en -ieren et les préfixes inséparables]
Deux catégories de verbes forment leur participe passé \textbf{sans le préfixe ge-} :
\begin{itemize}
\item les verbes en \textbf{-ieren} (souvent d'origine française) : \all{studieren} $\rightarrow$ \all{studiert}, \all{organisieren} $\rightarrow$ \all{organisiert}, \all{dekorieren} $\rightarrow$ \all{dekoriert} ;
\item les verbes à \textbf{préfixe inséparable} (\all{be-, ver-, er-, ent-, emp-, ge-, zer-}) : \all{besuchen} $\rightarrow$ \all{besucht}, \all{verkleiden} $\rightarrow$ \all{verkleidet}, \all{erzählen} $\rightarrow$ \all{erzählt}.
\end{itemize}
\all{Wir haben die Großeltern besucht.} « Nous avons rendu visite aux grands-parents. »

\all{Die Kinder haben sich verkleidet.} « Les enfants se sont déguisés. »
\end{propriete}

\begin{propriete}[Les verbes séparables : ge- au milieu]
Pour les verbes à \textbf{particule séparable}, le \all{ge-} se glisse \textbf{entre la particule et le radical} : particule + ge + radical + terminaison.

\all{stattfinden} $\rightarrow$ \all{stattgefunden} : \all{Das Fest hat im Dorf stattgefunden.} « La fête a eu lieu au village. »

\all{teilnehmen} $\rightarrow$ \all{teilgenommen} : \all{Viele Leute haben an der Feier teilgenommen.} « Beaucoup de gens ont participé à la fête. »

\all{mitmachen} $\rightarrow$ \all{mitgemacht} : \all{Alle Kinder haben mitgemacht.} « Tous les enfants ont participé. »
\end{propriete}

\courssub{Le choix de l'auxiliaire : haben ou sein ?}

C'est LA grande question du Perfekt. La règle est simple si on l'apprend avec des images :

\begin{propriete}[sein + participe pour les verbes de déplacement et de changement d'état]
On utilise l'auxiliaire \all{sein} avec :
\begin{itemize}
\item les verbes de \textbf{déplacement} (mouvement d'un point A vers un point B) : \all{gehen, fahren, kommen, reisen, fliegen, laufen, steigen} ;
\item les verbes de \textbf{changement d'état} : \all{sterben} (mourir), \all{aufstehen} (se lever), \all{einschlafen} (s'endormir), \all{werden} (devenir) ;
\item les deux verbes \all{sein} et \all{bleiben} eux-mêmes.
\end{itemize}
\all{Wir sind nach München gefahren.} « Nous sommes allés à Munich en voiture. »

\all{Die Besucher sind aus der ganzen Welt gekommen.} « Les visiteurs sont venus du monde entier. »

\all{Ich bin lange geblieben.} « Je suis resté longtemps. »

Pour \textbf{tous les autres verbes} (actions, verbes transitifs, verbes réfléchis), on utilise \all{haben} : \all{Ich habe gefeiert. Wir haben gesungen. Sie hat gekocht.}
\end{propriete}

Rappel de la conjugaison des deux auxiliaires au présent :

\begin{center}
\begin{tabular}{|l|l|l|}
\hline
\rowcolor{popblueL} Personne & \all{haben} & \all{sein} \\
\hline
ich & habe & bin \\
\hline
du & hast & bist \\
\hline
er/sie/es & hat & ist \\
\hline
wir & haben & sind \\
\hline
ihr & habt & seid \\
\hline
sie/Sie & haben & sind \\
\hline
\end{tabular}
\end{center}

Bonne nouvelle pour les francophones : la logique est proche de celle du français ! On dit « je suis allé » et l'allemand dit \all{ich bin gegangen} ; on dit « j'ai mangé » et l'allemand dit \all{ich habe gegessen}. La différence, c'est que l'allemand est plus strict : \emph{tous} les verbes de mouvement prennent \all{sein}, sans exception, alors que le français ne connaît qu'une petite liste (aller, venir, arriver, partir\dots).

\courssub{Emploi du Perfekt : le temps du récit oral}

\begin{aretenir}[Quand utilise-t-on le Perfekt ?]
Le Perfekt est le temps du passé utilisé \textbf{à l'oral et dans les récits personnels} : conversations, lettres à des amis, récits de vacances, de fêtes, de week-ends. C'est l'équivalent direct du passé composé français : \all{Ich habe das Fest gefeiert.} = « J'ai fêté la fête. » À la question, l'auxiliaire passe en première position et le participe reste à la fin : \all{Hast du das Fest gesehen?} « As-tu vu la fête ? »
\end{aretenir}

\begin{popfigure}
\begin{tikzpicture}[fleche/.style={-{Stealth[length=3mm]}, very thick, popdark}]
\draw[fleche] (0,0) -- (12.5,0);
\foreach \x/\nom/\couleur in {2/{Präsens\\ (présent)}/popgreen, 6.2/{Perfekt\\ (passé : oral,\\ récit personnel)}/poporange, 10.4/{Präteritum\\ (passé : textes\\ écrits formels)}/popblue}{
\node[circle, draw=\couleur, fill=\couleur!25, thick, inner sep=2.5pt] at (\x,0) {};
\node[align=center, font=\small, above=6pt] at (\x,0) {\\om};}
\node[align=center, font=\footnotesize, text=popmuted, below=8pt] at (6.2,0) {\all{Ich habe gefeiert.} / \all{Ich feierte.}};
\end{tikzpicture}
\legende{Frise des temps : le Perfekt est le passé de la langue parlée et du récit personnel.}
\end{popfigure}

\begin{center}
\begin{tabularx}{\linewidth}{|X|X|X|}
\hline
\rowcolor{popblueL} Français & Deutsch & Remarque \\
\hline
J'ai fêté la fête. & \all{Ich habe das Fest gefeiert.} & même logique : auxiliaire + participe \\
\hline
Nous sommes allés à Munich. & \all{Wir sind nach München gefahren.} & « être » = \all{sein} pour le mouvement \\
\hline
Elle a dansé avec son frère. & \all{Sie hat mit ihrem Bruder getanzt.} & le participe va à la fin, pas après l'auxiliaire \\
\hline
Ils ont visité le marché. & \all{Sie haben den Markt besucht.} & \all{besuchen} : préfixe inséparable, pas de ge- \\
\hline
\end{tabularx}
\end{center}

\begin{exemplebox}[Exemples traduits]
\all{Ich habe eine Tracht getragen.} « J'ai porté un costume traditionnel. »

\all{Wir haben gesungen und getanzt.} « Nous avons chanté et dansé. »

\all{Das Fest hat zwei Wochen gedauert.} « La fête a duré deux semaines. »

\all{Sie sind zum Umzug gegangen.} « Ils sont allés au défilé. »

\all{Hast du das Fest gesehen?} « As-tu vu la fête ? » (à la question, l'auxiliaire passe en première position, le participe reste à la fin.)
\end{exemplebox}

\begin{attention}[Erreurs fréquentes des francophones]
\begin{itemize}
\item \textbf{Ne colle jamais le participe à l'auxiliaire comme en français !} « Ich habe gefeiert das Fest » est FAUX. En allemand : \all{Ich habe das Fest gefeiert.} ✓ — le participe ferme la phrase.
\item \textbf{N'ajoute pas de ge- aux verbes en -ieren !} \all{Wir haben getanzt} ✓ mais \all{Wir haben studiert} ✓ (et jamais « gestudiert »).
\item \textbf{N'oublie pas sein pour les verbes de mouvement :} « Ich habe nach Douala gefahren » est faux ; il faut \all{Ich bin nach Douala gefahren.}
\item Attention au faux ami : \all{bekommen} signifie « recevoir », pas « devenir » : \all{Ich habe ein Geschenk bekommen.} « J'ai reçu un cadeau. » (« devenir » se dit \all{werden}.)
\end{itemize}
\end{attention}

\begin{exoresolu}[Mets les phrases au Perfekt]
Énoncé — Mets les phrases suivantes au Perfekt : 1. \all{Wir feiern das Fest.} 2. \all{Die Gäste kommen aus der Stadt.} 3. \all{Ich besuche meine Großeltern.} 4. \all{Das Fest findet im Dorf statt.} 5. \all{Sie fährt nach Yaoundé.}
\tcblower
Solution détaillée :
\begin{enumerate}
\item \all{Wir haben das Fest gefeiert.} — \all{feiern} est un verbe faible (ge + feier + t), action ordinaire : auxiliaire \all{haben}.
\item \all{Die Gäste sind aus der Stadt gekommen.} — \all{kommen} est un verbe fort de mouvement : auxiliaire \all{sein}, participe \all{gekommen}.
\item \all{Ich habe meine Großeltern besucht.} — \all{besuchen} a un préfixe inséparable \all{be-} : pas de ge-, donc \all{besucht}.
\item \all{Das Fest hat im Dorf stattgefunden.} — verbe séparable : le ge- se place entre \all{statt} et \all{gefunden}.
\item \all{Sie ist nach Yaoundé gefahren.} — \all{fahren} exprime ici un déplacement : auxiliaire \all{sein}, donc \all{ist gefahren}.
\end{enumerate}
\end{exoresolu}

\begin{savaistu}[Le Perfekt au Bac]
Au baccalauréat, le Perfekt est le temps attendu pour raconter un événement vécu : des vacances, une fête, un voyage, un week-end. Le Präteritum (\all{ich feierte, ich ging}) est réservé aux textes écrits formels (articles de presse, récits littéraires) — à l'exception des verbes \all{sein} et \all{haben}, qu'on met presque toujours au Präteritum même à l'oral : on dit naturellement \all{ich war müde} et \all{ich hatte keine Zeit} plutôt que \all{ich bin müde gewesen} ou \all{ich habe keine Zeit gehabt}. Retiens donc : raconter au passé = Perfekt + \all{war} / \all{hatte}.
\end{savaistu}

\begin{exercice}[À toi de jouer !]
1. Conjugue au Perfekt : \all{(feiern) ich \dots\ das Fest \dots} ; \all{(gehen) wir \dots\ zum Umzug \dots} ; \all{(trinken) sie \dots\ Tee \dots} ; \all{(organisieren) die Schule \dots\ eine Feier \dots} ; \all{(teilnehmen) viele Kinder \dots\ an dem Fest \dots}.

2. Traduis en allemand : « Le week-end dernier, nous avons célébré une fête traditionnelle. Ma grand-mère a cuisiné, nous avons chanté et dansé, et beaucoup de visiteurs sont venus de la ville. »
\end{exercice}

\begin{corrige}[Exercice « À toi de jouer ! »]
1. \all{Ich habe das Fest gefeiert.} / \all{Wir sind zum Umzug gegangen.} / \all{Sie hat Tee getrunken.} / \all{Die Schule hat eine Feier organisiert.} / \all{Viele Kinder haben an dem Fest teilgenommen.}

2. \all{Letztes Wochenende haben wir ein traditionelles Fest gefeiert. Meine Großmutter hat gekocht, wir haben gesungen und getanzt, und viele Besucher sind aus der Stadt gekommen.}
\end{corrige}

Tu maîtrises maintenant le temps indispensable pour raconter tes souvenirs. Dans la prochaine section, nous verrons comment exploiter méthodiquement un texte — compréhension globale, compréhension détaillée et commentaire — avant de rédiger toi-même le récit de ta fête traditionnelle.

\coursec{Méthode du commentaire et commentaire modèle}

Après avoir étudié le \cle{passé composé (Perfekt)} pour raconter des événements passés, nous abordons maintenant la compétence d'\emph{expression écrite} : rédiger un commentaire structuré sur un texte lu. En classe de Première, le commentaire ne demande pas une analyse littéraire complexe, mais une \cle{présentation claire}, un \cle{résumé ordonné} des idées et un \cle{avis personnel argumenté}, le tout en allemand correct (niveau A2).

\courssub{Les cinq étapes du commentaire de texte}

Un bon commentaire suit une démarche rigoureuse en cinq temps. Ne sautez aucune étape : la qualité de votre production écrite en dépend.

\begin{methode}[Comment rédiger un commentaire de texte (niveau A2)]
\begin{enumerate}
    \item \textbf{Présenter le texte (Einleitung).} Indiquez le \cle{thème} (Thema), le \cle{type de texte} (Textart : Artikel, Blogbeitrag, E-Mail, Dialog), et si possible la \cle{source} (Quelle) ou l'auteur. Formule type : \all{Der Text handelt von ...} (Le texte traite de ...).
    \item \textbf{Dégager l'idée principale (Kernaussage).} En une phrase, résumez le message global. Que veut nous dire l'auteur ?
    \item \textbf{Analyser la structure (Gliederung).} Repérez les \cle{parties logiques} du texte (souvent 2 ou 3). Notez les \cle{mots de liaison} (zuerst, dann, schließlich) qui guident la lecture.
    \item \textbf{Relever le vocabulaire et les procédés (Wortschatz \& Stil).} Sélectionnez 5 à 8 mots-clés (noms, verbes, adjectifs) liés au thème. Repérez les \cle{temps de verbe} (Präsens, Perfekt) et les \cle{prépositions de lieu} (in, auf, bei + Dativ) vus en grammaire.
    \item \textbf{Donner son avis argumenté (Stellungnahme).} Exprimez votre opinion personnelle en la justifiant. Comparez avec votre culture (Cameroun). Formules : \all{Meiner Meinung nach ...}, \all{Ich finde das interessant, weil ...}.
\end{enumerate}
\end{methode}

\begin{aretenir}[Règles d'or pour le commentaire A2]
\begin{itemize}
    \item \textbf{Ne recopiez jamais} des phrases entières du texte : reformulez avec vos mots (Paraphrase).
    \item \textbf{Une idée = une phrase.} Évitez les phrases trop longues avec beaucoup de « und ».
    \item \textbf{Utilisez les connecteurs} (zuerst, dann, außerdem, deshalb) pour structurer.
    \item \textbf{Terminez toujours} par votre avis personnel (Schluss), même court (2--3 phrases).
    \item \textbf{Relisez-vous :} vérifiez la majuscule des noms, la place du verbe (position 2), le Dativ après les prépositions de lieu, et la formation du Perfekt (haben/sein + Partizip II).
\end{itemize}
\end{aretenir}

\courssub{Plan type en allemand et connecteurs utiles}

Pour structurer votre paragraphe unique (ou vos trois paragraphes courts), suivez ce canevas standard. Le tableau ci-dessous oppose les parties du commentaire, leurs fonctions et les \cle{connecteurs logiques} (Konnektoren) indispensables.

\begin{center}
\begin{tabularx}{\linewidth}{|X|X|X|}
\hline
\rowcolor{popblueL}
\textbf{Partie du commentaire} & \textbf{Fonction / Phrases types (Deutsch)} & \textbf{Connecteurs clés (Konnektoren)} \\
\hline
\textbf{Einleitung} (Introduction) & \all{Der Text handelt von traditionellen Festen.} \\
\all{Es ist ein Blogbeitrag / Artikel von ...} & \cle{zuerst} (d'abord), \cle{im Text wird beschrieben} (le texte décrit) \\
\hline
\textbf{Hauptteil} (Développement) & \all{Zuerst beschreibt der Autor das Oktoberfest.} \\
\all{Dann spricht er vom Karneval in Köln.} \\
\all{Schließlich erwähnt er die Dorffeste.} & \cle{zuerst} / \cle{zunächst} (premièrement), \cle{dann} / \cle{danach} (ensuite), \cle{außerdem} (de plus), \cle{schließlich} / \cle{zuletzt} (finalement) \\
\hline
\textbf{Schluss} (Conclusion / Avis) & \all{Meiner Meinung nach sind diese Feste wichtig.} \\
\all{Das erinnert mich an das Ngondo-Fest in Kamerun.} & \cle{meiner Meinung nach} (à mon avis), \cle{ich finde, dass ...} (je trouve que), \cle{deshalb} (c'est pourquoi), \cle{zusammenfassend} (en résumé) \\
\hline
\end{tabularx}
\end{center}

\begin{exemplebox}[Phrases-modèles pour chaque partie]
\begin{itemize}
    \item \all{Der Text handelt von traditionellen Festen in Deutschland.} (Le texte traite des fêtes traditionnelles en Allemagne.)
    \item \all{Zuerst geht es um das Oktoberfest in München.} (D'abord, il s'agit de l'Oktoberfest à Munich.)
    \item \all{Dann beschreibt der Autor den Karneval in Köln.} (Ensuite, l'auteur décrit le carnaval de Cologne.)
    \item \all{Schließlich werden Feste auf dem Land genannt.} (Finalement, on cite les fêtes à la campagne.)
    \item \all{Meiner Meinung nach bringen diese Feste die Menschen zusammen.} (À mon avis, ces fêtes rassemblent les gens.)
    \item \all{Das finde ich schön, weil man die Tradition bewahrt.} (Je trouve cela beau car on préserve la tradition.)
\end{itemize}
\end{exemplebox}

\courssub{Commentaire modèle sur le texte support}

Le texte support étudié dans les sections précédentes (compréhension + lexique) présentait trois types de fêtes allemandes : l'Oktoberfest, le Karneval de Cologne et les fêtes villageoises (Schützenfest, Erntedankfest). Voici un commentaire modèle respectant la méthode (environ 9 phrases, niveau A2+).

\begin{textebox}[Commentaire modèle — Traditionelle Feste in Deutschland]
Der Text handelt von traditionellen Festen in Deutschland. Es ist ein informativer Artikel für Lernende. Zuerst beschreibt der Autor das Oktoberfest in München: die Besucher essen Brezeln, trinken Bier und hören Blasmusik in den großen Zelten. Dann spricht er vom Karneval in Köln, wo sich die Menschen verkleiden, tanzen und „Kölle Alaaf“ rufen. Schließlich erwähnt er die Feste in den Dörfern, zum Beispiel das Schützenfest mit dem Vogelschießen und das Erntedankfest mit der Erntekrone. Meiner Meinung nach sind diese Feste sehr wichtig für die Kultur, weil sie die Gemeinschaft stärken. Das erinnert mich an das Ngondo-Fest in Kamerun, wo die Sawa-Völker am Wouri-Fluss tanzen, singen und traditionelle Kleidung tragen. Ich finde es schön, dass man überall auf der Welt Feste feiert, um die Identität zu bewahren.
\end{textebox}

\begin{definition}[Glossaire du commentaire modèle]
\begin{itemize}
    \item \all{der Artikel, die Artikel} : l'article (de presse)
    \item \all{informativ} : informatif
    \item \all{die Brezel, die Brezeln} : la bretzel
    \item \all{die Blasmusik} : la musique de fanfare / cuivres
    \item \all{sich verkleiden} (Perfekt: \all{hat sich verkleidet}) : se déguiser
    \item \all{rufen} (rief, gerufen) : crier, appeler
    \item \all{das Vogelschießen} : le tir à l'oiseau (jeu de tir traditionnel)
    \item \all{die Erntekrone} : la couronne de récolte
    \item \all{stärken} : renforcer
    \item \all{die Identität, -en} : l'identité
    \item \all{bewahren} : préserver, conserver
\end{itemize}
\end{definition}

\begin{exemplebox}[Analyse des phrases clés (Allemand $\rightarrow$ Français)]
\begin{center}
\begin{tabular}{p{0.48\linewidth}p{0.48\linewidth}}
\all{Der Text handelt von traditionellen Festen.} & Le texte traite des fêtes traditionnelles. \\
\all{Zuerst beschreibt der Autor das Oktoberfest.} & D'abord, l'auteur décrit l'Oktoberfest. \\
\all{Meiner Meinung nach sind diese Feste sehr wichtig.} & À mon avis, ces fêtes sont très importantes. \\
\all{Das erinnert mich an das Ngondo-Fest.} & Cela me rappelle la fête du Ngondo. \\
\all{Ich finde es schön, dass man Feste feiert.} & Je trouve beau qu'on célèbre des fêtes.
\end{tabular}
\end{center}
\end{exemplebox}

\begin{attention}[Piège majeur : \cle{handeln von} + Dativ]
Le verbe \all{handeln} (traiter de, parler de) s'utilise \textbf{toujours} avec la préposition \cle{von} + \cle{Dativ}.
\begin{itemize}
    \item \textbf{Correct :} \all{Der Text handelt \underline{von dem} (=\underline{vom}) Fest.} (Le texte traite de la fête.)
    \item \textbf{Faux :} \all{Der Text handelt das Fest.} (Confusion avec « handeln » = agir / commercer).
    \item \textbf{Faux :} \all{Der Text handelt über das Fest.} (Anglicisme « handeln über » $\leftarrow$ *to deal with* ; en allemand standard : \all{handeln von} ou \all{gehen um}).
\end{itemize}
Autre piège : \all{erinnern an} + \cle{Akkusativ} (\all{mich an das Fest}), pas « von ».
\end{attention}

\courssub{Organigramme de la structure du commentaire}

La figure ci-dessous résume visuellement l'architecture du commentaire. Mémorisez ce cheminement : \cle{Einleitung} $\rightarrow$ \cle{Hauptteil} (chronologie ou thèmes) $\rightarrow$ \cle{Schluss} (avis + ouverture culturelle).

\begin{popfigure}
\begin{tikzpicture}[
    node distance=1.2cm and 2cm,
    box/.style={draw=popblue, thick, rounded corners, fill=popblueL, text width=3.8cm, align=center, minimum height=1.8cm, font=\small},
    arrow/.style={-Stealth, thick, popdark},
    labelbox/.style={font=\tiny, text=popmuted, align=center}
]
% Nœuds
\node[box, align=center] (einl){\textbf{Einleitung}\\\textit{(Présentation)}\\\all{Der Text handelt von ...}};
\node[box, below=of einl, align=center] (haupt){\textbf{Hauptteil}\\\textit{(Résumé structuré)}\\\all{Zuerst ... Dann ... Schließlich ...}};
\node[box, below=of haupt, align=center] (schluss){\textbf{Schluss}\\\textit{(Avis personnel)}\\\all{Meiner Meinung nach ...}};

% Flèches
\draw[arrow] (einl) -- (haupt) node[midway, right, labelbox, align=center]{Structure\\chronologique};
\draw[arrow] (haupt) -- (schluss) node[midway, right, labelbox, align=center]{Transition\\avis personnel};

% Annotations latérales
\node[left=0.3cm of einl, labelbox, text width=2cm] {Thème + Type + Source};
\node[left=0.3cm of haupt, labelbox, text width=2cm] {Idées principales\\+ Vocabulaire clé};
\node[left=0.3cm of schluss, labelbox, text width=2cm] {Opinion + Comparaison\\Cameroun/Allemagne};
\end{tikzpicture}
\legende{Structure type d'un commentaire de texte niveau A2 : trois blocs logiques reliés par des connecteurs.}
\end{popfigure}

\courssub{Entraînement oral : reformulation et appropriation}

L'écrit prépare l'oral. En binôme, reprenez le commentaire modèle sans le lire mot pour mot.

\begin{experience}[Jeu de rôle : « Journaliste culturel »]
\begin{enumerate}
    \item \textbf{Élève A (Journaliste) :} Pose les questions guides en allemand :
    \begin{itemize}
        \item \all{Wovon handelt der Text?} (De quoi traite le texte ?)
        \item \all{Was wird zuerst / dann / schließlich beschrieben?} (Qu'est-ce qui est décrit en premier / ensuite / enfin ?)
        \item \all{Was ist deine Meinung? Warum?} (Quel est ton avis ? Pourquoi ?)
        \item \all{Gibt es ein ähnliches Fest in Kamerun?} (Y a-t-il une fête similaire au Cameroun ?)
    \end{itemize}
    \item \textbf{Élève B (Expert) :} Répond en phrases complètes, en utilisant \textbf{ses propres mots} (reformulation), les connecteurs (\all{zuerst, dann, schließlich}) et le \cle{Perfekt} si nécessaire (\all{Der Autor hat beschrieben ...}).
    \item \textbf{Échange :} Inversez les rôles. Notez les difficultés (vocabulaire manquant, ordre des mots) pour les revoir.
\end{enumerate}
\end{experience}

\begin{savaistu}[Culture : Le Ngondo et les fêtes allemandes]
Le \cle{Ngondo} (fête des peuples Sawa, Douala, fin novembre/début décembre) partage avec l'\cle{Oktoberfest} et l'\cle{Erntedankfest} une fonction \cle{identitaire} et \cle{communautaire} : danses rituelles (\all{der Tanz}), chants (\all{das Lied}), tenues traditionnelles (\all{die Tracht, die traditionelle Kleidung}), offrandes (\all{das Opfer}) au fleuve Wouri (comme les couronnes de récolte \all{Erntekrone} en Allemagne). En Allemagne, le \cle{Karneval} (Cologne, Mayence, Düsseldorf) est la « cinquième saison » : inversion des rôles, humour, liberté — un peu comme les \cle{bendskins} ou le \cle{Makossa} festif au Cameroun où la musique libère la parole.
\end{savaistu}

\courssub{Exercice résolu : analyser la structure d'un commentaire}

\begin{exoresolu}[Identifier les parties et les connecteurs]
\textbf{Énoncé :} Lisez le commentaire suivant (niveau A2) sur un texte imaginaire sur la fête de la musique en Allemagne (\all{Tag der Musik}). Repérez : 1) la phrase d'introduction, 2) les trois idées du développement avec leurs connecteurs, 3) la phrase de conclusion avec l'avis personnel. Surlignez les connecteurs.

\all{Der Text handelt vom „Tag der Musik“ in Berlin. Zuerst erklärt der Autor, dass am 21. Juni viele Musiker auf der Straße spielen. Dann beschreibt er die Stimmung: die Leute tanzen und singen zusammen. Schließlich sagt er, dass das Fest kostenlos ist. Meiner Meinung nach ist das eine tolle Idee, weil Musik alle Menschen verbindet. In Kamerun feiern wir auch die Musik, zum Beispiel beim Festival Mützig Star.}

\tcblower

\textbf{Solution détaillée :}
\begin{enumerate}
    \item \textbf{Einleitung :} \all{Der Text handelt vom „Tag der Musik“ in Berlin.} (Présentation thème + lieu ; \cle{vom} = \cle{von dem} Dativ).
    \item \textbf{Hauptteil (3 idées + connecteurs) :}
    \begin{itemize}
        \item Idée 1 : \all{\textbf{Zuerst} erklärt der Autor, dass am 21. Juni viele Musiker auf der Straße spielen.} (Connecteur : \cle{Zuerst}).
        \item Idée 2 : \all{\textbf{Dann} beschreibt er die Stimmung: die Leute tanzen und singen zusammen.} (Connecteur : \cle{Dann}).
        \item Idée 3 : \all{\textbf{Schließlich} sagt er, dass das Fest kostenlos ist.} (Connecteur : \cle{Schließlich}).
    \end{itemize}
    \item \textbf{Schluss (Avis + comparaison) :} \all{\textbf{Meiner Meinung nach} ist das eine tolle Idee, \textbf{weil} Musik alle Menschen verbindet. In Kamerun feiern wir auch die Musik, zum Beispiel beim Festival Mützig Star.}
    \begin{itemize}
        \item Marqueur d'opinion : \cle{Meiner Meinung nach}.
        \item Justification : \cle{weil} + verbe à la fin (\all{verbindet}).
        \item Ouverture culturelle : \cle{In Kamerun ...} (Präsens pour fait général).
    \end{itemize}
\end{enumerate}
\textbf{Bilan :} Structure parfaite (Einleitung -- Hauptteil avec 3 étapes chronologiques -- Schluss avec avis justifié + comparaison). Vocabulaire précis (\all{die Stimmung, kostenlos, verbinden}). Grammaire : \cle{Perfekt} absent ici (texte au \cle{Präsens} car description d'un fait général), \cle{weil}-phrase correcte (verbe final).
\end{exoresolu}

\begin{exercice}[À vous de jouer : Rédigez votre commentaire]
Sur la base du texte support « Les fêtes traditionnelles » (sections 1--2) et du modèle ci-dessus, rédigez \textbf{votre propre commentaire} (8--10 phrases) dans votre cahier.
\textbf{Consignes :}
\begin{enumerate}
    \item Respectez la structure : \all{Einleitung -- Hauptteil (3 parties) -- Schluss}.
    \item Utilisez au moins 4 connecteurs différents (\all{zuerst, dann, außerdem, schließlich, meiner Meinung nach}).
    \item Intégrez 3 mots du lexique de la leçon (\all{die Tracht, das Schützenfest, der Umzug, die Blasmusik, sich verkleiden}).
    \item Terminez par une comparaison personnelle avec une fête camerounaise (Ngondo, Tabaski, Noël, Fête de la Jeunesse, etc.).
    \item Vérifiez : majuscules noms, ordre verbe (position 2), \cle{Dativ} après \all{von, bei, in, auf, mit}, \cle{Perfekt} si vous racontez un souvenir personnel.
\end{enumerate}
\end{exercice}

\begin{corrige}[Exercice : Rédigez votre commentaire]
\textbf{Proposition de correction (exemple élève) :}
\all{Der Text handelt von traditionellen Festen in Deutschland. Zuerst geht es um das Oktoberfest in München mit der Blasmusik und den Brezeln. Dann beschreibt der Autor den Karneval in Köln, wo sich die Leute verkleiden und beim Umzug tanzen. Außerdem erwähnt er das Schützenfest auf dem Land mit dem Vogelschießen. Meiner Meinung nach sind diese Traditionen sehr wichtig für die Identität. Das erinnert mich an den Jugendtag in Kamerun am 11. Februar, wo Umzüge und Tänze unseren nationalen Stolz zeigen. Ich finde das gut, weil Feste die Kultur bewahren.}
\\
\textbf{Points forts :} Structure respectée, connecteurs variés (\all{zuerst, dann, außerdem}), vocabulaire ciblé (\all{Blasmusik, Umzug, sich verkleiden, Schützenfest, Vogelschießen}), avis justifié (\all{weil}), comparaison culturelle pertinente (\all{der Jugendtag}). \textbf{Attention :} \all{leben lassen} (laisser vivre) est correct mais \all{bewahren} (préserver) ou \all{pflegen} (entretenir) est plus précis pour la culture. Vérifiez toujours l'orthographe des noms propres allemands (\all{Jugendtag} et non \all{la Fête de la Jeunesse}).
\end{corrige}

\coursec{Production guidée : ma fête traditionnelle}

Après avoir appris à lire, comprendre et commenter un texte sur les fêtes traditionnelles, il est temps de passer à l'étape la plus concrète de notre leçon : \cle{produire} ton propre texte. C'est exactement ce que l'on te demandera au baccalauréat et dans la vie réelle : savoir \all{beschreiben} (décrire) et \all{erzählen} (raconter) une fête de ton pays en allemand. Le Cameroun est un trésor de fêtes traditionnelles — le Ngondo des Sawa à Douala, le Nguon des Bamoun à Foumban, la fête de la jeunesse du 11 février, les mariages traditionnels — et tu vas apprendre à en parler comme un vrai germanophone. Pas de panique : nous allons procéder pas à pas, avec un plan, des phrases modèles et une grille de vérification.

\courssub{Objectif de la production}

\begin{definition}[La tâche à accomplir]
Tu dois rédiger un \cle{paragraphe} de \textbf{6 à 8 phrases} en allemand sur une fête traditionnelle camerounaise de ton choix (le Ngondo, le Nguon, la fête de la jeunesse, un mariage traditionnel), en suivant un \cle{plan-guide} en cinq étapes. Variante : raconter la dernière édition de la fête au \all{Perfekt}. Puis tu présentes ton texte \textbf{oralement} en une minute devant la classe, avec une affiche ou une photo.
\end{definition}

\textbf{Le plan-guide en cinq étapes :}

\begin{enumerate}
\item \all{Name und Ort} : le nom de la fête et le lieu où elle a lieu.
\item \all{Datum und Dauer} : la date et la durée.
\item \all{Aktivitäten} : les activités (danses, musique, repas).
\item \all{Kleidung} : les vêtements traditionnels.
\item \all{Bedeutung} : l'importance de la fête et ton avis personnel.
\end{enumerate}

\begin{popfigure}
\begin{center}
\begin{tikzpicture}[font=\small]
\node[draw=popblue, fill=popblueL, rounded corners, minimum width=3.4cm, minimum height=0.9cm, align=center] (s1) at (0,0) {\textbf{1. Name und Ort}\\ \all{Das Fest findet in ... statt.}};
\node[draw=popgreen, fill=popgreenL, rounded corners, minimum width=3.4cm, minimum height=0.9cm, align=center] (s2) at (4.2,-1.2) {\textbf{2. Datum und Dauer}\\ \all{Es dauert ... Tage.}};
\node[draw=poporange, fill=poporangeL, rounded corners, minimum width=3.4cm, minimum height=0.9cm, align=center] (s3) at (8.4,-2.4) {\textbf{3. Aktivitäten}\\ \all{Man tanzt, isst, singt ...}};
\node[draw=poppurple, fill=poppurpleL, rounded corners, minimum width=3.4cm, minimum height=0.9cm, align=center] (s4) at (4.2,-3.6) {\textbf{4. Kleidung}\\ \all{Die Menschen tragen ...}};
\node[draw=poppink, fill=poppinkL, rounded corners, minimum width=3.4cm, minimum height=0.9cm, align=center] (s5) at (0,-4.8) {\textbf{5. Bedeutung}\\ \all{Ich finde dieses Fest ..., weil ...}};
\draw[-{Stealth}, thick, popdark] (s1.east) -- (s2.west);
\draw[-{Stealth}, thick, popdark] (s2.east) -- (s3.west);
\draw[-{Stealth}, thick, popdark] (s3.west) -- (s4.east);
\draw[-{Stealth}, thick, popdark] (s4.west) -- (s5.east);
\end{tikzpicture}
\end{center}
\legende{Les cinq étapes du paragraphe, en escalier : du nom de la fête jusqu'à ton avis personnel.}
\end{popfigure}

\courssub{La banque de phrases modèles}

Voici ta boîte à outils. Ces phrases à trous sont des \cle{gabarits} : tu remplaces les pointillés par tes propres informations. Apprends-les par cœur, elles te serviront aussi à l'oral.

\begin{center}
\begin{tabularx}{\linewidth}{|X|X|}
\hline
\rowcolor{popblueL} \textbf{Phrase modèle (Deutsch)} & \textbf{Français} \\
\hline
\all{Das Ngondo-Fest findet in Douala statt.} & La fête Ngondo a lieu à Douala. \\
\hline
\all{Es findet jedes Jahr im Dezember statt.} & Elle a lieu chaque année en décembre. \\
\hline
\all{Das Fest dauert eine Woche.} & La fête dure une semaine. \\
\hline
\all{Die Menschen tanzen und singen am Wasser.} & Les gens dansent et chantent au bord de l'eau. \\
\hline
\all{Man isst Fisch und trinkt Palmwein.} & On mange du poisson et on boit du vin de palme. \\
\hline
\all{Die Menschen tragen traditionelle Kleidung.} & Les gens portent des vêtements traditionnels. \\
\hline
\all{Ich finde dieses Fest wichtig, weil es die Menschen zusammenbringt.} & Je trouve cette fête importante, parce qu'elle rassemble les gens. \\
\hline
\end{tabularx}
\end{center}

\begin{exemplebox}[Exemples traduits à imiter]
\all{Das Ngondo-Fest findet in Douala am Wasser statt.} « La fête Ngondo a lieu à Douala au bord de l'eau. » — Remarque le verbe \all{stattfinden} (avoir lieu), verbe à particule séparable : la particule \all{statt} va à la fin de la phrase.\par
\all{Letztes Jahr haben wir traditionelle Kleidung getragen und getanzt.} « L'année dernière, nous avons porté des vêtements traditionnels et dansé. » — Au \all{Perfekt}, l'auxiliaire \all{haben} est en deuxième position et les participes (\all{getragen}, \all{getanzt}) vont à la fin.
\end{exemplebox}

\begin{attention}[Le piège de « weil » : le verbe va à la fin !]
La conjonction \all{weil} (parce que) introduit une \textbf{subordonnée} : le verbe conjugué se place \textbf{à la toute fin} de la proposition.\par
Correct : \all{Ich finde dieses Fest wichtig, weil es die Menschen zusammenbringt.}\par
Faux : \emph{weil es bringt die Menschen zusammen.}\par
En français, l'ordre ne change pas après « parce que » ; en allemand, si. C'est l'erreur numéro un des francophones. Vérifie toujours : après \all{weil}, \all{dass}, \all{obwohl}, le verbe est à la fin.
\end{attention}

\courssub{Un modèle de production écrite}

Lis d'abord ce texte modèle. Il fait exactement 8 phrases et respecte le plan en cinq étapes. Observe la place des verbes, le datif après les prépositions de lieu et le vocabulaire du thème.

\begin{textebox}[Mein Lieblingsfest : das Ngondo-Fest]
\all{Das Ngondo-Fest findet in Douala statt, am Wasser des Flusses Wouri. Es findet jedes Jahr im Dezember statt und dauert eine Woche. Viele Menschen kommen zu dem Fest, auch Touristen aus Deutschland und Frankreich. Die Sawa tanzen traditionelle Tänze und singen alte Lieder. Es gibt auch ein Rennen mit den Pirogen auf dem Fluss. Man isst Fisch, Kochbananen und man trinkt Palmwein. Die Menschen tragen schöne traditionelle Kleidung, zum Beispiel den Kaba. Ich finde dieses Fest sehr wichtig, weil es die Traditionen unserer Väter zeigt und die Menschen zusammenbringt.}
\end{textebox}

\textbf{Glossaire :} \all{der Fluss, die Flüsse} (le fleuve) ; \all{das Rennen, die Rennen} (la course) ; \all{die Piroge, die Pirogen} (la pirogue) ; \all{die Kochbanane, die Kochbananen} (la banane plantain) ; \all{der Palmwein} (le vin de palme, pas de pluriel) ; \all{der Kaba, die Kaba} (le kaba, robe traditionnelle) ; \all{zeigen} (montrer) ; \all{zusammenbringen} (rassembler, particule séparable : \all{bringt ... zusammen}).

Observe bien : \all{zu dem Fest} et \all{auf dem Fluss} sont des \textbf{datifs} après des prépositions de lieu (\all{zu}, \all{auf} + datif, car il s'agit d'un lieu où l'on est, pas d'un mouvement). \all{mit den Pirogen} : \all{mit} est toujours suivi du datif. \all{vor dem Palast} (dans l'exercice ci-dessous) : \all{vor} + datif (lieu).

\courssub{La variante au passé : raconter au Perfekt}

Pour raconter la \textbf{dernière édition} de la fête, on utilise le \all{Perfekt} : auxiliaire \all{haben} ou \all{sein} (pour les verbes de mouvement et changement d'état) en deuxième position + participe passé à la fin.

\begin{center}
\begin{tabularx}{\linewidth}{|X|X|X|}
\hline
\rowcolor{popblueL} \textbf{Präsens (habitude)} & \textbf{Perfekt (dernière édition)} & \textbf{Français} \\
\hline
\all{Ich feiere das Fest.} & \all{Ich habe das Fest gefeiert.} & J'ai fêté la fête. \\
\hline
\all{Wir gehen nach Douala.} & \all{Wir sind nach Douala gegangen.} & Nous sommes allés à Douala. \\
\hline
\all{Die Menschen tanzen.} & \all{Die Menschen haben getanzt.} & Les gens ont dansé. \\
\hline
\all{Wir tragen Kleidung.} & \all{Wir haben Kleidung getragen.} & Nous avons porté des vêtements. \\
\hline
\all{Man isst Fisch.} & \all{Man hat Fisch gegessen.} & On a mangé du poisson. \\
\hline
\end{tabularx}
\end{center}

\begin{exemplebox}[Deux phrases de départ pour ta variante au passé]
\all{Letztes Jahr habe ich das Ngondo-Fest mit meiner Familie gefeiert.} « L'année dernière, j'ai fêté la fête Ngondo avec ma famille. »\par
\all{Wir sind nach Douala gegangen und haben den ganzen Tag getanzt.} « Nous sommes allés à Douala et avons dansé toute la journée. »
\end{exemplebox}

\begin{methode}[Réussir sa production en quatre étapes]
\begin{enumerate}
\item \textbf{Avant d'écrire}, note sur ton brouillon \textbf{5 mots de vocabulaire} du thème que tu veux utiliser (par exemple : \all{das Fest, der Tanz, die Kleidung, das Essen, feiern}).
\item \textbf{Écris des phrases courtes} : sujet + verbe en deuxième position + compléments. Une idée par phrase. Mieux vaut 8 phrases simples et correctes que 3 phrases longues et fausses.
\item \textbf{Vérifie la place du verbe} : verbe conjugué en position 2 ; participe passé à la fin au \all{Perfekt} ; verbe à la fin après \all{weil}.
\item \textbf{Relis-toi avec la grille d'évaluation} ci-dessous : as-tu bien 3 mots du thème, 2 datifs, 3 verbes au \all{Perfekt} (pour la variante) ?
\end{enumerate}
\end{methode}

\begin{exoresolu}[Complète puis rédige]
\textbf{Énoncé.} Complète les phrases avec les éléments entre parenthèses, puis assemble-les pour obtenir un paragraphe sur le Nguon.\par
1. Das Nguon-Fest \_\_\_ (findet / statt) in Foumban. 2. Es \_\_\_ (dauern) zwei Tage. 3. Die Menschen \_\_\_ (tragen) traditionelle Kleidung. 4. Man \_\_\_ (tanzen) und \_\_\_ (singen) vor dem Palast des Sultans. 5. Ich finde dieses Fest interessant, weil es sehr alt \_\_\_ (sein).
\tcblower
\textbf{Solution.} 1. \all{Das Nguon-Fest findet in Foumban statt.} (verbe à particule séparable : \all{statt} à la fin). 2. \all{Es dauert zwei Tage.} 3. \all{Die Menschen tragen traditionelle Kleidung.} (\all{tragen}, verbe fort : \all{sie tragen}). 4. \all{Man tanzt und singt vor dem Palast des Sultans.} (\all{vor dem Palast} : datif de lieu). 5. \all{Ich finde dieses Fest interessant, weil es sehr alt ist.} (après \all{weil}, le verbe \all{ist} va à la fin).\par
\textbf{Paragraphe final :} \all{Das Nguon-Fest findet in Foumban statt. Es dauert zwei Tage. Die Menschen tragen traditionelle Kleidung. Man tanzt und singt vor dem Palast des Sultans. Ich finde dieses Fest interessant, weil es sehr alt ist.} — Cinq phrases simples, correctes, avec le vocabulaire du thème, un datif et une subordonnée avec \all{weil} : c'est déjà une excellente base que tu peux enrichir.
\end{exoresolu}

\courssub{La production orale : une minute devant la classe}

\begin{experience}[Présentation orale guidée]
Par groupes de deux, puis individuellement :
\begin{enumerate}
\item Choisis une fête et prépare ton paragraphe écrit (6 à 8 phrases).
\item Réalise une petite \textbf{affiche} ou apporte une \textbf{photo} : le nom de la fête en allemand, trois images (danses, repas, vêtements) et trois mots-clés.
\item Présente ta fête en \textbf{une minute} devant la classe, sans lire tout le texte : regarde tes camarades, montre les images sur ton affiche.
\item La classe pose ensuite une question : \all{Wann findet das Fest statt?} (Quand la fête a-t-elle lieu ?) ou \all{Was isst man?} (Qu'est-ce qu'on mange ?).
\end{enumerate}
Conseil de prononciation : \all{Ngondo} se dit « n-go-ndo », \all{Kleidung} « klaï-doung », \all{wichtig} « vikh-tikh ».
\end{experience}

\courssub{La grille d'évaluation}

Utilise cette grille pour t'auto-évaluer avant de rendre ton texte ou de passer à l'oral.

\begin{center}
\begin{tabularx}{\linewidth}{|X|c|}
\hline
\rowcolor{popblueL} \textbf{Critère} & \textbf{Points} \\
\hline
Vocabulaire du thème : au moins 3 mots (\all{das Fest, die Tradition, der Tanz, das Essen, die Kleidung, feiern}...) & 0 à 2 \\
\hline
Datif correct : au moins 2 exemples (\all{in Douala, zu dem Fest, auf dem Fluss, mit den Pirogen, vor dem Palast}) & 0 à 2 \\
\hline
\all{Perfekt} correct : au moins 3 verbes (variante au passé) & 0 à 2 \\
\hline
Place du verbe respectée (position 2, participe à la fin, verbe final après \all{weil}) & 0 à 2 \\
\hline
Fluidité et prononciation à l'oral (une minute, sans lire) & 0 à 2 \\
\hline
\textbf{Total} & \textbf{10} \\
\hline
\end{tabularx}
\end{center}

\begin{exercice}[À toi de jouer !]
Rédige un paragraphe de 6 à 8 phrases sur la fête traditionnelle de ton choix (Ngondo, Nguon, fête de la jeunesse, mariage traditionnel) en suivant le plan en cinq étapes. Utilise au moins trois phrases de la banque de phrases modèles, deux datifs et une phrase avec \all{weil}. Variante : raconte la dernière édition au \all{Perfekt} avec au moins trois verbes conjugués.
\end{exercice}

\begin{savaistu}[Le Nguon, une fête vieille de plus de 600 ans]
Le \all{Nguon} des Bamoun, à Foumban dans la région de l'Ouest, est une fête vieille de plus de 600 ans, célébrée autour du sultan-roi. Elle est inscrite au patrimoine culturel immatériel de l'humanité. Sache que décrire une fête camerounaise comme le Nguon ou le Ngondo est un sujet classique et très apprécié au baccalauréat : c'est l'occasion de montrer ton vocabulaire, ton datif et ton \all{Perfekt} tout en valorisant la culture de ton pays. Un sujet en or !
\end{savaistu}

\begin{aretenir}[L'essentiel de la production guidée]
\begin{itemize}
\item Suis le plan en cinq étapes : nom et lieu, date et durée, activités, vêtements, avis personnel.
\item Utilise les phrases gabarits : \all{Das Fest findet in ... statt. / Es dauert ... / Die Menschen tragen ... / Man isst ... / Ich finde dieses Fest ..., weil ...}
\item Au présent : verbe en position 2 ; au \all{Perfekt} : auxiliaire en position 2, participe à la fin ; après \all{weil} : verbe à la fin.
\item Prépare 5 mots de vocabulaire avant d'écrire et relis-toi avec la grille d'évaluation.
\end{itemize}
\end{aretenir}

\newpage

\coursec{Activité d'intégration}

\coursec{ALA\_21 : Les fêtes traditionnelles}

\courssub{Objectifs de l'activité}
Cette activité de synthèse vise à mobiliser l'ensemble des savoirs et savoir-faire travaillés dans la leçon ALA\_21 dans une tâche de communication authentique et motivante. À l'issue de cette séance, les élèves doivent être capables de :
\begin{itemize}
    \item \textbf{Comprendre} une consigne complexe et un document de cadrage en allemand (niveau A1+/A2).
    \item \textbf{Mobiliser} le lexique thématique des fêtes traditionnelles (\all{das Fest, die Tradition, der Tanz, das Essen, die Kleidung, feiern, der Brauch, die Musik, der Umzug}) avec les articles et pluriels corrects.
    \item \textbf{Appliquer} la grammaire du chapitre : le \textbf{datif} après les prépositions de lieu (\all{in, an, auf, vor, hinter, unter, über, neben, zwischen}) pour localiser l'événement et décrire les positions statiques (\all{Wo?}), et le \textbf{passé composé (Perfekt)} avec auxiliaires \all{haben} et \all{sein} pour raconter le déroulement de la dernière édition.
    \item \textbf{Produire} un texte cohérent de 8 à 10 phrases (description au présent + récit au Perfekt) et le présenter à l'oral en 2 minutes.
    \item \textbf{Coopérer} en groupe (répartition des rôles : scripteur, correcteur, présentateur, chercheur de vocabulaire) et évaluer les productions des pairs selon des critères clairs.
\end{itemize}

\courssub{Situation}
\begin{textebox}[Contexte : Projet d'échange scolaire Yaoundé – Leipzig]
\emph{Contexte :} Le \all{Gymnasium Leibniz} à Leipzig (Saxe) et le Lycée Général Leclerc de Yaoundé préparent un échange scolaire pour l'année prochaine. Les correspondants allemands souhaitent découvrir la richesse culturelle du Cameroun avant leur voyage. Votre classe de Première A a pour mission de réaliser un « Guide des fêtes camerounaises » (Titre : \all{Feste in Kamerun – Ein Reiseführer für Freunde}).

\all{Aufgabe: Arbeitet in Gruppen zu viert. Wählt eine bekannte kamerunische Tradition (z. B. Ngondo in Douala, Nguon bei den Bamoun, Nyem-Nyem in der Region Adamaoua, die Jugendfeier am 11. Februar, das Ngondo-Fest der Sawa, das Medumba-Fest). Erstellt eine Seite für den Guide.}

\all{Die Seite muss enthalten:}
\begin{enumerate}
    \item \all{Einen Titel mit dem Namen des Festes und der Region/Volksgruppe.}
    \item \all{4–5 Sätze im \textbf{Präsens}: Beschreibung des Ortes (Wo? – Dativ mit Präpositionen), der Zeit, der Kleidung, der Musik, des Essens, der Tänze.}
    \item \all{4–5 Sätze im \textbf{Perfekt}: Bericht über die letzte Ausgabe (Letztes Jahr / Vor zwei Wochen …). Was habt ihr gemacht? Was habt ihr gegessen? Wo seid ihr hingegangen?}
\end{enumerate}
\all{Präsentiert eure Seite der Klasse (ca. 2 Minuten pro Gruppe). Die Klasse stimmt ab: Welches Fest ist am attraktivsten? Alle Seiten werden zu einem Klassen-Guide zusammengefügt und im Flur ausgehängt.}
\end{textebox}

\begin{savaistu}[Contexte culturel : Le Ngondo et le Nguon]
Le \textbf{Ngondo} (peuple Sawa, Douala) est une grande fête annuelle sur les rives du Wouri : pirogues décorées, culte des ancêtres (Jengu), danses, lutte traditionnelle. Le \textbf{Nguon} (peuple Bamoun, Foumban) est une fête triennale majeure : couronnement symbolique, défilés, masques, marché artisanat. La \textbf{Fête de la Jeunesse} (11 février) est nationale : défilés scolaires, chants, sports, discours officiels. Ces fêtes illustrent la diversité des \all{die Traditionen} au Cameroun.
\end{savaistu}

\courssub{Consignes}
\emph{Travail en groupes de 4 (20 min rédaction + 10 min répétition + présentations).}

\begin{enumerate}
    \item \textbf{Choix et préparation (5 min) :} Choisissez votre fête. Désignez les rôles : \textit{Rédacteur principal} (écrit la version finale), \textit{Correcteur grammaire} (vérifie Dativ/Perfekt), \textit{Recherche vocabulaire} (dictionnaire/lexique leçon), \textit{Présentateur} (parle au nom du groupe).
    \item \textbf{Rédaction – Phase 1 : Description au présent (10 min).} Rédigez 4–5 phrases. Utilisez \textbf{impérativement} :
    \begin{itemize}
        \item 3 prépositions de lieu différentes (\all{in, an, auf, vor, hinter, neben}) avec le \textbf{Datif} (\all{Wo?}).
        \item Du vocabulaire : \all{die traditionelle Kleidung, die Trommel, der Tanz, das Gericht, die Menge, die Atmosphäre}.
        \item Exemple structure : \all{Das Ngondo-Fest findet \textbf{am Ufer des Wouri} statt. Die Leute tragen \textbf{bunte traditionelle Kleidung}.}
    \end{itemize}
    \item \textbf{Rédaction – Phase 2 : Récit au Perfekt (10 min).} Rédigez 4–5 phrases sur la \textbf{dernière édition} (\all{letztes Jahr} ou \all{vor zwei Wochen}). Utilisez \textbf{impérativement} :
    \begin{itemize}
        \item 3 verbes au Perfekt avec \all{haben} (ex. \all{gegessen, getanzt, gesehen, gekauft, fotografiert}).
        \item 2 verbes de mouvement au Perfekt avec \all{sein} (ex. \all{sind ... gegangen, sind ... gefahren, sind ... gelaufen}).
        \item Marqueurs temporels : \all{letztes Jahr, damals, am Festtag, abends}.
    \end{itemize}
    \item \textbf{Mise en page (5 min) :} Recopiez proprement sur une feuille A4 (titre, nom du groupe, petites illustrations bienvenues). Relisez : majuscules aux noms, Umlaute (ä, ö, ü), ß, ordre des mots (verbe en 2e position, participe à la fin).
    \item \textbf{Présentation orale (2 min/groupe) :} Le présentateur lit ou dit le texte. Les autres groupes notent sur une grille : \textit{Vocabulaire riche ? / Dativ correct ? / Perfekt correct ? / Intérêt culturel ? / Prononciation ?}
    \item \textbf{Vote et assemblage (5 min) :} Vote à main levée pour la fête la plus attrayante. Assemblage des pages en un guide de classe affiché au tableau/au couloir.
\end{enumerate}

\courssub{Ressources à mobiliser}
\begin{propriete}[Rappel : Datif \& Prépositions de lieu (Wo? – position statique)]
Les prépositions \cle{in, an, auf, vor, hinter, unter, über, neben, zwischen} gouvernent le \textbf{Datif} pour indiquer \textbf{OÙ} quelque chose/quelqu'un se trouve (pas de mouvement vers).
\begin{center}
\begin{tabularx}{\linewidth}{|X|X|X|}
\hline
\rowcolor{popblueL} \textbf{Préposition} & \textbf{Article Datif (masc./neutre)} & \textbf{Exemple (Fête)} \\
\hline
\all{in} & \all{im} (= in dem) & \all{im Dorf, im Stadion, im Februar} \\
\hline
\all{an} & \all{am} (= an dem) & \all{am Fluss, am Ufer, am Platz} \\
\hline
\all{auf} & \all{auf dem} & \all{auf dem Platz, auf der Bühne (fem.)} \\
\hline
\all{vor} & \all{vor dem} & \all{vor dem Palast, vor der Bühne} \\
\hline
\all{hinter} & \all{hinter dem} & \all{hinter dem Haus} \\
\hline
\all{neben} & \all{neben dem} & \all{neben dem Markt} \\
\hline
\all{zwischen} & \all{zwischen dem ... und dem ...} & \all{zwischen dem Fluss und dem Wald} \\
\hline
\end{tabularx}
\end{center}
\emph{Attention :} \all{in der} (fém.), \all{an der} (fém.), \all{auf der} (fém.), \all{vor der} (fém.). Pas de contraction \all{in der} $\to$ \all{inder} (n'existe pas).
\end{propriete}

\begin{propriete}[Rappel : Le Perfekt (Passé composé) – Structure]
\all{Subjekt + Hilfsverb (haben/sein, konjugué P2) + ... + Partizip II (à la fin).}
\begin{itemize}
    \item \textbf{Auxiliaire \all{haben}} : verbes transitifs + la plupart intransitifs (\all{feiern $\to$ gefeiert, essen $\to$ gegessen, tanzen $\to$ getanzt, sehen $\to$ gesehen, kaufen $\to$ gekauft, singen $\to$ gesungen}).
    \item \textbf{Auxiliaire \all{sein}} : verbes de \textbf{mouvement} (aller quelque part) + changement d'état (\all{gehen $\to$ gegangen, fahren $\to$ gefahren, laufen $\to$ gelaufen, fliegen $\to$ geflogen, kommen $\to$ gekommen, ankommen $\to$ angekommen, aufstehen $\to$ aufgestanden}).
    \item \textbf{Participe II} : verbes faibles \all{ge...t} (\all{gefeiert}), verbes forts \all{ge...en} + souvent voyelle modifiée (\all{gegessen, gesehen, gefunden, geschrieben}), verbes mixtes (\all{gebracht, gedacht}). Verbes en \all{-ieren} : pas de \all{ge-} (\all{fotografiert, telefoniert, studiert}).
\end{itemize}
\end{propriete}

\begin{definition}[Lexique utile – Fêtes camerounaises (Article + Pluriel)]
\begin{center}
\begin{tabularx}{\linewidth}{|l|l|X|}
\hline
\rowcolor{poporangeL} \textbf{Deutsch} & \textbf{Français} & \textbf{Exemple d'usage} \\
\hline
\all{das Fest, die Feste} & la fête & \all{Ein großes Fest} \\
\all{die Tradition, die Traditionen} & la tradition & \all{alte Traditionen pflegen} \\
\all{der Brauch, die Bräuche} & la coutume & \all{ein alter Brauch} \\
\all{der Tanz, die Tänze} & la danse & \all{traditionelle Tänze aufführen} \\
\all{die Musik} & la musique & \all{livemusik, Trommelmusik} \\
\all{die Trommel, die Trommeln} & le tam-tam & \all{auf der Trommel spielen} \\
\all{das Essen, -} / \all{das Gericht, die Gerichte} & le repas / le plat & \all{typische Gerichte: Ndolé, Koki, Poulet DG} \\
\all{die Kleidung} / \all{das Gewand, die Gewänder} & le vêtement / le costume & \all{traditionelle Kleidung tragen} \\
\all{der Umzug, die Umzüge} & le défilé & \all{ein bunter Umzug} \\
\all{die Menge, die Mengen} & la foule & \all{viele Menschen, die Menge jubelt} \\
\all{der Fluss, die Flüsse} & le fleuve & \all{der Wouri, der Sanaga} \\
\all{das Ufer, die Ufer} & la rive & \all{am Ufer stehen} \\
\all{der Palast, die Paläste} & le palais & \all{der Sultanpalast in Foumban} \\
\all{feiern} & fêter & \all{wir feiern das Fest} \\
\all{aufführen} & présenter (spectacle) & \all{einen Tanz aufführen} \\
\all{schmecken} & avoir bon goût & \all{Das Essen schmeckt gut.} \\
\hline
\end{tabularx}
\end{center}
\end{definition}

\begin{attention}[Pièges fréquents pour francophones]
\begin{itemize}
    \item \textbf{Oubli du Datif :} \all{*in der Dorf} $\to$ \all{im Dorf} ; \all{*an der Fluss} $\to$ \all{am Fluss} (masc.). Pensez \textbf{dem $\to$ m} (im, am, vom, zum).
    \item \textbf{Confusion haben/sein au Perfekt :} \all{*Wir haben nach Douala gefahren} $\to$ \all{Wir sind nach Douala gefahren} (mouvement). \all{*Wir sind das Fest gefeiert} $\to$ \all{Wir haben das Fest gefeiert} (transitif).
    \item \textbf{Ordre des mots (Perfekt) :} \all{*Wir haben gegessen das Ndolé} $\to$ \all{Wir haben das Ndolé gegessen} (Participe II tout à la fin).
    \item \textbf{Participe II verbes forts :} \all{sehen $\to$ gesehen} (pas *gesehn), \all{geben $\to$ gegeben}, \all{finden $\to$ gefunden}, \all{lesen $\to$ gelesen}.
    \item \textbf{Majuscules :} Tous les noms : \all{das Fest, die Tradition, der Tanz}. Pas de majuscule aux verbes/in adjectifs : \all{wir feiern traditionell}.
\end{itemize}
\end{attention}

\courssub{Proposition de production et corrigé commenté}
\begin{exoresolu}[Production modèle : Le Ngondo (Groupe « Wouri »)]
\emph{Énoncé :} Rédigez la page « Ngondo » pour le guide „Feste in Kamerun“.
\tcblower
\textbf{Production élève (niveau A1+/A2 visé) :}

\all{Unsere Seite: Das Ngondo-Fest der Sawa in Douala.}

\all{1. Das Ngondo-Fest findet \textbf{jedes Jahr im August} \textbf{am Ufer des Wouri} statt.} \\
\all{2. Viele Menschen \textbf{tragen bunte traditionelle Kleidung} und \textbf{schmuckvolle Gewänder}.} \\
\all{3. Auf dem Wasser \textbf{fahren geschmückte Pirogen} und \textbf{trommeln die Musiker}.} \\
\all{4. Die Menge \textbf{tanzt zu lauter Musik} und \textbf{isst typische Gerichte wie Ndolé und Plantains}.} \\
\all{5. Die Atmosphäre \textbf{ist fröhlich und feierlich}.}

\all{6. \textbf{Letztes Jahr sind wir mit dem Bus nach Douala gefahren}.} \\
\all{7. Wir \textbf{sind früh am Morgen am Ufer angekommen}.} \\
\all{8. Wir \textbf{haben den großen Umzug mit den Masken gesehen}.} \\
\all{9. Meine Freunde und ich \textbf{haben Ndolé und gegrillten Fisch gegessen}.} \\
\all{10. Abends \textbf{sind wir müde, aber glücklich nach Hause gegangen}.}

\textbf{Commentaire détaillé des choix de langue (pour l'enseignant / auto-correction élève) :}

\begin{itemize}
    \item \textbf{Phrases 1–5 (Présent – Description) :}
    \begin{itemize}
        \item \all{jedes Jahr im August} : complément de temps (accusatif de durée/fréquence, pas de préposition, ou \all{im} = in dem mois).
        \item \all{am Ufer des Wouri} : \textbf{Dativ} après \all{an} (Wo? position statique). \all{des} = Genitiv (appartenance : rive \textbf{du} Wouri). \textbf{Point clé :} distinction Dativ (lieu) / Genitiv (appartenance).
        \item \all{tragen ... Kleidung} : verbe \all{tragen} (irrégulier : trägt, trug, getragen) + accusatif.
        \item \all{fahren ... Pirogen} : verbe \all{fahren} au présent (fährt) mais ici sujet pluriel \all{Pirogen} $\to$ \all{fahren}. Mouvement sur l'eau $\to$ Perfekt avec \all{sein} plus tard.
        \item \all{tanzt ... isst} : verbes forts au présent (\all{tanzen} régulier, \all{essen} $\to$ isst).
    \end{itemize}
    \item \textbf{Phrases 6–10 (Perfekt – Récit) :}
    \begin{itemize}
        \item \all{sind ... gefahren} (Phrase 6) : \textbf{sein} + Participe II \all{gefahren} (verbe de mouvement \all{fahren}). \all{mit dem Bus} (instrument/moyen, Dativ). \all{nach Douala} (direction, pas d'article pays/ville).
        \item \all{sind ... angekommen} (Phrase 7) : \textbf{sein} + \all{angekommen} (verbe à particule inséparable \all{an|kommen} $\to$ pas de \all{ge-}, juste \all{angekommen}). \all{früh am Morgen} (temps). \all{am Ufer} (Dativ lieu).
        \item \all{haben ... gesehen} (Phrase 8) : \textbf{haben} + \all{gesehen} (verbe fort \all{sehen}). Objet \all{den großen Umzug} (Accusatif). \all{mit den Masken} (Dativ pluriel \all{den} + \all{n} : \all{die Masken $\to$ den Masken}).
        \item \all{haben ... gegessen} (Phrase 9) : \textbf{haben} + \all{gegessen} (verbe fort \all{essen}). Objets \all{Ndolé} (invariable, masc.) et \all{gegrillten Fisch} (Accusatif adjectif \all{-en} sans article).
        \item \all{sind ... gegangen} (Phrase 10) : \textbf{sein} + \all{gegangen} (verbe fort \all{gehen}). \all{müde, aber glücklich} (adjectifs prédicatifs, pas de déclinaison). \all{nach Hause} (direction fixe, sans article).
    \end{itemize}
    \item \textbf{Orthographe / Ponctuation :} Majuscules à tous les noms (\all{Fest, Ufer, Wouri, Kleidung, Pirogen, Musik, Menge, Gerichte, Atmosphäre, Bus, Morgen, Umzug, Masken, Freunde, Ndolé, Fisch, Abends, Hause}). Umlaute présents dans le texte : \all{fröhlich, müde, glücklich, Kleidung, Freunde, Gewänder, Masken, Bräuche, Flüsse, Paläste, Abends}. \all{ß} après voyelle longue/diphtongue : \all{groß, heiß} (rappel, non présents dans ce texte). Virgule avant \all{und} seulement si proposition subordonnée (pas ici, coordination simple).
\end{itemize}

\textbf{Variantes acceptées / Différenciation :}
\begin{itemize}
    \item \textit{Niveau soutien :} Phrases plus courtes, vocabulaire donné (liste au tableau), modèle de phrase à trous : \all{Letztes Jahr sind wir nach ... gefahren. Wir haben ... gegessen.}
    \item \textit{Niveau approfondissement :} Ajout d'une relative (\all{Das Fest, das sehr alt ist, ...}), du \all{Präteritum} pour \all{sein/haben/modaux} (\all{war, hatte, konnte}), connecteurs logiques (\all{deshalb, danach, plötzlich}).
\end{itemize}
\end{exoresolu}

\courssub{Bilan de l'activité}
\begin{aretenir}[Bilan linguistique et communicatif]
Cette activité a permis de vérifier l'acquisition des trois piliers de la leçon ALA\_21 dans une tâche unique et signifiante :
\begin{enumerate}
    \item \textbf{Lexique :} Les élèves ont réinvesti \all{das Fest, die Tradition, der Tanz, das Essen, die Kleidung, der Brauch, der Umzug, feiern, aufführen, schmecken} avec genres et pluriels corrects. L'ancrage camerounais (Ngondo, Nguon, Ndolé, Wouri) a favorisé l'appropriation personnelle.
    \item \textbf{Grammaire – Dativ de lieu :} La contrainte « 3 prépositions de lieu différentes » a forcé la distinction \all{in/im} (espace clos), \all{an/am} (bord, rive, vertical), \all{auf/dem} (surface, place ouverte), \all{vor/hinter/neben} (position relative). Les erreurs résiduelles (\all{*in der Dorf}) sont corrigées collectivement lors de la relecture croisée.
    \item \textbf{Grammaire – Perfekt :} L'opposition \all{haben/sein} a été testée en situation réelle : verbes de mouvement (\all{fahren, gehen, kommen, ankommen}) $\to$ \all{sein} ; verbes d'action/transitifs (\all{essen, sehen, tanzen, feiern, kaufen}) $\to$ \all{haben}. Le positionnement du Participe II en fin de phrase a été respecté dans 90\% des productions (correction par le « Correcteur grammaire » du groupe).
    \item \textbf{Compétences transverses :} Expression orale continue (2 min), interaction (questions de la classe), écriture collaborative, culture générale (valorisation du patrimoine camerounais), citoyenneté (partage avec l'Autre germanophone).
\end{enumerate}
\end{aretenir}

\begin{methode}[Prolongements possibles (Devoirs / Projet)]
\begin{enumerate}
    \item \textbf{Version numérique :} Saisie des textes sur traitement de texte / Canva avec photos (libres de droits) $\to$ PDF „Feste in Kamerun“ envoyé au lycée partenaire de Leipzig.
    \item \textbf{Version audio/vidéo :} Enregistrement de la présentation (2 min) $\to$ QR code collé sur l'affiche au couloir.
    \item \textbf{Traduction inversée (Thème) :} Traduire la page modèle « Ngondo » en français (version) puis la retraduire en allemand sans regarder l'original (thème) pour fixer les structures.
    \item \textbf{Comparaison interculturelle :} Rechercher une fête allemande équivalente (ex. \all{Oktoberfest, Karneval, Weihnachtsmarkt, Schützenfest}) et remplir un tableau comparatif : Lieu / Date / Vêtements / Nourriture / Musique / Rôle social.
\end{enumerate}
\end{methode}

\newpage

\coursec{Méthodes et exercices résolus}

\courssub{Texte support et première lecture}

\begin{textebox}[Texte support : Le Ngondo, fête traditionnelle Sawa]
\all{Letztes Jahr habe ich mit meiner Familie das Ngondo-Fest in Douala gefeiert. Wir sind früh am Morgen zum Wasser gegangen. Viele Menschen haben traditionelle Kleidung getragen. Die Tänzer haben wunderbar getanzt, und alle haben zusammen gegessen. Am Abend sind wir müde, aber glücklich nach Hause gefahren.}
\end{textebox}

\begin{methode}[Lire et exploiter un texte de compréhension écrite]
\begin{enumerate}
\item \textbf{Première lecture globale.} Lis le texte en entier sans dictionnaire. Repère le thème (ici : une fête traditionnelle), le lieu (\all{Douala}, \all{zum Wasser}), les personnages (\all{meine Familie, die Tänzer, viele Menschen}), le moment (\all{Letztes Jahr} $\rightarrow$ passé). Souligne les mots du champ lexical : \all{das Fest, die Tradition, der Tanz, das Essen, die Kleidung, feiern}.
\item \textbf{Deuxième lecture détaillée.} Lis phrase par phrase. Devine les mots inconnus grâce au contexte et aux familles de mots (\all{feiern} $\rightarrow$ \all{das Fest, die Feier, der Feiertag} ; \all{tanzen} $\rightarrow$ \all{der Tanz, der Tänzer}).
\item \textbf{Repérage des indices grammaticaux.} Le \all{Perfekt} (\all{habe gefeiert, sind gegangen, haben getragen, haben getanzt, haben gegessen, sind gefahren}) signale un récit au passé ; le datif après \all{zu, mit, in} (\all{zum Wasser, mit meiner Familie, in Douala}) indique le lieu ou l'accompagnement.
\item \textbf{Réponse aux questions.} Réponds toujours par une phrase complète en allemand, en reprenant les mots de la question, jamais par un simple « oui » ou « non ». Vérifie la place du verbe en deuxième position.
\item \textbf{Résumé.} Pour résumer, garde uniquement : qui ? quoi ? où ? quand ? Rédige 3 à 5 phrases au présent ou au Perfekt, avec tes propres mots.
\end{enumerate}
\end{methode}

\courssub{Compréhension du texte}

\begin{exercice}[Questions de compréhension détaillée]
Réponds en allemand par des phrases complètes.
\begin{enumerate}
\item \all{Wo hat die Familie das Fest gefeiert?}
\item \all{Wann ist die Familie zum Wasser gegangen?}
\item \all{Was haben die Menschen getragen?}
\item \all{Was haben die Tänzer gemacht?}
\item \all{Was haben alle zusammen gemacht?}
\item \all{Wie sind sie am Abend nach Hause gefahren?}
\item \all{Wie war die Stimmung am Ende?}
\end{enumerate}
\end{exercice}

\begin{corrige}[Exercice 1 : Compréhension écrite]
\begin{enumerate}
\item \all{Die Familie hat das Fest in Douala gefeiert.}
\item \all{Die Familie ist früh am Morgen zum Wasser gegangen.}
\item \all{Die Menschen haben traditionelle Kleidung getragen.}
\item \all{Die Tänzer haben wunderbar getanzt.}
\item \all{Alle haben zusammen gegessen.}
\item \all{Sie sind müde, aber glücklich nach Hause gefahren.}
\item \all{Die Stimmung war wunderbar / Alle waren glücklich.}
\end{enumerate}
\end{corrige}

\begin{exoresolu}[Modèle de résumé en 5 phrases]
\textbf{Consigne :} Résume le texte en 5 phrases (Perfekt ou Présent).
\tcblower
\textbf{Proposition (Perfekt) :}
\all{Letztes Jahr hat der Erzähler das Ngondo-Fest in Douala mit seiner Familie gefeiert. Zuerst sind sie zum Wasser gegangen. Die Menschen haben traditionelle Kleidung getragen und die Tänzer haben getanzt. Danach haben alle zusammen gegessen. Am Ende sind sie glücklich nach Hause gefahren.}

\textbf{Proposition (Présent - description habituelle) :}
\all{Das Ngondo-Fest ist eine traditionelle Feier in Douala. Die Menschen gehen zum Wasser und tragen schöne Kleidung. Es gibt Tänze und ein gemeinsames Essen. Am Abend fahren alle glücklich nach Hause. Diese Tradition ist sehr wichtig.}
\end{exoresolu}

\courssub{Lexique thématique : les fêtes traditionnelles}

\begin{definition}[Vocabulaire de base : noms, articles, pluriels]
\begin{center}
\begin{tabularx}{\linewidth}{|X|X|X|}
\hline
\rowcolor{popblueL} \textbf{Allemand (Nom + Article + Pluriel)} & \textbf{Français} & \textbf{Famille de mots / Exemple} \\
\hline
\all{das Fest, die Feste} & la fête, la festival & \all{feiern (v.), die Feier, der Feiertag} \\
\all{die Tradition, die Traditionen} & la tradition & \all{traditionell (adj.), traditionell feiern} \\
\all{der Tanz, die Tänze} & la danse & \all{tanzen (v.), der Tänzer / die Tänzerin} \\
\all{das Essen, -} & le repas, la nourriture & \all{essen (v., unregelmäßig), das Mittagessen} \\
\all{die Kleidung (kein Plural / nur Singular)} & les vêtements, l'habillement & \all{tragen (v., unregelmäßig), kleiden} \\
\all{die Feier, die Feiern} & la célébration, la fête (événement) & \all{feiern, das Fest} \\
\all{der Feiertag, die Feiertage} & le jour férié & \all{feiern, der Tag} \\
\all{der Brauch, die Bräuche} & la coutume, l'usage & \all{brauchen (v. \emph{avoir besoin}), üblich (adj.)} \\
\all{das Volk, die Völker} & le peuple & \all{völkisch (adj.), die Volksmusik} \\
\all{der Tanz, die Tänze} & la danse & \all{tanzen, der Tanzstil} \\
\hline
\end{tabularx}
\end{center}
\emph{Note :} \all{die Kleidung} est un collectif, généralement au singulier. Pour « différents vêtements », on utilise \all{die Kleidungsstücke}.
\end{definition}

\begin{methode}[Mémoriser et utiliser le lexique thématique]
\begin{enumerate}
\item \textbf{Apprends chaque nom avec son article et son pluriel} (voir tableau ci-dessus).
\item \textbf{Regroupe par familles de mots} : \all{feiern} (verbe) $\rightarrow$ \all{das Fest, die Feier, der Feiertag, feierlich} ; \all{tanzen} $\rightarrow$ \all{der Tanz, der Tänzer, die Tänzerin, tänzerisch} ; \all{essen} $\rightarrow$ \all{das Essen, der Esser, essbar}.
\item \textbf{Construis des phrases modèles} réutilisables : \all{Wir feiern das Fest mit der Familie.} ; \all{Die Menschen tragen traditionelle Kleidung.} ; \all{Man tanzt und isst zusammen.} ; \all{Der Brauch ist sehr alt.}
\item \textbf{Attention aux genres} : un nom mal genré (\all{\textbf{die} Fest} au lieu de \all{\textbf{das} Fest}) entraîne des fautes de déclinaison en cascade (article, adjectif, pronom).
\item \textbf{Réutilise le lexique dans ta production} : une bonne copie emploie au moins cinq mots du thème, correctement orthographiés avec la majuscule aux noms.
\end{enumerate}
\end{methode}

\begin{exoresolu}[Vocabulaire : compléter un texte à trous]
\textbf{Énoncé :} Complète avec le mot qui convient (forme correcte) : \all{Fest -- Tradition -- Tanz -- Essen -- Kleidung -- feiern -- Brauch}.

\all{In Kamerun gibt es viele traditionelle (1) \dots . Beim Ngondo (2) \dots\ die Sawa am Wasser. Die Menschen tragen schöne (3) \dots\ und der (4) \dots\ ist sehr wichtig. Nach dem (4) \dots\ gibt es ein großes (5) \dots\ für alle. Dieser (6) \dots\ ist sehr alt. Man (7) \dots\ gerne zusammen.}
\tcblower
\textbf{Raisonnement :} On analyse chaque trou : article/indicateur, genre, nombre, fonction (sujet/objet), sens du contexte.

\textbf{Solutions :}
\begin{enumerate}
\item \all{Feste} : après \all{viele} (pluriel), il faut le pluriel de \all{das Fest} $\rightarrow$ \all{Feste}.
\item \all{feiern} : sujet pluriel \all{die Sawa}, verbe \all{feiern} au présent, 3\textsuperscript{e} personne du pluriel $\rightarrow$ \all{feiern}.
\item \all{Kleidung} : \all{schöne} (accusatif féminin singulier) + nom féminin ; on parle de ce que l'on porte (\all{tragen}).
\item \all{Tanz} : article \all{der} (nominatif masculin) $\rightarrow$ nom masculin ; le contexte (\all{Nach dem Tanz}) confirme (datif \all{dem Tanz}).
\item \all{Essen} : \all{ein großes} (accusatif neutre singulier) $\rightarrow$ nom neutre ; après la danse, on mange.
\item \all{Brauch} (ou \all{Tradition}) : \all{Dieser} (nominatif masculin) $\rightarrow$ nom masculin ; conclusion du texte sur la coutume.
\item \all{feiern} : sujet \all{Man} (3\textsuperscript{e} pers. sg) $\rightarrow$ \all{feiert}. \emph{Attention : conjugaison !} $\rightarrow$ \all{feiert}.
\end{enumerate}

\textbf{Texte complet :}
\all{In Kamerun gibt es viele traditionelle Feste. Beim Ngondo feiern die Sawa am Wasser. Die Menschen tragen schöne Kleidung und der Tanz ist sehr wichtig. Nach dem Tanz gibt es ein großes Essen für alle. Dieser Brauch ist sehr alt. Man feiert gerne zusammen.}

\textbf{Conclusion :} Vérifier l'article devant le trou (ou la terminaison de l'adjectif) et le sens de la phrase permet de choisir le bon mot et la bonne forme sans hésitation.
\end{exoresolu}

\courssub{Grammaire 1 : le datif et les prépositions de lieu}

\begin{propriete}[Le cas datif : formes des articles et adjectifs]
\textbf{Tableau de déclinaison (Datif singulier \& pluriel) :}
\begin{center}
\begin{tabular}{|l|c|c|c|c|}
\hline
\rowcolor{popblueL} & \textbf{Masculin} & \textbf{Féminin} & \textbf{Neutre} & \textbf{Pluriel} \\
\hline
\textbf{Article défini} & \all{dem} & \all{der} & \all{dem} & \all{den} \\
\hline
\textbf{Article indéfini} & \all{einem} & \all{einer} & \all{einem} & \all{---} \\
\hline
\textbf{Adj. (décl. faible)} & \all{dem guten Mann} & \all{der guten Frau} & \all{dem guten Kind} & \all{den guten Leuten} \\
\hline
\textbf{Adj. (décl. forte/sans art.)} & \all{gutem Mann} & \all{guter Frau} & \all{gutem Kind} & \all{guten Leuten} \\
\hline
\textbf{Nom en \cle{-n} (pluriel)} & --- & --- & --- & \all{den Kindern, den Leuten, den Frauen} \\
\hline
\end{tabular}
\end{center}
\cle{Règle d'or :} Au datif pluriel, l'article est toujours \all{den} et le nom prend un \cle{-n} final (sauf s'il finit déjà par \cle{-n} ou \cle{-s} : \all{den Frauen, den Autos}).
\end{propriete}

\begin{propriete}[Prépositions et cas : Datif vs Accusatif]
\begin{center}
\begin{tabularx}{\linewidth}{|X|X|X|}
\hline
\rowcolor{popblueL} \textbf{Type} & \textbf{Prépositions} & \textbf{Cas / Question} \\
\hline
\textbf{Toujours Datif} & \all{aus, bei, mit, nach, von, zu, seit, gegenüber} & \all{Wo?} (Position) / \all{Wem?} \\
\hline
\textbf{Toujours Accusatif} & \all{durch, für, gegen, ohne, um, entlang} & \all{Wohin?} (Direction) / \all{Wen?} \\
\hline
\textbf{Deux cas (Wechselpräpositionen)} & \all{in, an, auf, neben, vor, hinter, über, unter, zwischen} & \all{Wo?} $\rightarrow$ \textbf{Datif} (Position statique) \\
& & \all{Wohin?} $\rightarrow$ \textbf{Accusatif} (Mouvement vers) \\
\hline
\end{tabularx}
\end{center}
\end{propriete}

\begin{methode}[Employer le datif avec les prépositions de lieu]
\begin{enumerate}
\item \textbf{Apprends par cœur les prépositions toujours suivies du datif} : \all{aus, bei, mit, nach, von, zu, seit, gegenüber}.
\item \textbf{Prépositions de lieu à deux cas (Wechselpräpositionen)} : \all{in, an, auf, neben, vor, hinter, über, unter, zwischen}. Question \all{Wo?} (où ? position statique) $\rightarrow$ \textbf{datif} ; question \all{Wohin?} (où va-t-on ? mouvement vers) $\rightarrow$ \textbf{accusatif}.
\item \textbf{Connais les formes du datif} : masc. \all{dem/einem}, fém. \all{der/einer}, neutre \all{dem/einem}, pluriel \all{den} + nom en \cle{-n} (\all{mit den Kindern}).
\item \textbf{Contractions utiles (obligatoires à l'oral, fréquentes à l'écrit)} :
    \begin{itemize}
    \item \all{zu dem = zum} (\all{zum Markt, zum Fest})
    \item \all{bei dem = beim} (\all{beim Ngondo})
    \item \all{an dem = am} (\all{am Platz})
    \item \all{in dem = im} (\all{im Dorf, im Wasser})
    \item \all{zu der = zur} (\all{zur Feier})
    \item \all{in das = ins} (\all{ins Dorf} -- mouvement !)
    \end{itemize}
\item \textbf{Vérifie toujours} : après avoir écrit la préposition, demande-toi « mouvement (\all{Wohin?}) ou position (\all{Wo?}) ? » avant de choisir le cas.
\end{enumerate}
\end{methode}

\begin{exoresolu}[Grammaire : choisir le bon cas après une préposition de lieu]
\textbf{Énoncé :} Complète avec l'article au cas correct (datif ou accusatif). Attention aux contractions.

\begin{enumerate}
\item \all{Wir feiern das Fest in \dots\ Dorf. (das, position)}
\item \all{Die Familie fährt in \dots\ Stadt Douala. (die, mouvement)}
\item \all{Die Tänzer stehen auf \dots\ Platz. (der, position)}
\item \all{Die Kinder gehen zu \dots\ Großeltern. (die, Plural, mouvement vers personne)}
\item \all{Wir sprechen mit \dots\ Leuten. (die, Plural, accompagnement)}
\item \all{Das Essen steht auf \dots\ Tisch. (der, position)}
\end{enumerate}
\tcblower
\textbf{Raisonnement :} Pour chaque phrase, on pose la question : \all{Wo?} (position, datif) ou \all{Wohin?} (mouvement, accusatif). \all{zu, mit} sont toujours datif.

\textbf{Solutions :}
\begin{enumerate}
\item \all{in dem (im) Dorf} — \all{feiern} + lieu $\rightarrow$ \all{Wo?} Position $\rightarrow$ datif : \all{dem Dorf} $\rightarrow$ \all{im Dorf}.
\item \all{in die Stadt Douala} — \all{fahren} exprime un mouvement vers un lieu $\rightarrow$ \all{Wohin?} Accusatif : \all{die Stadt} (fém. acc. = nom.).
\item \all{auf dem Platz} — \all{stehen} indique une position statique $\rightarrow$ \all{Wo?} Datif : \all{dem Platz}.
\item \all{zu den Großeltern} — \all{zu} + personne $\rightarrow$ toujours datif ; pluriel : \all{den} + \cle{-n} $\rightarrow$ \all{den Großeltern}.
\item \all{mit den Leuten} — \all{mit} $\rightarrow$ toujours datif ; pluriel : \all{den Leuten} (\all{Leute} prend \cle{-n} au datif pluriel).
\item \all{auf dem Tisch} — \all{stehen} + \all{Wo?} $\rightarrow$ datif : \all{dem Tisch}.
\end{enumerate}

\textbf{Conclusion :} La distinction \all{Wo?/Wohin?} est le réflexe essentiel : position = datif, mouvement = accusatif. Et \all{zu, mit, bei, nach, von, aus, seit, gegenüber} sont toujours au datif, sans exception.
\end{exoresolu}

\courssub{Grammaire 2 : le passé composé (das Perfekt)}

\begin{propriete}[Structure et formation du Perfekt]
\textbf{Structure :} Sujet + Auxiliaire conjugué (2\textsuperscript{e} position) + \dots + Participe Passé (à la fin de la phrase).
\medskip

\textbf{Choix de l'auxiliaire :}
\begin{itemize}
\item \all{sein} : verbes de mouvement (\all{gehen, kommen, fahren, reisen, laufen, fliegen, schwimmen}) + verbes de changement d'état (\all{aufstehen, einschlafen, aufwachen, sterben, werden, bleiben} -- \emph{rester/devenir}). \emph{Astuce :} « Aller quelque part » ou « Changer d'état ».
\item \all{haben} : tous les autres verbes (transitifs, réflexifs, modaux au Perfekt \emph{Ersatzinfinitiv}, verbes d'état comme \all{sein, haben, wohnen}).
\end{itemize}

\textbf{Formation du Participe Passé (Partizip II) :}
\begin{center}
\begin{tabular}{|l|l|l|l|}
\hline
\rowcolor{popblueL} \textbf{Type} & \textbf{Règle} & \textbf{Exemple (Infinitif)} & \textbf{Partizip II} \\
\hline
\all{Verbes faibles (réguliers)} & \all{ge-} + radical + \all{-t} & \all{feiern, tanzen, kaufen, machen} & \all{gefeiert, getanzt, gekauft, gemacht} \\
\hline
\all{Verbes forts (irréguliers)} & \all{ge-} + radical modifié + \all{-en} & \all{essen, tragen, geben, singen, trinken} & \all{gegessen, getragen, gegeben, gesungen, getrunken} \\
\hline
\all{Verbes mixtes} & \all{ge-} + radical modifié + \all{-t} & \all{denken, bringen, kennen, brennen} & \all{gedacht, gebracht, gekannt, gebrannt} \\
\hline
\all{Verbes à particule séparable} & Particule + \all{ge-} + radical + \all{-t/en} & \all{ankommen, aufstehen, mitbringen} & \all{angekommen, aufgestanden, mitgebracht} \\
\hline
\all{Verbes en \cle{-ieren} / préfixes inséparables} & Sans \all{ge-} & \all{studieren, besuchen, verlieren, verstehen} & \all{studiert, besucht, verloren, verstanden} \\
\hline
\end{tabular}
\end{center}
\end{propriete}

\begin{methode}[Construire et employer le Perfekt pour raconter une fête]
\begin{enumerate}
\item \textbf{Identifie le verbe} et choisis l'auxiliaire (\all{sein} = mouvement/changement d'état, \all{haben} = reste).
\item \textbf{Forme le participe passé} selon le type (faible \cle{-t}, fort \cle{-en}, séparable \cle{ge-} au milieu).
\item \textbf{Place l'auxiliaire conjugué en 2\textsuperscript{e} position} et le participe à la \textbf{fin de la phrase} (ou de la proposition).
\item \textbf{Plan de récit d'une fête} (connecteurs temporels) : \all{Letztes Jahr / Letztes Wochenende} (date) $\rightarrow$ \all{zuerst / zuerst einmal} (préparatifs) $\rightarrow$ \all{dann / danach} (déroulement : danses, repas) $\rightarrow$ \all{am Ende / zum Schluss} (fin) $\rightarrow$ \all{weil / denn} (impression/raison).
\end{enumerate}
\end{methode}

\begin{exoresolu}[Conjugaison : mettre un texte au Perfekt]
\textbf{Énoncé :} Récris les phrases au Perfekt.

\begin{enumerate}
\item \all{Wir feiern das Fest im Dorf.}
\item \all{Die Gäste kommen am Abend an.}
\item \all{Die Frauen tragen traditionelle Kleidung.}
\item \all{Alle essen zusammen.}
\item \all{Ich trinke ein Glas Wasser.}
\item \all{Wir schlafen im Zelt.}
\end{enumerate}
\tcblower
\textbf{Raisonnement :} Pour chaque verbe : choix de l'auxiliaire, formation du participe, place du participe à la fin.

\textbf{Solutions :}
\begin{enumerate}
\item \all{Wir haben das Fest im Dorf gefeiert.} — \all{feiern} (faible, avoir) $\rightarrow$ \all{haben gefeiert}.
\item \all{Die Gäste sind am Abend angekommen.} — \all{ankommen} (séparable, mouvement, être) $\rightarrow$ \all{sind angekommen} (\all{an+ge+kommen}).
\item \all{Die Frauen haben traditionelle Kleidung getragen.} — \all{tragen} (fort, avoir) $\rightarrow$ \all{haben getragen}.
\item \all{Alle haben zusammen gegessen.} — \all{essen} (fort, avoir) $\rightarrow$ \all{haben gegessen} (doublement \cle{e}).
\item \all{Ich habe ein Glas Wasser getrunken.} — \all{trinken} (fort, avoir) $\rightarrow$ \all{habe getrunken}.
\item \all{Wir haben im Zelt geschlafen.} — \all{schlafen} (fort, avoir -- pas de changement de lieu \emph{vers}, action statique) $\rightarrow$ \all{haben geschlafen}. \emph{Note :} \all{schlafen} prend \all{haben} (sauf sens « s'endormir » \all{einschlafen} $\rightarrow$ \all{sein}).
\end{enumerate}

\textbf{Conclusion :} Deux pièges classiques sont évités : l'auxiliaire \all{sein} avec les verbes de mouvement (\all{gehen, kommen, fahren}) et le \all{-ge-} inséré dans les verbes séparables (\all{ankommen $\rightarrow$ angekommen}). Le participe se place toujours en fin de phrase.
\end{exoresolu}

\courssub{Méthode du commentaire et commentaire modèle}

\begin{methode}[Commenter un texte sur une fête traditionnelle (Écrit / Oral)]
\begin{enumerate}
\item \textbf{Introduction (1 phrase) :} Type de texte, thème, source/contexte. \all{Der Text „Das Ngondo-Fest“ ist ein Erlebnisbericht. Er handelt von einer traditionellen Feier in Douala.}
\item \textbf{Résumé structuré (3--4 phrases) :} Qui ? Quoi ? Où ? Quand ? Déroulement. Utilise \all{zuerst, dann, danach, am Ende}.
\item \textbf{Analyse lexicale (2--3 phrases) :} Relevé du champ lexical (fête, tradition, danse, vêtements, repas). Cite des exemples du texte. \all{Der Wortschatz gehört zum Themenkreis „Feste“: das Fest, die Kleidung, der Tanz, das Essen, feiern.}
\item \textbf{Analyse grammaticale (2 phrases) :} Repérage du temps (Perfekt = récit passé), des prépositions de lieu (datif/accusatif), de la structure verbe 2\textsuperscript{e} position. \all{Der Text steht im Perfekt (Vergangenheit). Man sieht Wechselpräpositionen: in Douala (Dativ), zum Wasser (Dativ), nach Hause (ohne Artikel).}
\item \textbf{Appréciation personnelle / Lien culturel (1--2 phrases) :} Comparaison avec une fête connue (Noël, Tabaski, Ngondo, Oktoberfest). \all{Das Fest erinnert an das Oktoberfest in Deutschland: Gemeinschaft, Tanz, Essen, Tradition.}
\end{enumerate}
\end{methode}

\begin{exoresolu}[Commentaire modèle intégré (niveau A2/B1)]
\textbf{Sujet :} Commente le texte du Ngondo en suivant la méthode.
\tcblower
\textbf{Commentaire :}
\all{Der Text „Das Ngondo-Fest“ ist ein persönlicher Erlebnisbericht aus der Ich-Perspektive. Er beschreibt die Feier des traditionellen Ngondo-Festes der Sawa in Douala.}
\all{Inhaltlich: Letztes Jahr ist der Erzähler mit seiner Familie zum Wasser gegangen. Viele Menschen trugen traditionelle Kleidung. Die Tänzer tanzten wunderbar. Alle aßen zusammen. Am Abend fuhren sie glücklich nach Hause.}
\all{Sprachlich: Der Text nutzt das Perfekt für die Vergangenheit (habe gefeiert, sind gegangen, haben getragen, haben getanzt, haben gegessen, sind gefahren). Wichtige Präpositionen mit Dativ: mit meiner Familie, zum Wasser (zu dem), in Douala. Der Wortschatz ist thematisch passend: das Fest, die Tradition, die Kleidung, der Tanz, das Essen, feiern.}
\all{Das Fest ähnelt deutschen Volksfesten wie dem Oktoberfest: Gemeinschaft, Tanz, Essen und Tradition stehen im Mittelpunkt. Es zeigt, wie wichtig kulturelle Identität ist.}
\end{exoresolu}

\courssub{Traduction (Thème) : du français vers l'allemand}

\begin{methode}[Traduire du français vers l'allemand sans calque]
\begin{enumerate}
\item \textbf{Analyse la phrase française} : sujet, verbe, temps, compléments. Repère les temps passés $\rightarrow$ Perfekt en allemand (langue parlée/écrite courante).
\item \textbf{Place le verbe conjugué en deuxième position}, quoi qu'il y ait au début de la phrase (\all{Letztes Jahr \textbf{haben} wir...} et non \all{Letztes Jahr wir haben...}).
\item \textbf{Choisis le bon cas} après chaque préposition (\all{mit der Familie} [Datif], \all{in das/ins Dorf} [Akk. mouvement], \all{im Dorf} [Dat. position], \all{zum Fest} [Datif]).
\item \textbf{Méfie-toi des faux amis} :
    \begin{itemize}
    \item « Assister à » $\neq$ \all{assistieren} $\rightarrow$ \all{teilnehmen an} (+ Dat) ou \all{dabei sein} ou \all{besuchen}.
    \item « Porter (des vêtements) » $\neq$ \all{portieren} $\rightarrow$ \all{tragen} (irrégulier : \all{trägt, trug, getragen}).
    \item « Une fête » : \all{das Fest} (événement), \all{die Feier} (célébration), \all{der Feiertag} (jour férié).
    \item « Traditionnel » $\rightarrow$ \all{traditionell} (pas \all{traditionnel}).
    \end{itemize}
\item \textbf{Relis} : majuscule à \textbf{tous} les noms, Umlaute (\all{ä, ö, ü}), \all{ß} après voyelle longue/diphtongue, participe en fin de phrase, accord de l'auxiliaire avec le sujet.
\end{enumerate}
\end{methode}

\begin{exoresolu}[Thème : traduction de phrases sur une fête traditionnelle]
\textbf{Énoncé :} Traduis en allemand.

\begin{enumerate}
\item L'année dernière, nous avons célébré la fête avec nos grands-parents.
\item Beaucoup de gens sont venus au village.
\item Les femmes ont porté des vêtements traditionnels et les hommes ont dansé.
\item Après la danse, nous avons mangé un bon repas ensemble.
\item Nous sommes rentrés chez nous tard le soir.
\end{enumerate}
\tcblower
\textbf{Raisonnement :} Chaque phrase est au passé $\rightarrow$ Perfekt. On identifie le sujet, l'auxiliaire, le participe et les cas (prépositions).

\textbf{Solutions détaillées :}
\begin{enumerate}
\item \all{Letztes Jahr haben wir das Fest mit unseren Großeltern gefeiert.}
    \begin{itemize}
    \item \all{Letztes Jahr} (complément de temps en tête) $\rightarrow$ inversion : verbe \all{haben} (2\textsuperscript{e} pos.), sujet \all{wir}.
    \item \all{mit} + Datif pluriel : \all{mit unseren Großeltern} (article \all{unseren}, nom \all{Großeltern} déjà en \cle{-n}).
    \item \all{gefeiert} (faible) en fin de phrase.
    \end{itemize}
\item \all{Viele Menschen sind ins Dorf gekommen.}
    \begin{itemize}
    \item \all{kommen} = mouvement $\rightarrow$ auxiliaire \all{sein} (\all{sind}).
    \item « Au village » = direction \all{Wohin?} $\rightarrow$ Accusatif : \all{in das Dorf = ins Dorf}.
    \item \all{gekommen} (fort, \all{ge-} + \all{kommen}) en fin.
    \end{itemize}
\item \all{Die Frauen haben traditionelle Kleidung getragen und die Männer haben getanzt.}
    \begin{itemize}
    \item \all{tragen} (fort) $\rightarrow$ \all{haben getragen}.
    \item \all{Kleidung} = singulier collectif (accusatif fém. : \all{die/eine traditionelle Kleidung}).
    \item \all{tanzen} (faible) $\rightarrow$ \all{haben getanzt}.
    \end{itemize}
\item \all{Nach dem Tanz haben wir zusammen ein gutes Essen gegessen.} (ou \all{gegessen} + \all{gegessen} répétition $\rightarrow$ mieux : \all{zusammen gut gegessen} ou \all{ein gutes Essen gegessen}).
    \begin{itemize}
    \item \all{Nach} + Datif : \all{dem Tanz}.
    \item \all{zuerst/ensuite} $\rightarrow$ \all{dann/nachdem} (mais \all{nach dem} + nom ici).
    \item \all{essen} $\rightarrow$ \all{haben gegessen}.
    \end{itemize}
\item \all{Wir sind spät am Abend nach Hause gekommen / gefahren.}
    \begin{itemize}
    \item « Rentrer » = \all{nach Hause kommen/fahren/gehen} (mouvement $\rightarrow$ \all{sein}).
    \item \all{nach Hause} (sans article, direction fixe).
    \item \all{spät am Abend} (temps).
    \end{itemize}
\end{enumerate}

\textbf{Conclusion :} Les trois points de vigilance du thème : inversion après complément initial (V2), choix de l'auxiliaire du Perfekt (\all{sein}/\all{haben}), et cas après prépositions (\all{mit} + Dat, \all{in} + Akk/Dat, \all{nach} + Dat).
\end{exoresolu}

\courssub{Production guidée : ma fête traditionnelle}

\begin{methode}[Rédiger un paragraphe guidé sur une fête traditionnelle]
\begin{enumerate}
\item \textbf{Prépare un plan en 4 points} : (a) Nom et lieu de la fête ; (b) Avec qui et quand ; (c) Déroulement (vêtements, danses, repas, chants) ; (d) Ton impression / sentiment final.
\item \textbf{Choisis un temps et tiens-toi-y} : Récit d'une fête passée $\rightarrow$ \textbf{Perfekt} ; Description d'une fête habituelle/récurrente $\rightarrow$ \textbf{Présent}.
\item \textbf{Utilise des connecteurs simples} : \all{zuerst, dann, danach, am Ende, und, aber, weil} (\textbf{Attention :} après \all{weil}, le verbe va à la fin de la proposition).
\item \textbf{Insère au moins 5 mots du lexique} : \all{das Fest, die Tradition, der Tanz, das Essen, die Kleidung, feiern, der Brauch, die Feier}.
\item \textbf{Relis avec la grille de contrôle} :
    \begin{itemize}
    \item Verbe conjugué en 2\textsuperscript{e} position ?
    \item Participe passé à la fin (si Perfekt) ?
    \item Majuscules à \textbf{tous} les noms ?
    \item Cas correct après prépositions (Datif / Accusatif) ?
    \item Auxiliaire correct (\all{sein} / \all{haben}) ?
    \item Umlaute (\all{ä, ö, ü}) et \all{ß} corrects ?
    \end{itemize}
\end{enumerate}
\end{methode}

\begin{exoresolu}[Expression écrite : rédiger un court paragraphe (Modèle corrigé)]
\textbf{Énoncé :} Décris en 6 à 8 phrases une fête traditionnelle que tu as célébrée avec ta famille. Utilise le Perfekt, au moins cinq mots du vocabulaire du thème et deux prépositions de lieu (dont une \all{Wechselpräposition}).
\tcblower
\textbf{Raisonnement :} On suit le plan en 4 points : présentation de la fête, accompagnement, déroulement, impression. On vérifie chaque phrase : auxiliaire correct, participe final, cas après prépositions, majuscules.

\textbf{Production modèle :}
\all{Letztes Jahr habe ich mit meiner Familie das Ngondo-Fest in Douala gefeiert. Zuerst sind wir früh am Morgen zum Wasser gegangen. Viele Menschen haben dort traditionelle Kleidung getragen. Dann haben die Tänzer auf dem Platz wunderbar getanzt. Danach haben alle zusammen ein großes Essen gegessen und gesungen. Am Ende bin ich müde, aber sehr glücklich nach Hause gefahren. Diese Tradition ist wunderschön und ich werde sie nie vergessen.}

\textbf{Vérification de la grille :}
\begin{itemize}
\item \textbf{Lexique du thème (6) :} \all{Fest, feiern, Kleidung, Tänzer, getanzt, Essen, Tradition} (\all{Tradition} dans dernière phrase).
\item \textbf{Perfekt correct :}
    \begin{itemize}
    \item \all{habe ... gefeiert} (faible, avoir)
    \item \all{sind ... gegangen} (mouvement, être)
    \item \all{haben ... getragen} (fort, avoir)
    \item \all{haben ... getanzt} (faible, avoir)
    \item \all{haben ... gegessen} (fort, avoir) / \all{haben ... gesungen} (fort, avoir)
    \item \all{bin ... gefahren} (mouvement, être)
    \end{itemize}
\item \textbf{Prépositions de lieu (3) :}
    \begin{itemize}
    \item \all{in Douala} (Wo? Position $\rightarrow$ Datif, nom propre sans article).
    \item \all{zum Wasser} = \all{zu dem Wasser} (\all{zu} + Datif, mouvement vers personne/lieu $\rightarrow$ Datif fixe).
    \item \all{auf dem Platz} (\all{auf} + Datif, Wo? Position statique).
    \item \all{nach Hause} (direction fixe, sans article, Datif historique).
    \end{itemize}
\item \textbf{Subordonnée :} \all{weil ... vergessen} (verbe à la fin). \all{und} (coordination, verbe 2\textsuperscript{e} pos).
\item \textbf{Orthographe :} Majuscules aux noms (\all{Fest, Wasser, Menschen, Kleidung, Tänzer, Platz, Essen, Ende, Tradition}), Umlaute (\all{schön, müde, glücklich, wunderbar, Grö...} -- ici \all{größten} absent mais \all{wunderschön}), \all{ß} (\all{groß} absent, \all{Straße} absent -- vérifier \all{gefreut}? non. \all{weiß}? non. \all{ß} pas dans ce texte, correct).
\end{itemize}

\textbf{Conclusion :} Un paragraphe court mais complet : plan respecté, grammaire contrôlée (Perfekt, cas, V2, subordonnée), vocabulaire riche et varié. C'est exactement ce que l'examinateur attend au Bac.
\end{exoresolu}

\begin{experience}[Activité orale : Interview « Meine Lieblingsfeste »]
\textbf{Consigne :} En binôme. L'élève A est journaliste, l'élève B est invité(e).
\textbf{Questions types (Élève A) :}
\begin{enumerate}
\item \all{Welches traditionelle Fest feierst du gerne?}
\item \all{Wann und wo findet es statt?}
\item \all{Mit wem feierst du?}
\item \all{Was isst und trinkt man dort?}
\item \all{Was trägt man? (Kleidung)}
\item \all{Gibt es Tanz oder Musik?}
\item \all{Warum ist dieses Fest wichtig für dich?}
\end{enumerate}
\textbf{Réponses types (Élève B) :} Utiliser le Présent (habitude) ou Perfekt (souvenir précis). \all{Ich feiere gerne Weihnachten / das Ngondo / das Erntedankfest. ... Wir essen ... Wir singen ... Es ist wichtig, weil ...}
\textbf{Objectif :} Interaction spontanée, réemploi lexique, Perfekt/Présent, prépositions de lieu.
\end{experience}

\begin{savaistu}[Culture : Fêtes au Cameroun et dans l'espace germanophone]
\textbf{Cameroun :} Le \all{Ngondo} (peuple Sawa, Douala, novembre/décembre) : culte des ancêtres, pirogues, danse, immersion dans le Wouri. Le \all{Nguon} (Bamoun, Foumban, tous les 2 ans) : grande cérémonie royale, masques, danses, marché. \all{Tabaski} (Aïd el-Kébir) : fête musulmane du sacrifice, partage mouton, famille. \all{Noël} / \all{Fin d'année} : fêtes chrétiennes et civiles, repas de famille, \all{le réveillon}.

\textbf{Allemagne / Autriche / Suisse (DACH) :}
\begin{itemize}
\item \all{Weihnachten} (Noël) : 24/25 déc., \all{Weihnachtsmarkt}, \all{Christbaum}, \all{Plätzchen}, \all{Bescherung}.
\item \all{Silvester} (Saint-Sylvestre / 31 déc.) : feux d'artifice (\all{Feuerwerk}), \all{Bleigießen} (plomb fondu), \all{Berliner/Pfannkuchen} (beignets).
\item \all{Karneval / Fasching / Fastnacht} (Carnaval) : fév., costumes (\all{Kostüme}), défilés (\all{Umzüge}), \all{Alaaf!} (Cologne) / \all{Helau!} (Düsseldorf/Mayence).
\item \all{Oktoberfest} (Munich, sept.-oct.) : plus grande fête foraine mondiale, \all{Dirndl} / \all{Lederhosen}, \all{Maß Bier}, \all{Brezn}.
\item \all{Erntedankfest} (Fête des moissons) : sept./oct., église décorée fruits/légumes, procession.
\item \all{Tag der Deutschen Einheit} (3 oct.) : fête nationale (réunification 1990).
\end{itemize}
\textbf{Point commun :} Communauté (\all{Gemeinschaft}), tradition (\all{Tradition}), nourriture (\all{Essen}), musique/danse (\all{Musik/Tanz}), vêtements spéciaux (\all{Tracht/Kleidung}).
\end{savaistu}

\begin{attention}[Les 5 erreurs qui coûtent des points au Bac]
\begin{enumerate}
\item \textbf{Oublier la majuscule aux noms.} \all{das fest, die tradition, der tanz} sont \textbf{fautifs} : écris \all{das Fest, die Tradition, der Tanz}. En allemand, \emph{tous} les noms (substantifs) prennent la majuscule, partout dans la phrase (sujet, objet, après préposition).
\item \textbf{Mauvais auxiliaire au Perfekt.} Ne dis jamais \all{wir haben gegangen} ou \all{ich habe gefahren} : les verbes de mouvement (\all{gehen, kommen, fahren, laufen, reisen}) et changement d'état exigent \all{sein} : \all{wir \textbf{sind} gegangen, sie \textbf{ist} gekommen, ich \textbf{bin} gefahren}.
\item \textbf{Verbe mal placé (Règle V2 et subordonnée).} Après un complément en tête de phrase, le verbe conjugué reste en \textbf{deuxième position} et le sujet suit : \all{Letztes Jahr \textbf{haben wir} gefeiert} (et non \all{Letztes Jahr wir haben gefeiert}). Dans une subordonnée (\all{weil, dass, wenn...}), le verbe conjugué va à la \textbf{fin} : \all{weil das Fest wunderbar \textbf{war}}.
\item \textbf{Confondre datif et accusatif après les prépositions de lieu.} \all{in dem Dorf} / \all{im Dorf} (\all{Wo?} position statique $\rightarrow$ Datif) mais \all{in das Dorf} / \all{ins Dorf} (\all{Wohin?} mouvement vers $\rightarrow$ Accusatif). Et les prépositions \all{zu, mit, bei, nach, von, aus, seit, gegenüber} : \textbf{toujours le datif} (\all{mit \textbf{der} Familie}, jamais \all{mit die Familie} ; \all{zu \textbf{dem}} / \all{zum Markt}).
\item \textbf{Calques et faux amis du français.}
    \begin{itemize}
    \item « Assister à une fête » $\neq$ \all{assistieren} $\rightarrow$ \all{an einem Fest \textbf{teilnehmen}} (+\all{an} + Dat) ou \all{ein Fest \textbf{besuchen}} ou \all{dabei \textbf{sein}}.
    \item « Porter des vêtements » $\neq$ \all{portieren} $\rightarrow$ \all{Kleidung \textbf{tragen}} (verbe fort : \all{trägt, trug, getragen}).
    \item « Fête » : jour férié = \all{der Feiertag} ; événement célébré = \all{das Fest} / \all{die Feier}.
    \item « Traditionnel » $\rightarrow$ \all{traditionell} (pas \all{traditionnel}).
    \item « Village » $\rightarrow$ \all{das Dorf} (pas \all{das Village}).
    \end{itemize}
\end{enumerate}
\end{attention}

\newpage

\coursec{Exercices}

\courssub{Je vérifie mes connaissances}

\begin{exercice}[QCM : le bon article]
Pour chaque phrase, choisis la bonne réponse (a, b ou c) et recopie la phrase complète.
\begin{enumerate}
\item \all{Wir feiern das Fest mit \_\_\_ Familie.} \quad a) \all{die} \quad b) \all{der} \quad c) \all{dem}
\item \all{Die Kinder tanzen auf \_\_\_ Wiese.} \quad a) \all{der} \quad b) \all{die} \quad c) \all{dem}
\item \all{Am Wochenende fahren wir zu \_\_\_ Großeltern.} \quad a) \all{die} \quad b) \all{den} \quad c) \all{der}
\item \all{In \_\_\_ Straßen gibt es viele Menschen.} \quad a) \all{die} \quad b) \all{den} \quad c) \all{dem}
\item \all{Ich habe \_\_\_ Fest besucht.} \quad a) \all{das} \quad b) \all{dem} \quad c) \all{den}
\end{enumerate}
\end{exercice}

\begin{exercice}[Vrai ou faux ? Justifie ta réponse]
Réponds par \all{richtig} ou \all{falsch} et justifie chaque réponse par une phrase complète en français ou en allemand.
\begin{enumerate}
\item \all{Bei} est toujours suivi du datif.
\item Le verbe \all{kommen} forme son Perfekt avec l'auxiliaire \all{haben}.
\item \all{Der Tanz} signifie « la danse ».
\item Après \all{in} avec la question \all{wo ?}, on emploie le datif.
\item Le participe passé de \all{essen} est \all{geesst}.
\end{enumerate}
\end{exercice}

\begin{exercice}[Vocabulaire : à toi de compléter]
Complète chaque phrase avec le mot qui convient, choisi dans la liste : \all{die Tradition, das Essen, die Kleidung, der Umzug, feiern}.
\begin{enumerate}
\item Die Frauen tragen schöne traditionelle \_\_\_\_\_\_\_\_\_.
\item Beim Fest gibt es gutes \_\_\_\_\_\_\_\_\_ : Fisch, Reis und Bananen.
\item Viele Familien \_\_\_\_\_\_\_\_\_ das Fest jedes Jahr im Dorf.
\item Der \_\_\_\_\_\_\_\_\_ geht durch die Straßen der Stadt.
\item Dieses Fest ist eine alte \_\_\_\_\_\_\_\_\_ unserer Region.
\end{enumerate}
\end{exercice}

\begin{exercice}[Conjugaison : le Perfekt]
Mets les verbes entre parenthèses au Perfekt. Attention au choix de l'auxiliaire !
\begin{enumerate}
\item Wir \_\_\_\_\_\_\_\_\_ das Fest mit der ganzen Familie \_\_\_\_\_\_\_\_\_. (feiern)
\item Die Gäste \_\_\_\_\_\_\_\_\_ nach Douala \_\_\_\_\_\_\_\_\_. (kommen)
\item Ich \_\_\_\_\_\_\_\_\_ viel Fisch und Reis \_\_\_\_\_\_\_\_\_. (essen)
\item Die Jugendlichen \_\_\_\_\_\_\_\_\_ die ganze Nacht \_\_\_\_\_\_\_\_\_. (tanzen)
\item Wir \_\_\_\_\_\_\_\_\_ unsere Großeltern im Dorf \_\_\_\_\_\_\_\_\_. (besuchen)
\end{enumerate}
\end{exercice}

\courssub{Je m'entraîne}

\begin{exercice}[Texte à trous : le datif]
Complète le texte suivant avec l'article qui convient au \textbf{datif} (\all{dem, der, den, einem, einer} ou l'article zéro si nécessaire).

\begin{textebox}[Das Fest in meinem Dorf]
Am Wochenende feiern wir mit \_\_\_ Familie ein großes Fest. In \_\_\_ Straßen tanzen viele Menschen. Bei \_\_\_ Umzug tragen die Frauen schöne Kleider. Auf \_\_\_ Wiese hinter dem Haus essen wir mit \_\_\_ Nachbarn. Die Kinder spielen mit \_\_\_ Freunden, und am Abend sitzen wir bei \_\_\_ Großeltern und erzählen Geschichten.
\end{textebox}
\end{exercice}

\begin{exercice}[Le bon auxiliaire : haben ou sein ?]
Mets les phrases suivantes au Perfekt en choisissant le bon auxiliaire et le bon participe passé.
\begin{enumerate}
\item Das Fest (stattfinden) \_\_\_\_\_\_\_\_\_ im Dorf \_\_\_\_\_\_\_\_\_.
\item Wir (fahren) \_\_\_\_\_\_\_\_\_ nach Yaoundé \_\_\_\_\_\_\_\_\_.
\item Die Gäste (bleiben) \_\_\_\_\_\_\_\_\_ bis Mitternacht \_\_\_\_\_\_\_\_\_.
\item Ich (essen) \_\_\_\_\_\_\_\_\_ eine Banane \_\_\_\_\_\_\_\_\_.
\item Die Kinder (tanzen) \_\_\_\_\_\_\_\_\_ auf der Wiese \_\_\_\_\_\_\_\_\_.
\item Meine Tante (kommen) \_\_\_\_\_\_\_\_\_ aus Bafoussam \_\_\_\_\_\_\_\_\_.
\item Wir (feiern) \_\_\_\_\_\_\_\_\_ drei Tage lang \_\_\_\_\_\_\_\_\_.
\item Er (besuchen) \_\_\_\_\_\_\_\_\_ das Fest \_\_\_\_\_\_\_\_\_.
\end{enumerate}
\end{exercice}

\begin{exercice}[Appariements : mots et définitions]
Relie chaque mot allemand (1 à 6) à sa définition française (a à f).
\begin{center}
\begin{tabularx}{\linewidth}{|l|X|}
\hline
\rowcolor{popblueL} \textbf{Deutsch} & \textbf{Français} \\
\hline
1. \all{der Umzug, die Umzüge} & a) les habits que l'on porte pour la fête \\
\hline
2. \all{die Tradition, die Traditionen} & b) un défilé dans les rues \\
\hline
3. \all{die Kleidung (Sg.)} & c) une coutume transmise de génération en génération \\
\hline
4. \all{der Gast, die Gäste} & d) une personne invitée à une fête \\
\hline
5. \all{das Feuerwerk, die Feuerwerke} & e) un repas pris ensemble \\
\hline
6. \all{das Festessen, die Festessen} & f) des fusées colorées dans le ciel la nuit \\
\hline
\end{tabularx}
\end{center}
\end{exercice}

\begin{exercice}[Datif ou accusatif ? Discrimination]
Dans le paragraphe suivant, \textbf{souligne} les compléments au datif et \textbf{entoure} les compléments à l'accusatif. Justifie ensuite chaque choix avec la question \all{wo ?} (datif) ou \all{wohin ?} (accusatif).

\begin{textebox}[Der Festtag]
Am Morgen gehen wir in die Kirche. Nach dem Gottesdienst tanzen die Kinder auf dem Platz. Am Nachmittag fahren wir zu unseren Großeltern ins Dorf. Dort sitzen wir bei der Familie und essen mit den Nachbarn. Am Abend gehen alle in den großen Saal und tanzen mit den Musikern bis Mitternacht.
\end{textebox}
\end{exercice}

\courssub{Je me prépare au Bac}

\begin{exercice}[Compréhension de texte type Bac : le carnaval de Cologne]
Lis le texte suivant, puis réponds aux questions \textbf{en allemand, par des phrases complètes}.

\begin{textebox}[Der Kölner Karneval]
Jedes Jahr im Februar feiert die Stadt Köln ihren berühmten Karneval. Das Fest beginnt am Donnerstag und dauert sechs Tage. In den Straßen der Stadt tragen die Menschen bunte Kostüme: Clowns, Tiere, Prinzen und Prinzessinnen. Viele Familien feiern mit ihren Freunden und Nachbarn.

Am Rosenmontag gibt es einen großen Umzug. Mehr als eine Million Besucher kommen aus ganz Deutschland und aus dem Ausland nach Köln. Die Menschen auf den Wagen werfen Bonbons und Blumen in die Menge, und alle rufen „Kölle Alaaf!“. Die Kinder freuen sich besonders über die Süßigkeiten.

Beim Karneval gibt es auch viel Musik. In den Kneipen und auf den Plätzen singen und tanzen die Leute bis spät in die Nacht. Am Ende des Festes verbrennen die Kölner eine große Puppe aus Stroh: das ist eine alte Tradition. Dann ist der Karneval für dieses Jahr zu Ende, aber die Menschen warten schon auf das nächste Fest.
\end{textebox}

\textbf{Fragen zum Text :}
\begin{enumerate}
\item Wann beginnt der Karneval in Köln und wie lange dauert er ?
\item Was tragen die Menschen in den Straßen ?
\item Wer kommt am Rosenmontag nach Köln ?
\item Was werfen die Menschen auf den Wagen in die Menge ?
\item Was passiert am Ende des Festes mit der Puppe aus Stroh ?
\end{enumerate}

\textbf{Fragen zur Sprache :}
\begin{enumerate}
\item Relève dans le texte trois prépositions suivies du datif et recopie le groupe nominal complet.
\item Mets au Perfekt : \all{Die Menschen tragen bunte Kostüme und tanzen in den Straßen.}
\end{enumerate}
\end{exercice}

\begin{exercice}[Thème et version type Bac]
\textbf{A. Thème.} Traduis en allemand les phrases suivantes. Attention au datif, au Perfekt et à la place du verbe !
\begin{enumerate}
\item Nous avons fêté la fête avec toute la famille.
\item Les visiteurs sont venus à Douala.
\item Les femmes portent des vêtements traditionnels.
\item Les enfants ont dansé dans les rues jusqu'à minuit.
\item Chez mes grands-parents, nous avons mangé du poisson avec du riz.
\end{enumerate}

\textbf{B. Version.} Traduis en français le paragraphe suivant.

\begin{textebox}[Das Erntedankfest]
Im Herbst feiern viele Familien in Deutschland das Erntedankfest. An diesem Tag danken die Menschen für das Essen aus den Gärten und von den Feldern. In den Kirchen liegen Obst, Gemüse und Brot auf den Tischen. Nach dem Gottesdienst essen die Familien mit ihren Gästen ein großes Festessen. Die Kinder haben oft Ferien und helfen den Großeltern im Garten.
\end{textebox}
\end{exercice}

\begin{exercice}[Expression écrite type Bac]
\textbf{Sujet :} \all{Erzählen Sie von einem traditionellen Fest in Ihrer Region.}

Rédige un texte de \textbf{60 à 80 mots} en allemand. Tu raconteras une fête traditionnelle de ta région (par exemple le Ngondo, le Nyem-Nyem, une fête de village ou une fête de famille) en répondant aux questions suivantes :
\begin{itemize}
\item Wann und wo findet das Fest statt ?
\item Wer kommt zu dem Fest ?
\item Was essen und trinken die Gäste ?
\item Was machen die Menschen ? (tanzen, singen, spielen...)
\item Wie war das Fest im letzten Jahr ? (Perfekt !)
\end{itemize}

\textbf{Grille d'évaluation :}
\begin{center}
\begin{tabularx}{\linewidth}{|X|c|}
\hline
\rowcolor{popblueL} \textbf{Critère} & \textbf{Points} \\
\hline
Respect du sujet et du nombre de mots (60--80) & 4 \\
\hline
Vocabulaire du thème « les fêtes traditionnelles » & 4 \\
\hline
Emploi correct du datif après les prépositions de lieu & 4 \\
\hline
Emploi correct du Perfekt (auxiliaire + participe) & 4 \\
\hline
Orthographe, majuscules des noms, place du verbe & 4 \\
\hline
\end{tabularx}
\end{center}
\end{exercice}

\newpage

\coursec{Corrigés des exercices}

\courssub{Corrigés des activités et exercices}

\begin{corrige}[Exercice 1 — Cas, articles et prépositions de lieu]
\begin{enumerate}
\item \all{Wir feiern das Fest mit \textbf{der} Familie.} \quad (Préposition \all{mit} + datif : \all{die Familie} $\rightarrow$ \all{der Familie})
\item \all{Die Kinder tanzen auf \textbf{der} Wiese.} \quad (Préposition \all{auf} + datif \emph{(Wo?)} : \all{die Wiese} $\rightarrow$ \all{der Wiese})
\item \all{Am Wochenende fahren wir zu \textbf{den} Großeltern.} \quad (Préposition \all{zu} + datif pluriel : \all{die Großeltern} $\rightarrow$ \all{den Großeltern})
\item \all{In \textbf{den} Straßen gibt es viele Menschen.} \quad (Préposition \all{in} + datif pluriel \emph{(Wo?)} : \all{die Straßen} $\rightarrow$ \all{den Straßen})
\item \all{Ich habe \textbf{das} Fest besucht.} \quad (Accusatif objet direct : \all{das Fest} reste \all{das})
\end{enumerate}
\end{corrige}

\begin{corrige}[Exercice 2 — Vrai ou faux ? Justifie ta réponse]
\begin{enumerate}
\item \all{richtig}. \all{Bei} est une préposition qui gouverne toujours le datif (ex. : \all{bei dem Fest, bei der Familie}).
\item \all{falsch}. Le verbe \all{kommen} (verbe de mouvement vers un lieu) forme son Perfekt avec l'auxiliaire \all{sein} : \all{ich bin gekommen}.
\item \all{richtig}. \all{der Tanz, die Tänze} signifie bien « la danse / les danses ».
\item \all{richtig}. La préposition \all{in} (comme \all{an, auf, hinter, in, neben, über, unter, vor, zwischen}) prend le datif quand elle répond à la question \all{wo?} (position statique).
\item \all{falsch}. Le participe passé de \all{essen} (verbe fort, alternance vocalique \cle{e $\rightarrow$ a $\rightarrow$ e}) est \all{gegessen} (et non *geesst).
\end{enumerate}
\end{corrige}

\begin{corrige}[Exercice 3 — Vocabulaire : complétion et genres]
\begin{enumerate}
\item Die Frauen tragen schöne traditionelle \textbf{Kleidung}.
\item Beim Fest gibt es gutes \textbf{Essen} : Fisch, Reis und Bananen.
\item Viele Familien \textbf{feiern} das Fest jedes Jahr im Dorf.
\item Der \textbf{Umzug} geht durch die Straßen der Stadt.
\item Dieses Fest ist eine alte \textbf{Tradition} unserer Region.
\end{enumerate}
\emph{Rappel des genres et pluriels :} \all{die Kleidung} (Sg. seulement, collectif), \all{das Essen} (Sg., collectif), \all{der Umzug, die Umzüge}, \all{die Tradition, die Traditionen}, \all{feiern} (verbe régulier, auxiliaire \all{haben}).
\end{corrige}

\begin{corrige}[Exercice 4 — Conjugaison : le Perfekt]
\begin{enumerate}
\item Wir \textbf{haben} das Fest mit der ganzen Familie \textbf{gefeiert}. (verbe régulier, auxiliaire \all{haben})
\item Die Gäste \textbf{sind} nach Douala \textbf{gekommen}. (verbe de mouvement \all{kommen} $\rightarrow$ auxiliaire \all{sein}, participe \all{gekommen})
\item Ich \textbf{habe} viel Fisch und Reis \textbf{gegessen}. (verbe fort \all{essen} $\rightarrow$ \all{gegessen}, auxiliaire \all{haben} car transitif)
\item Die Jugendlichen \textbf{haben} die ganze Nacht \textbf{getanzt}. (verbe régulier \all{tanzen} $\rightarrow$ \all{getanzt}, auxiliaire \all{haben} — pas de déplacement de A vers B)
\item Wir \textbf{haben} unsere Großeltern im Dorf \textbf{besucht}. (verbe régulier \all{besuchen} $\rightarrow$ \all{besucht}, auxiliaire \all{haben})
\end{enumerate}
\end{corrige}

\begin{corrige}[Exercice 5 — Texte à trous : le datif]
\begin{textebox}[Das Fest in meinem Dorf — Solution]
Am Wochenende feiern wir mit \textbf{der} Familie ein großes Fest. In \textbf{den} Straßen tanzen viele Menschen. Bei \textbf{dem} Umzug tragen die Frauen schöne Kleider. Auf \textbf{der} Wiese hinter \textbf{dem} Haus essen wir mit \textbf{den} Nachbarn. Die Kinder spielen mit \textbf{den} Freunden, und am Abend sitzen wir bei \textbf{den} Großeltern und erzählen Geschichten.
\end{textebox}

\textbf{Détail des cas (tous datifs) :}
\begin{itemize}
\item \all{mit der Familie} (\all{mit} + datif sing. fém.)
\item \all{In den Straßen} (\all{in} + datif pluriel \emph{Wo?})
\item \all{Bei dem Umzug} (\all{bei} + datif sing. masc. — souvent contracté en \all{beim})
\item \all{Auf der Wiese} (\all{auf} + datif sing. fém. \emph{Wo?})
\item \all{hinter dem Haus} (\all{hinter} + datif sing. neutre \emph{Wo?})
\item \all{mit den Nachbarn} (\all{mit} + datif pluriel)
\item \all{mit den Freunden} (\all{mit} + datif pluriel)
\item \all{bei den Großeltern} (\all{bei} + datif pluriel)
\end{itemize}
\end{corrige}

\begin{corrige}[Exercice 6 — Le bon auxiliaire : haben ou sein ?]
\begin{enumerate}
\item Das Fest \textbf{hat} im Dorf \textbf{stattgefunden}. (\all{stattfinden} $\rightarrow$ auxiliaire \all{haben} pour un événement organisé)
\item Wir \textbf{sind} nach Yaoundé \textbf{gefahren}. (\all{fahren} mouvement $\rightarrow$ \all{sein})
\item Die Gäste \textbf{sind} bis Mitternacht \textbf{geblieben}. (\all{bleiben} changement d'état / position $\rightarrow$ \all{sein})
\item Ich \textbf{habe} eine Banane \textbf{gegessen}. (\all{essen} transitif $\rightarrow$ \all{haben})
\item Die Kinder \textbf{haben} auf der Wiese \textbf{getanzt}. (\all{tanzen} activité sans déplacement de A vers B $\rightarrow$ \all{haben})
\item Meine Tante \textbf{ist} aus Bafoussam \textbf{gekommen}. (\all{kommen} mouvement $\rightarrow$ \all{sein})
\item Wir \textbf{haben} drei Tage lang \textbf{gefeiert}. (\all{feiern} $\rightarrow$ \all{haben})
\item Er \textbf{hat} das Fest \textbf{besucht}. (\all{besuchen} transitif $\rightarrow$ \all{haben})
\end{enumerate}
\end{corrige}

\begin{corrige}[Exercice 7 — Appariements : mots et définitions]
\begin{center}
\begin{tabularx}{\linewidth}{|l|X|}
\hline
\rowcolor{popblueL} \textbf{Deutsch} & \textbf{Français} \\
\hline
1. \all{der Umzug, die Umzüge} & b) un défilé dans les rues \\
\hline
2. \all{die Tradition, die Traditionen} & c) une coutume transmise de génération en génération \\
\hline
3. \all{die Kleidung (Sg.)} & a) les habits que l'on porte pour la fête \\
\hline
4. \all{der Gast, die Gäste} & d) une personne invitée à une fête \\
\hline
5. \all{das Feuerwerk, die Feuerwerke} & f) des fusées colorées dans le ciel la nuit \\
\hline
6. \all{das Festessen, die Festessen} & e) un repas pris ensemble \\
\hline
\end{tabularx}
\end{center}
\end{corrige}

\begin{corrige}[Exercice 8 — Datif ou accusatif ? Discrimination]
\textbf{Analyse du paragraphe :}

\begin{textebox}[Der Festtag — Analyse]
Am Morgen gehen wir \textbf{in die Kirche} (Akk: \all{wohin?}). Nach \underline{dem Gottesdienst} (Dat: \all{nach} + Dat) tanzen die Kinder \underline{auf dem Platz} (Dat: \all{wo?}). Am Nachmittag fahren wir \underline{zu unseren Großeltern} (Dat: \all{zu} + Dat) \textbf{ins Dorf} (Akk: \all{in das} $\rightarrow$ \all{wohin?}). Dort sitzen wir \underline{bei der Familie} (Dat: \all{bei} + Dat) und essen \underline{mit den Nachbarn} (Dat: \all{mit} + Dat). Am Abend gehen alle \textbf{in den großen Saal} (Akk: \all{wohin?}) und tanzen \underline{mit den Musikern} (Dat: \all{mit} + Dat) bis Mitternacht.
\end{textebox}

\textbf{Légende :} \underline{souligné} = Datif (question \all{wo?} ou préposition imposant le datif) | \textbf{gras} = Accusatif (question \all{wohin?}).

\textbf{Justifications :}
\begin{itemize}
\item \all{in die Kirche} : mouvement vers un lieu (\all{wohin?}) $\rightarrow$ Accusatif.
\item \all{dem Gottesdienst} : \all{nach} impose toujours le Datif.
\item \all{auf dem Platz} : position statique (\all{wo?}) $\rightarrow$ Datif.
\item \all{zu unseren Großeltern} : \all{zu} impose toujours le Datif.
\item \all{ins Dorf} (= \all{in das Dorf}) : mouvement vers un lieu (\all{wohin?}) $\rightarrow$ Accusatif.
\item \all{bei der Familie} : \all{bei} impose toujours le Datif.
\item \all{mit den Nachbarn} : \all{mit} impose toujours le Datif.
\item \all{in den großen Saal} : mouvement vers un lieu (\all{wohin?}) $\rightarrow$ Accusatif.
\item \all{mit den Musikern} : \all{mit} impose toujours le Datif.
\end{itemize}
\end{corrige}

\begin{corrige}[Exercice 9 — Compréhension de texte type Bac : le carnaval de Cologne]
\textbf{Fragen zum Text (réponses modèles) :}
\begin{enumerate}
\item Der Karneval in Köln beginnt am Donnerstag und dauert sechs Tage.
\item Die Menschen in den Straßen tragen bunte Kostüme (Clowns, Tiere, Prinzen und Prinzessinnen).
\item Am Rosenmontag kommen mehr als eine Million Besucher aus ganz Deutschland und aus dem Ausland nach Köln.
\item Die Menschen auf den Wagen werfen Bonbons und Blumen in die Menge.
\item Am Ende des Festes verbrennen die Kölner eine große Puppe aus Stroh.
\end{enumerate}

\textbf{Fragen zur Sprache :}
\begin{enumerate}
\item Trois prépositions suivies du datif (groupe nominal complet) :
\begin{itemize}
\item \all{in den Straßen} (l. 3 : \all{in} + datif pluriel \emph{Wo?})
\item \all{auf den Wagen} (l. 6 : \all{auf} + datif pluriel \emph{Wo?})
\item \all{in den Kneipen} (l. 8 : \all{in} + datif pluriel \emph{Wo?}) — ou \all{auf den Plätzen} / \all{aus Stroh} (\all{aus} + datif).
\end{itemize}
\item Perfekt : \all{Die Menschen \textbf{haben} bunte Kostüme \textbf{getragen} und \textbf{haben} in den Straßen \textbf{getanzt}.}
(Verbes \all{tragen} [fort, \all{a $\rightarrow$ a}] et \all{tanzen} [régulier] $\rightarrow$ auxiliaire \all{haben} ; participes \all{getragen} et \all{getanzt}).
\end{enumerate}

\textbf{Barème indicatif (sur 10) :} Compréhension 5 pts (1 pt par réponse correcte et complète) + Langue 5 pts (1,5 pt par groupe nominal correct + 2 pts pour la phrase au Perfekt correcte).
\end{corrige}

\begin{corrige}[Exercice 10 — Thème et version type Bac]
\textbf{A. Thème (propositions de traduction) :}
\begin{enumerate}
\item \all{Wir haben das Fest mit der ganzen Familie gefeiert.}
\item \all{Die Besucher sind nach Douala gekommen.} (ou \all{Die Gäste sind...})
\item \all{Die Frauen tragen traditionelle Kleidung.} (ou \all{traditionelle Kleider} — présent de narration)
\item \all{Die Kinder haben in den Straßen bis Mitternacht getanzt.}
\item \all{Bei meinen Großeltern haben wir Fisch mit Reis gegessen.}
\end{enumerate}

\textbf{Points de vigilance :}
\begin{itemize}
\item Place du verbe en position 2 (phrase principale) : auxiliaire \all{haben/sind} en 2\textsuperscript{e} place, participe à la fin.
\item Choix de l'auxiliaire : \all{sein} pour \all{kommen} (mouvement), \all{haben} pour \all{feiern, tragen, tanzen, essen, besuchen}.
\item Datif après \all{mit} (\all{der Familie}), \all{bei} (\all{meinen Großeltern}), \all{in} + datif (\all{den Straßen} \emph{Wo?}).
\item Majuscules aux noms : \all{Fest, Familie, Besucher, Douala, Frauen, Kleidung, Kinder, Straßen, Mitternacht, Großeltern, Fisch, Reis}.
\end{itemize}

\textbf{B. Version (traduction proposée du texte allemand suivant) :}
\begin{textebox}[Quelltext : Erntedankfest]
Im Herbst feiern viele Familien in Deutschland das Erntedankfest. An diesem Tag danken die Menschen für das Essen, das aus den Gärten und Feldern kommt. In den Kirchen stellen sie Obst, Gemüse und Brot auf die Tische. Nach dem Gottesdienst essen die Familien ein großes Festessen mit ihren Gästen. Die Kinder haben oft Ferien und helfen den Großeltern im Garten.
\end{textebox}

\begin{quote}
« À l'automne, de nombreuses familles en Allemagne fêtent la fête des moissons (l'Action de grâce). Ce jour-là, les gens remercient pour la nourriture qui vient des jardins et des champs. Dans les églises, des fruits, des légumes et du pain sont posés sur les tables. Après l'office, les familles mangent un grand repas de fête avec leurs invités. Les enfants ont souvent des vacances et aident les grands-parents au jardin. »
\end{quote}
\end{corrige}

\begin{corrige}[Exercice 11 — Expression écrite type Bac]
\textbf{Proposition de rédaction modèle (environ 70--80 mots) :}

\begin{textebox}[Modelltext : Das Ngondo-Fest]
In meiner Region in Kamerun feiern wir das Ngondo-Fest. Es findet jedes Jahr im Dezember in Douala am Wouri-Fluss statt. Viele Menschen aus dem ganzen Land und Touristen kommen zu dem Fest. Die Gäste essen Fisch, Maniok und Bananen und trinken Palmwein. Die Menschen tanzen, singen und fahren mit Booten auf dem Fluss. Letztes Jahr habe ich das Fest mit meiner Familie besucht. Wir haben viel gegessen und bis spät in die Nacht getanzt. Es war sehr schön.
\end{textebox}

\textbf{Grille d'évaluation appliquée au modèle :}
\begin{center}
\begin{tabularx}{\linewidth}{|X|c|}
\hline
\rowcolor{popblueL} \textbf{Critère} & \textbf{Points} \\
\hline
Respect du sujet et du nombre de mots (60--80) & 4/4 (env. 78 mots) \\
\hline
Vocabulaire du thème (Fest, Tradition, essen, trinken, tanzen, singen, Gäste, Umzug...) & 4/4 \\
\hline
Datif correct (mit der Familie, in den Straßen, bei meinen Großeltern, am Wouri-Fluss, mit Booten, auf dem Fluss) & 4/4 \\
\hline
Perfekt correct (habe ... besucht, haben ... gegessen, haben ... getanzt) & 4/4 \\
\hline
Orthographe, majuscules, place du verbe (verbe 2\textsuperscript{e} position, participe final) & 4/4 \\
\hline
\textbf{Total} & \textbf{20/20} \\
\hline
\end{tabularx}
\end{center}

\textbf{Points de vigilance pour l'élève :}
\begin{itemize}
\item Vérifier le nombre de mots (compter : \all{In(1) meiner2 Region3 in4 Kamerun5 feiern6 wir7 das8 Ngondo-Fest9. Es10 findet11 jedes12 Jahr13 im14 Dezember15 in16 Douala17 am18 Wouri-Fluss19 statt20. Viele21 Menschen22 aus23 dem24 ganzen25 Land26 und27 Touristen28 kommen29 zu30 dem31 Fest32. Die33 Gäste34 essen35 Fisch,36 Maniok37 und38 Bananen39 und40 trinken41 Palmwein42. Die43 Menschen44 tanzen,45 singen46 und47 fahren48 mit49 Booten50 auf51 dem52 Fluss53. Letztes54 Jahr55 habe56 ich57 das58 Fest59 mit60 meiner61 Familie62 besucht63. Wir64 haben65 viel66 gegessen67 und68 bis69 spät70 in71 die72 Nacht73 getanzt74. Es75 war76 sehr77 schön78.} = 78 mots).
\item Ne pas oublier la majuscule à \textbf{tous} les noms (\all{Fest, Jahr, Dezember, Douala, Fluss, Menschen, Land, Touristen, Gäste, Fisch, Maniok, Bananen, Palmwein, Booten, Jahr, Familie, Nacht}).
\item Structure \all{Perfekt} : auxiliaire en position 2, participe à la fin (\all{habe ... besucht, haben ... gegessen, haben ... getanzt}).
\item Prépositions de lieu + Datif : \all{in meiner Region, in Kamerun, in Douala, am (= an dem) Wouri-Fluss, zu dem Fest, mit Booten, auf dem Fluss, bei meiner Familie, in der Nacht}.
\end{itemize}
\end{corrige}

\newpage

\coursec{Fiche bilan}

\courssub{Fiche Bilan — Les fêtes traditionnelles}

\begin{popfigure}
\begin{tikzpicture}[
    mindmap,
    grow cyclic,
    every node/.style={concept, execute at begin node=\hskip0pt},
    concept color=popblue,
    level 1/.style={level distance=4.5cm, sibling angle=360/6},
    level 2/.style={level distance=3cm, sibling angle=360/4},
    texte/.style={concept, font=\small, text width=2.2cm, align=center},
    fete/.style={concept color=poporange, font=\small, text width=2.2cm, align=center},
    vocab/.style={concept color=popgreen, font=\small, text width=2.2cm, align=center},
    gram/.style={concept color=poppurple, font=\small, text width=2.2cm, align=center},
    expr/.style={concept color=poppink, font=\small, text width=2.2cm, align=center},
    cult/.style={concept color=popgold, font=\small, text width=2.2cm, align=center}
]
\node[texte, font=\bfseries\large, text width=3cm] {Fêtes\\ traditionnelles}
    child[clockwise from=60] { node[fete, align=center]{Types de fêtes\\ (Weihnachten, Ostern,\\ Kirchweih, Ngondo)} }
    child[clockwise from=0] { node[vocab, align=center]{Vocabulaire clé\\ (das Fest, die Tradition,\\ feiern, der Brauch)} }
    child[clockwise from=-60] { node[gram, align=center]{Grammaire 1\\ Datif \& Prépositions\\ de lieu (in, an, auf)} }
    child[clockwise from=-120] { node[gram, align=center]{Grammaire 2\\ Perfekt\\ (haben/sein + Partizip II)} }
    child[clockwise from=-180] { node[expr, align=center]{Expression\\ Résumé, description,\\ avis personnel} }
    child[clockwise from=120] { node[cult, align=center]{Culture\\ Allemagne / Cameroun\\ Coutumes comparées} };
\end{tikzpicture}
\legende{Carte mentale récapitulative de la leçon ALA21}
\end{popfigure}

\begin{definition}[Lexique clé — Les fêtes traditionnelles]
\begin{center}
\begin{tabularx}{\linewidth}{|X|X|X|}
\hline
\rowcolor{popblueL} \textbf{Deutsch (Artikel + Plural)} & \textbf{Français} & \textbf{Exemple / Famille de mots} \\
\hline
das Fest, die Feste & la fête & das Dorffest, das Familienfest, festlich \\
die Tradition, die Traditionen & la tradition & eine alte Tradition pflegen, traditionell \\
der Brauch, die Bräuche & la coutume, l'usage & nach altem Brauch, bräuchlich \\
feiern (regelmäßig) & fêter, célébrer & wir feiern Weihnachten, die Feier \\
der Tanz, die Tänze & la danse & der traditionelle Tanz, tanzen \\
das Essen, - & la nourriture, le repas & das Festessen, essen (isst, aß, gegessen) \\
die Kleidung, - & le vêtement, l'habillement & die Tracht (costume trad.), sich kleiden \\
der Schmuck, - & le bijou, la parure & der Kopfschmuck, schmücken \\
die Musik, - & la musique & Volksmusik spielen, musikalisch \\
der Gast, die Gäste & l'invité & die Gäste bewirten, gastfreundlich \\
die Feier, die Feiern & la célébration, la fête & die Geburtstagsfeier, feierlich \\
einladen (lädt ein, lud ein, hat eingeladen) & inviter & wir laden Gäste ein, die Einladung \\
besuchen (regelmäßig) & rendre visite & wir besuchen die Großeltern, der Besuch \\
\hline
\end{tabularx}
\end{center}
\end{definition}

\begin{propriete}[Règles de grammaire essentielles]
\begin{center}
\begin{tabularx}{\linewidth}{|X|X|}
\hline
\rowcolor{poppurpleL} \textbf{Thème} & \textbf{Règle essentielle + Exemples} \\
\hline
\textbf{Datif (cas du destinataire / lieu où)} & \textbf{Article défini :} dem (m/n), der (f), den (pl). \\
& \textbf{Article indéfini :} einem (m/n), einer (f), einem (n), - (pl). \\
& \textbf{Ex. (Destinataire) :} \all{Wir helfen \textbf{dem} Mann.} (Nous aidons l'homme). \\
& \textbf{Ex. (Lieu où) :} \all{Wir sind \textbf{im} (in dem) Saal.} (Nous sommes dans la salle). \\
\hline
\textbf{Prépositions de lieu (Wechselpräpositionen)} & \textbf{in, an, auf, unter, über, vor, hinter, neben, zwischen} + \textbf{Datif} = \textbf{WO?} (Lieu statique / Position). \\
& \all{Wir feiern \textbf{im} Garten.} (dans le jardin — \emph{in dem}). \\
& \all{Das Bild hängt \textbf{am} (an dem) Baum.} (à l'arbre). \\
& \all{Wir sitzen \textbf{auf dem} Boden.} (par terre). \\
& \textbf{Contraste (Accusatif = WOHIN? Mouvement) :} \all{Wir gehen \textbf{in den} Garten.} \\
\hline
\textbf{Perfekt (Passé composé)} & \textbf{Structure :} Auxiliaire conjugué (Pos. 2) + Partizip II (à la fin). \\
& \textbf{Auxiliaire \textbf{sein} :} Verbes de mouvement (gehen, fahren, kommen, laufen, fliegen) \& changement d'état (aufwachen, einschlafen, sterben). \\
& \all{Ich \textbf{bin} nach Hause \textbf{gefahren}.} \\
& \textbf{Auxiliaire \textbf{haben} :} Verbes transitifs, réflexifs, la plupart des autres. \\
& \all{Wir \textbf{haben} viel \textbf{gegessen}.} \\
& \textbf{Partizip II régulier :} ge- + radical + -t (\all{spielen $\to$ gespielt}, \all{feiern $\to$ gefeiert}). \\
& \textbf{Partizip II irrégulier (verbes forts) :} ge- + radical modifié + -en (\all{singen $\to$ gesungen}, \all{essen $\to$ gegessen}, \all{trinken $\to$ getrunken}, \all{tanzen $\to$ getanzt}). \\
& \textbf{Verbes en -ieren / préfixes inséparables (be-, ent-, er-, ver-, zer-, ge-, miss-) :} \textbf{SANS ge-}. \\
& \all{studieren $\to$ studiert}, \all{besuchen $\to$ besucht}, \all{erzählen $\to$ erzählt}. \\
& \textbf{Verbes à particule séparable :} \textbf{ge-} s'insère entre la particule et le radical. \\
& \all{einladen $\to$ eingeladen}, \all{aufstehen $\to$ aufgestanden}, \all{mitbringen $\to$ mitgebracht}. \\
\hline
\end{tabularx}
\end{center}
\end{propriete}

\begin{methode}[Méthodes express pour l'examen]
\begin{enumerate}
    \item \textbf{Compréhension de texte (écrit/oral) :}
    \begin{enumerate}
        \item Lire le titre, regarder images/photos $\to$ formuler des hypothèses.
        \item Lecture globale : identifier le sujet général, le type de texte, l'origine.
        \item Lecture détaillée : surligner mots-clés, connecteurs (und, aber, weil, dann), repérer les temps verbaux (Präsens, Perfekt, Präteritum).
        \item Répondre aux questions \textbf{en allemand}, phrases complètes, en reprenant le vocabulaire du texte.
    \end{enumerate}
    \item \textbf{Traduction Thème (Français $\to$ Allemand) :}
    \begin{enumerate}
        \item Identifier le verbe principal et son temps (Présent ou Perfekt).
        \item Conjuguer le verbe $\to$ \textbf{Position 2}.
        \item Déterminer les cas : Sujet = Nominatif, Objet direct = Accusatif, Objet indirect / Lieu où = Datif.
        \item Appliquer l'ordre \textbf{TEKAMOLO} (Temps — Causalité — Modal — Lieu — Objet).
        \item Si Perfekt : placer le \textbf{Partizip II à la fin}. Vérifier l'auxiliaire (haben/sein) et la formation du Partizip II (ge- / pas de ge- / ge- au milieu).
    \end{enumerate}
    \item \textbf{Production écrite « Ma fête traditionnelle » :}
    \begin{enumerate}
        \item \textbf{Présentation :} Nom de la fête, date, lieu, origine (Ngondo, Nguon, Weihnachten, Kirchweih…).
        \item \textbf{Préparatifs (Perfekt) :} \all{Wir haben das Haus geschmückt. Die Frauen haben traditionelles Essen gekocht.}
        \item \textbf{Déroulement (Präsens ou Präteritum) :} \all{Die Leute tanzen und singen. Es gibt Trommelmusik.}
        \item \textbf{Détails culturels :} Vêtements (\all{die Tracht}, \all{der Schmuck}), Nourriture (\all{das Essen}), Danses (\all{der Tanz}).
        \item \textbf{Avis personnel :} \all{Ich finde das Fest toll, weil die Stimmung fröhlich ist.}
    \end{enumerate}
\end{enumerate}
\end{methode}

\begin{aretenir}[Checklist avant le Bac]
$\square$ Je sais décliner l'article défini et indéfini au \textbf{Datif} (dem, der, dem, den / einem, einer, einem, -).\\
$\square$ Je distingue \textbf{Accusatif} (Wohin ? direction) et \textbf{Datif} (Wo ? position) avec les \cle{Wechselpräpositionen in, an, auf, unter, über, vor, hinter, neben, zwischen}.\\
$\square$ Je forme le \textbf{Partizip II} des verbes réguliers (\textbf{ge-...-t}) et je connais les formes irrégulières principales (\textbf{ge-...-en} : gegessen, getrunken, getanzt, gesungen, gefahren, gegangen, gekommen).\\
$\square$ Je forme correctement le \textbf{Partizip II des verbes à particule séparable} (\textbf{ge-} au milieu : \all{eingeladen, aufgestanden, mitgebracht}).\\
$\square$ Je sais que les verbes en \textbf{-ieren} et à \textbf{préfixe inséparable} ne prennent \textbf{PAS de ge-} (\all{studiert, besucht, erzählt}).\\
$\square$ Je choisis l'auxiliaire correct au \textbf{Perfekt} : \textbf{sein} pour mouvement/déplacement (gehen, fahren, kommen, laufen) et changement d'état (aufwachen, einschlafen), \textbf{haben} sinon.\\
$\square$ Je place le verbe conjugué en \textbf{position 2} et le \textbf{Partizip II à la fin} de la phrase principale.\\
$\square$ J'utilise le vocabulaire thématique (\cle{das Fest, die Tradition, der Brauch, feiern, der Tanz, das Essen, die Tracht, der Gast, die Einladung}) avec le bon genre et le bon pluriel.\\
$\square$ Je sais écrire un court paragraphe structuré (Présentation $\to$ Description $\to$ Avis) sur une fête camerounaise ou germanophone.\\
$\square$ Je repère les \textbf{faux amis} : \all{das Gift} = le poison (pas le cadeau \all{das Geschenk}), \all{der Chef} = le patron (pas le chef cuisinier \all{der Koch}), \all{das Event} = l'événement (anglicisme, préférer \all{die Veranstaltung}).\\
$\square$ Je respecte la \textbf{majuscule} sur tous les noms allemands (\all{der Tanz, die Freude, das Essen, die Musik}).\\
$\square$ Je sais transformer une phrase au Présent en Perfekt et inversement.\\
$\square$ Je peux expliquer une tradition camerounaise (ex. \all{Ngondo}, \all{Nguon}, fêtes des moissons) en allemand simple (niveau A2).
\end{aretenir}

\findecours

\end{document}
